\documentclass[11pt, a4paper, logo, copyright, nonumbering]{deepseek}
\usepackage[authoryear, sort&compress, round]{natbib}
\usepackage{dblfloatfix}
\usepackage{ulem}
\usepackage{caption}
\usepackage{dramatist}
\usepackage{xspace}
\usepackage{pifont} % http://ctan.org/pkg/pifont
\usepackage{multirow}
\usepackage{tcolorbox}
\usepackage{xltabular}
\usepackage{longtable}
\usepackage{hyperref}
\interfootnotelinepenalty=10000

\usepackage{amsfonts}
\usepackage{amsmath}
\usepackage{amssymb}
\usepackage{lineno}
\usepackage{multirow}
\usepackage{adjustbox}

\usepackage[bottom]{footmisc}
% \linenumbers
\usepackage{listings}
\usepackage{CJKutf8}
\usepackage{subfigure}
\usepackage{setspace}

% damai custom package
\usepackage{dsfont}
\usepackage{array} % 用于tabularx环境
\usepackage{tabularx} % 用于tabularx环境
\usepackage{subfigure} % 用于创建子图
\usepackage{xcolor} % 用于画有颜色的框
\usepackage{tabularx}
\usepackage{booktabs}

% ########
% 附录包
\usepackage{lipsum}  % 用于生成示例文本
\usepackage{multicol} % 引入 multicol 宏包

% ########


% 找 Claude-3.5 要的
% 创建自定义颜色主题
\definecolor{codebackground}{rgb}{0.95,0.95,0.95} % 浅灰色背景
\definecolor{codeframe}{rgb}{0.8,0.8,0.8}        % 边框颜色
\definecolor{keyword}{rgb}{0.0,0.0,0.55}         % 关键词颜色 (深蓝)
\definecolor{stringcolor}{rgb}{0.65,0.13,0.13}   % 字符串颜色 (深红)
\definecolor{commentcolor}{rgb}{0.0,0.5,0.0}     % 注释颜色 (深绿)


% 定义 listings 的样式
\lstdefinestyle{code}{
    backgroundcolor=\color{codebackground},    % 背景颜色
    frame=single,                              % 代码框样式：单线框
    rulecolor=\color{codeframe},               % 框线颜色
    basicstyle=\ttfamily\footnotesize,         % 基础字体：等宽、小字号
    % keywordstyle=\color{keyword}\bfseries,     % 关键词样式：深蓝色+加粗
    keywordstyle=\bfseries,     % 关键词样式：加粗
    stringstyle=\color{stringcolor},           % 字符串样式：深红色
    commentstyle=\itshape\color{commentcolor}, % 注释样式：斜体+深绿色
    % numberstyle=\tiny\color{gray},             % 行号样式：小字号灰色
    % numbers=left,                              % 行号位置：左侧
    % numbersep=7pt,                             % 行号与代码间的间距
    showspaces=false,                          % 不显示空格
    showstringspaces=false,                    % 字符串中的空格不显示特殊标记
    showtabs=false,                            % 不显示制表符特殊标记
    tabsize=4,                                 % Tab 等价于 4 个空格
    captionpos=b,                              % 标题位置：代码下方 (bottom)
    breaklines=true,                           % 自动换行
    breakatwhitespace=true,                    % 仅在空格处换行
    sensitive=true,                            % 区分大小写
    breakindent=0pt, % 换行后无缩进
}


\lstdefinestyle{data}{
  backgroundcolor=\color{codebackground},  % 设置背景色
  frame=topbottom,  % 只在顶部和底部加框
  basicstyle=\footnotesize\ttfamily,
  columns=fullflexible,
  breaklines=true,
  extendedchars=true,
  escapechar=@,
  literate={á}{{\'a}}1 {ã}{{\~a}}1 {é}{{\'e}}1 {£}{{\pounds}}1 {–}{{-}}1 {’}{{'}}1,
% % backgroundcolor=\color{backcolour}
  % postbreak=\usebox\mypostbreak,
}

% 设置默认的 lstlisting 样式
% \lstset{style=mystyle}

% \usepackage[finalizecache,cachedir=.]{minted}


\makeatletter
\def\@BTrule[#1]{%
  \ifx\longtable\undefined
    \let\@BTswitch\@BTnormal
  \else\ifx\hline\LT@hline
    \nobreak
    \let\@BTswitch\@BLTrule
  \else
     \let\@BTswitch\@BTnormal
  \fi\fi
  \global\@thisrulewidth=#1\relax
  \ifnum\@thisruleclass=\tw@\vskip\@aboverulesep\else
  \ifnum\@lastruleclass=\z@\vskip\@aboverulesep\else
  \ifnum\@lastruleclass=\@ne\vskip\doublerulesep\fi\fi\fi
  \@BTswitch}
\makeatother

\addto\extrasenglish{
    \def\chapterautorefname{Chapter}%
    \def\sectionautorefname{Section}%
    \def\paragraphautorefname{Section}
    \def\subparagraphautorefname{Section}
    \def\paragraphautorefname{Paragraph}%
    \def\tableautorefname{Table}%
    \def\equationautorefname{Equation}%
}

\newenvironment{indentchunk}[1]%
 {\begin{list}{}%
         {\setlength{\leftmargin}{#1}}%
         \item[]%
 }
 {\end{list}}
 
\bibliographystyle{abbrvnat}
% \bibliography{main}

\reportnumber{001} % Leave blank if n/a

% \renewcommand{\today}{2021-12-08}
\renewcommand{\today}{}

% \title{\centering DeepSeek-R1: Pushing the Boundaries of Reasoning via an Innovative Post-Training Paradigm}
\title{\centering DeepSeek-R1: Incentivizing Reasoning Capability in LLMs via Reinforcement Learning}
% o1 - human-leval reasoning

% \title{DeeepSeek-R1: Incentivizing Reasoning of LLMs with Reinforcement Learning}
% \title{DeeepSeek-R1: Achieving Human-level Reasoning of LLMs without Supervision}
% \title{DeeepSeek-R1: Scaling Inference Compute of LLMs with Reinforcement Learning}

\author[*]{
% \vspace{-0.5cm}
DeepSeek-AI
\\
\small
\texttt{research@deepseek.com}
% \vspace{-0.5cm}
}
\input{commands.tex}

\newcommand{\dsviii}{DeepSeek-V3}
\renewcommand{\figurename}{Supplementary Fig.}
\renewcommand{\tablename}{Supplementary Table}

\newcommand{\fix}{\marginpar{FIX}}
\newcommand{\new}{\marginpar{NEW}}
\newcommand{\dsri}{DeepSeek-R1}
\newcommand{\dsro}{DeepSeek-R1-Zero}
\newcommand{\todo}[1]{\textcolor[rgb]{0.8,0,0.2}{#1}}

\begin{abstract}

General reasoning represents a long-standing and formidable challenge in artificial intelligence. Recent breakthroughs, exemplified by large language models (LLMs) \citep{gpt3, gpt4} and chain-of-thought prompting \citep{wei2022chain}, have achieved considerable success on foundational reasoning tasks. However, this success is heavily contingent upon extensive human-annotated demonstrations, and models' capabilities are still insufficient for more complex problems.
Here we show that the reasoning abilities of LLMs can be incentivized through pure reinforcement learning (RL), obviating the need for human-labeled reasoning trajectories. The proposed RL framework facilitates the emergent development of advanced reasoning patterns, such as self-reflection, verification, and dynamic strategy adaptation. Consequently, the trained model achieves superior performance on verifiable tasks such as mathematics, coding competitions, and STEM fields, surpassing its counterparts trained via conventional supervised learning on human demonstrations. Moreover, the emergent reasoning patterns exhibited by these large-scale models can be systematically harnessed to guide and enhance the reasoning capabilities of smaller models.

\end{abstract}

\begin{document}
\begin{CJK*}{UTF8}{gbsn}
\maketitle
\keywords{Large Language Model, Reasoning, Reinforcement Learning}

%%\pacs[JEL Classification]{D8, H51}

%%\pacs[MSC Classification]{35A01, 65L10, 65L12, 65L20, 65L70}



\section{Introduction}\label{sec1}
Reasoning capability, the cornerstone of human intelligence, enables complex cognitive tasks ranging from mathematical problem-solving to logical deduction and programming.  Recent advances in artificial intelligence have demonstrated that large language models (LLMs) can exhibit emergent behaviors, including reasoning abilities, when scaled to a sufficient size \citep{wei2022emergent, kaplan2020scaling}. However, achieving such capabilities in pre-training typically demands substantial computational resources.
 In parallel, a complementary line of research has demonstrated that large language models can be effectively augmented through chain-of-thought (CoT) prompting. This technique, which involves either providing carefully designed few-shot examples or using minimalistic prompts such as  “Let’s think step by step”\citep{wei2022chain,kojima2022large}, enables models to produce intermediate reasoning steps, thereby substantially enhancing their performance on complex tasks.
 % cite: scratch-pads, deepseek-math, 
 Similarly, further performance gains have been observed when models learn high-quality, multi-step reasoning trajectories during the post-training phase \citep{gpt4,chung2024scaling}.
 Despite their effectiveness, these approaches exhibit notable limitations. Their dependence on human-annotated reasoning traces hinders scalability and introduces cognitive biases. Furthermore, by constraining models to replicate human thought processes, their performance is inherently capped by the human-provided exemplars, which prevents the exploration of superior, non-human-like reasoning pathways.


To tackle these issues, we aim to explore the potential of LLMs for developing reasoning abilities through self-evolution in an RL framework, with minimal reliance on human labeling efforts.
Specifically, we build upon DeepSeek-V3-Base \citep{dsviii} and employ Group Relative Policy Optimization (GRPO) \citep{deepseekmath} as our RL framework. The reward signal is solely based on the correctness of final predictions against ground-truth answers, without imposing constraints on the reasoning process itself. Notably, we bypass the conventional supervised fine-tuning (SFT) phase before RL training. This design choice stems from our hypothesis that human-defined reasoning patterns may limit model exploration, whereas unrestricted RL training can better incentivize the emergence of novel reasoning capabilities in LLMs. 
Through this process, detailed in Section \ref{sec:ds-zero}, our model (referred to as DeepSeek-R1-Zero) naturally developed diverse and sophisticated reasoning behaviors.
In solving reasoning problems, the model exhibits a tendency to generate longer responses, incorporating verification, reflection, and the exploration of alternative approaches within each response. Although we do not explicitly teach the model how to reason, it successfully learns improved reasoning strategies through reinforcement learning.

 
Although DeepSeek-R1-Zero demonstrates excellent reasoning capabilities, it faces challenges such as poor readability and language mixing, occasionally combining English and Chinese within a single chain-of-thought response. Furthermore, the rule-based RL training stage of DeepSeek-R1-Zero is narrowly focused on reasoning tasks, resulting in limited performance in broader areas such as writing and open-domain question answering.
To address these challenges, we introduce DeepSeek-R1, a model trained through a multi-stage learning framework that integrates rejection sampling, reinforcement learning, and supervised fine-tuning, detailed in Section \ref{sec:ds-r1}. This training pipeline enables DeepSeek-R1 to inherit the reasoning capabilities of its predecessor, DeepSeek-R1-Zero, while aligning model behavior with human preferences through additional non-reasoning data. 

To enable broader access to
 powerful AI at a lower energy cost, we have distilled several smaller models and made them publicly available. These distilled models exhibit strong reasoning capabilities, surpassing the performance of their original instruction-tuned counterparts. We believe that these instruction-tuned versions will also significantly contribute to the research community by providing a valuable resource for understanding the mechanisms underlying long chain-of-thought (CoT) reasoning models and for fostering the development of more powerful reasoning models. We release DeepSeek-R1 series  models to the public at \url{https://huggingface.co/deepseek-ai}. 

\section{ DeepSeek-R1-Zero}

\label{sec:ds-zero}
We begin by elaborating on the training of DeepSeek-R1-Zero, which relies exclusively on reinforcement learning without supervised fine-tuning. To facilitate large-scale RL efficiency, we adopt Group Relative Policy Optimization (GRPO) \citep{deepseekmath}.


\subsection{Group Relative Policy Optimization} \label{method:grpo}
GRPO \citep{deepseekmath} is the reinforcement learning algorithm that we adopt to train DeepSeek-R1-Zero and DeepSeek-R1. It was originally proposed to simplify the training process and reduce the resource consumption of Proximal Policy Optimization (PPO) \citep{schulman2017proximal}, which is widely used in the RL stage of LLMs \citep{ouyang2022training}. 

For each question $q$, GRPO samples a group of outputs $\{o_1, o_2, \cdots, o_G\}$ from the old policy $\pi_{\theta_{old}}$ and then optimizes the policy model $\pi_{\theta}$ by maximizing the following objective:
\begin{equation}
\begin{split}
    \footnotesize
    & \mathcal{J}_{GRPO}(\theta) = \mathbb{E}{[q \sim P(Q), \{o_i\}_{i=1}^G \sim \pi_{\theta_{old}}(O|q)]}  \\
    & \frac{1}{G}\sum_{i=1}^G \left( \min \left( \frac{\pi_\theta(o_i |q)}{\pi_{\theta_{old}}(o_i |q)} A_i, \text{clip} \left( \frac{\pi_\theta(o_i |q)}{\pi_{\theta_{old}}(o_i |q)}, 1 - \epsilon, 1 + \epsilon \right)  A_i \right) - \beta \mathbb{D}_{KL}\left(\pi_{\theta} || \pi_{ref}\right)\right) ,
\end{split}
\label{eq:GRPO-obj}
\end{equation}
\begin{equation}
    \mathbb{D}_{KL}\left(\pi_{\theta} || \pi_{ref}\right) = \frac{\pi_{ref}(o_i|q)}{\pi_{\theta}(o_i|q)}- \log\frac{\pi_{ref}(o_i|q)}{\pi_{\theta}(o_i|q)} - 1,
\end{equation}
where $\pi_{ref}$ is a reference policy, $\epsilon$ and $\beta$ are hyper-parameters, and $A_i$ is the advantage, computed using a group of rewards $\{r_1, r_2, \ldots, r_G\}$ corresponding to the outputs within each group:
\begin{equation}
    A_i = \frac{r_i - {\mathrm mean(\{r_1, r_2, \cdots, r_G\})}}{{\mathrm std(\{r_1, r_2, \cdots, r_G\})}}.
\end{equation}

We give a comparison of GRPO and PPO in Supplementary \ref{section:grpo}.
To train DeepSeek-R1-Zero, we set the learning rate to 3e-6, the KL coefficient to 0.001, and the sampling temperature to 1 for rollout. For each question, we sample 16 outputs with a maximum length of 32,768 tokens before the 8.2k step and 65,536 tokens afterward. As a result, both the performance and response length of DeepSeek-R1-Zero exhibit a significant jump at the 8.2k step, with training continuing for a total of 10,400 steps, corresponding to 1.6 training epochs. Each training step consists of 32 unique questions, resulting in a training batch size of 512. Every 400 steps, we replace the reference model with the latest policy model. To accelerate training, each rollout generates 8,192 outputs, which are randomly split into 16 mini-batches and trained for only a single inner epoch.
\begin{table}[h]
    \centering
    \small
    \caption{Template for DeepSeek-R1-Zero. \textcolor{red}{prompt} will be replaced with the specific reasoning question during training.}
    \begin{tabularx}{\textwidth}{X}
    \toprule
    A conversation between User and Assistant. The user asks a question, and the Assistant solves it. 
    The assistant first thinks about the reasoning process in the mind and then provides the user with the answer.
    The reasoning process and answer are enclosed within \texttt{<think>...</think>} and \texttt{<answer>...</answer>} tags, respectively, i.e., \texttt{<think>} reasoning process here \texttt{</think>} \texttt{<answer>} answer here \texttt{</answer>}. 
    User: \textcolor{red}{prompt}. Assistant: \\
    \bottomrule
    \end{tabularx}
    \label{tab:r0_template}
\end{table}

Our high-performance RL infrastructure is described in Supplementary \ref{sec:infra}, ensuring scalable and efficient training.
\subsection{Reward Design} 
\label{section:reward}
The reward is the source of the training signal, which decides the direction of RL optimization.
For DeepSeek-R1-Zero,  we employ rule-based rewards to 
deliver precise feedback for data in mathematical, coding, and logical reasoning domains. Our rule-based reward system mainly consists of two types of rewards: accuracy rewards and format rewards.

\textbf{Accuracy rewards} evaluate whether the response is correct. For example, in the case of math problems with deterministic results, the model is required to provide the final answer in a specified format (e.g., within a box), enabling reliable rule-based verification of correctness. Similarly, for code competition prompts, a compiler can be utilized to evaluate the model's responses against a suite of predefined test cases, thereby generating objective feedback on correctness.

\textbf{Format rewards} complement the accuracy reward model by enforcing specific formatting requirements. In particular, the model is incentivized to encapsulate its reasoning process within designated tags, specifically `\texttt{<think>} ' and `\texttt{</think>} '. This ensures that the model's thought process is explicitly delineated, enhancing interpretability and facilitating subsequent analysis. 
\begin{equation}
    Reward_\text{rule} = Reward_\text{acc} + Reward_\text{format}
\end{equation}

The accuracy, reward and format reward are combined with the same weight. 
Notably, we abstain from applying neural reward models—whether outcome-based or process-based—to reasoning tasks. This decision is predicated on our observation that neural reward models are susceptible to reward hacking during large-scale reinforcement learning. Moreover, retraining such models necessitates substantial computational resources and introduces additional complexity into the training pipeline, thereby complicating the overall optimization process.



\subsection{Incentivize Reasoning Capability in LLMs}
Specifically, we apply the RL technique on the DeepSeek-V3 base to train DeepSeek-R1-Zero. During training, we design a straightforward template, to require DeepSeek-R1-Zero to first produce a reasoning process, followed by the final answer. We intentionally limit our constraints to this structural format, avoiding any content-specific biases to ensure that we can accurately observe the model's natural progression during the RL process.



\begin{figure}[t]
\centering
\subfigure{
\includegraphics[width=0.47\textwidth]{figures/plot_aime_with_maj.pdf}
}
\subfigure{
\includegraphics[width=0.481\textwidth]{figures/plot_length.pdf}
}
\caption{(a) AIME accuracy of DeepSeek-R1-Zero during training. AIME takes a mathematical problem as input and a number as output, illustrated in Table  \ref{tab:math_eval_format_example}.  Pass@1 and Cons@16 are described in Supplementary \ref{sec: evaluation_set}. The baseline is the average score achieved by human participants in the AIME competition.   (b) The average response length of DeepSeek-R1-Zero on the training set during the RL process. DeepSeek-R1-Zero naturally learns to solve reasoning tasks with more thinking time. Note that a training step refers to a single policy update operation.}
\label{fig:r1-zero}
\end{figure}


Figure \ref{fig:r1-zero}(a) depicts the performance trajectory of DeepSeek-R1-Zero on the AIME 2024 benchmark throughout the RL training process, where the average pass@1 score on AIME 2024 shows a significant increase, jumping from an initial 15.6\% to 77.9\%.
In addition, by leveraging the self-consistency decoding \citep{wangself}, the model's performance can be further improved, achieving an accuracy of 86.7\%. This performance significantly surpasses the average performance across all human competitors.
Besides the math competitions, as shown in Figure \ref{fig:dsv3_performance}, DeepSeek-R1-Zero also achieves remarkable performance in coding competitions and graduate-level biology, physics, and chemistry problems.
These results underscore the effectiveness of RL in enhancing the reasoning capabilities of large language models.

\begin{table}[t]

\caption{An interesting ``aha moment'' of an intermediate version of DeepSeek-R1-Zero. The model learns to rethink using an anthropomorphic tone. This is also an aha moment for us, allowing us to witness the power and beauty of reinforcement learning.}\centering
\begin{tabular}{l}
\toprule 
Question: If \(a > 1\), then the sum of the real solutions of \(\sqrt{a - \sqrt{a + x}} = x\) is equal to \\
\midrule
Response: \texttt{<think>} \\
To solve the equation \(\sqrt{a - \sqrt{a + x}} = x\), let's start by squaring both $\cdots$ \\
$\left(\sqrt{a - \sqrt{a + x}}\right)^2 = x^2 \implies a - \sqrt{a + x} = x^2.$ \\
Rearrange to isolate the inner square root term: \\
$(a - x^2)^2 = a + x \implies a^2 - 2ax^2 + (x^2)^2 = a + x \implies x^4 - 2ax^2 - x + (a^2 - a) = 0$ \\
\ldots \\
\textcolor{red}{Wait, wait. Wait. That's an aha moment I can flag here}. \\
Let's reevaluate this step-by-step to identify if the correct sum can be $\cdots$ \\
We started with the equation: \\
$\sqrt{a - \sqrt{a + x}} = x$ \\
First, let's square both sides: \\
$a - \sqrt{a + x} = x^2 \implies \sqrt{a + x} = a - x^2$ \\
Next, I could square both sides again, treating the equation: 
$\cdots$
\\
\ldots \\
\hline
\end{tabular}

\label{tab:aha_moment}
\end{table}

The self-evolution of DeepSeek-R1-Zero exemplifies how RL can autonomously enhance a model’s reasoning capabilities. 
 
As shown in Figure \ref{fig:r1-zero}(b), DeepSeek-R1-Zero exhibits a steady increase in thinking time throughout training, driven solely by intrinsic adaptation rather than external modifications. Leveraging long CoT, the model progressively refines its reasoning, generating hundreds to thousands of tokens to explore and improve its problem-solving strategies.  

The increase in thinking time fosters the autonomous development of sophisticated behaviors.
Specifically, DeepSeek-R1-Zero increasingly exhibits advanced reasoning strategies such as reflective reasoning and systematic exploration of alternative solutions (see Figure \ref{fig:r1-zero-evol}(a) in Supplementary \ref{app:evo_zero} for details), significantly boosting its performance on verifiable tasks like math and coding.
Notably, during training, DeepSeek-R1-Zero exhibits an ``aha moment'' (Table \ref{tab:aha_moment}), characterized by a sudden increase in the use of the word ``wait'' during reflections (see Figure \ref{fig:r1-zero-evol}(b) in Supplementary \ref{app:evo_zero} for details). This moment marks a distinct change in reasoning patterns and clearly shows the self-evolution process of DeepSeek-R1-Zero.

The self-evolution of DeepSeek-R1-Zero underscores the power and beauty of RL: rather than explicitly teaching the model how to solve a problem, we simply provide it with the right incentives, and it autonomously develops advanced problem-solving strategies.
This serves as a reminder of the potential of RL to unlock higher levels of capabilities in LLMs, paving the way for more autonomous and adaptive models in the future.


\section{DeepSeek-R1}

\label{sec:ds-r1}
\begin{figure}
    \centering
    \includegraphics[width=\linewidth]{figures/R1Pipieline_v0603.pdf}
    \caption{The multi-stage pipeline of DeepSeek-R1. A detailed background on DeepSeek-V3 Base and DeepSeek-V3 is provided in Supplementary \ref{sec: background_v3}. The models DeepSeek-R1 Dev1, Dev2, and Dev3 represent intermediate checkpoints within this pipeline.}
    \label{fig:r1-pipeline}
\end{figure}
Although DeepSeek-R1-Zero exhibits strong reasoning capabilities, it faces several issues.  DeepSeek-R1-Zero struggles with challenges like poor readability, and language mixing, as DeepSeek-V3-Base is trained on multiple languages, especially English and Chinese. 
To address these issues, we develop DeepSeek-R1, whose pipeline is illustrated in Figure \ref{fig:r1-pipeline}.

In the initial stage, we collect thousands of cold-start data that exhibits a conversational, human-aligned thinking process. RL training is then applied to improve the model performance with the conversational thinking process and language consistency. Subsequently, we apply rejection sampling and SFT once more. This stage incorporates both reasoning and non-reasoning datasets into the SFT process, enabling the model to not only excel in reasoning tasks but also demonstrate advanced writing capabilities. To further align the model with human preferences, we implement a secondary RL stage designed to enhance the model's helpfulness and harmlessness while simultaneously refining its reasoning capabilities. 

The remainder of this section details the key components of this pipeline: Section \ref{sec: model_reward} introduces the Reward Model utilized in our RL stages, and Section \ref{sec:rl-detail} elaborates on the specific training methodologies and implementation details. Data we used in this stage is detailed in Supplementary \ref{sec:rej_sample}.

\subsection{Model-based Rewards}\label{sec: model_reward}

For general data, we resort to reward models to capture human preferences in complex and nuanced scenarios. We build
upon the DeepSeek-V3 pipeline and adopt a similar distribution of preference pairs and training prompts. For helpfulness, we focus exclusively on the final summary, ensuring that the
assessment emphasizes the utility and relevance of the response to the user while minimizing
interference with the underlying reasoning process. For harmlessness, we evaluate the entire
response of the model, including both the reasoning process and the summary, to identify and
mitigate any potential risks, biases, or harmful content that may arise during the generation
process.

\paragraph{Helpful Reward Model}


Regarding helpful reward model training, we first generate preference pairs by prompting DeepSeek-V3 using the arena-hard prompt format, listed in Supplementary \ref{section:reward}, where each pair consists of a user query along with two candidate responses. For each preference pair, we query DeepSeek-V3 four times, randomly assigning the responses as either Response A or Response B to mitigate positional bias. The final preference score is determined by averaging the four independent judgments, retaining only those pairs where the score difference ($\Delta$) exceeds 1 to ensure meaningful distinctions. Additionally, to minimize length-related biases, we ensure that the chosen and rejected responses of the whole dataset have comparable lengths. In total, we curated 66,000 data pairs for training the reward model. The prompts used in this dataset are all non-reasoning questions and are sourced either from publicly available open-source datasets or from users who have explicitly consented to share their data for the purpose of model improvement. The architecture of our reward model is consistent with that of DeepSeek-R1, with the addition of a reward head designed to predict scalar preference scores. 

\begin{equation} 
    Reward_{helpful} = RM_{helpful}(Response_A, Response_B)
\end{equation}

The helpful reward models were trained with a batch size of 256, a learning rate of 6e-6, and for a single epoch over the training dataset. The maximum sequence length during training is set to 8192 tokens, whereas no explicit limit is imposed during reward model inference. 



\paragraph{Safety Reward Model}
To assess and improve model safety, we curated a dataset of 106,000 prompts with model-generated responses annotated as ``safe" or ``unsafe" according to predefined safety guidelines. Unlike the pairwise loss employed in the helpfulness reward model, the safety reward model was trained using a point-wise methodology to distinguish between safe and unsafe responses. The training hyperparameters are the same as the helpful reward model. 

\begin{equation}
    Reward_{safety} = RM_{safety}(Response)
\end{equation}
For general queries, each instance is categorized as belonging to either the safety dataset or the helpfulness dataset. The general reward, $Reward_{General}$, assigned to each query corresponds to the respective reward defined within the associated dataset. 


\subsection{Training Details}

\label{sec:rl-detail}
\subsubsection{Training Details of the First RL Stage} \label{sec: r1-1-rl}

In the first stage of RL, we set the learning rate to 3e-6, the KL coefficient to 0.001, the GRPO clip ratio $\epsilon$ to 10, and the sampling temperature to 1 for rollout. For each question, we sample 16 outputs with a maximum length of 32,768. Each training step consists of 32 unique questions, resulting in a training batch size of 512 per step. Every 400 steps, we replace the reference model with the latest policy model. To accelerate training, each rollout generates 8,192 outputs, which are randomly split into 16 minibatches and trained for only a single inner epoch. However, to mitigate the issue of language mixing, we introduce a language consistency reward during RL training, which is calculated as the proportion of target language words in the CoT. 
\begin{equation}
    Reward_{language} = \frac{Num(Words_{target})}{Num(Words)}
\end{equation}
Although ablation experiments in Supplementary \ref{sec:LC_reward} show that such alignment results in a slight degradation in the model's performance, this reward aligns with human preferences, making it more readable. We apply the language consistency reward to both reasoning and non-reasoning data by directly adding it to the final reward. 

Note that the clip ratio plays a crucial role in training.  A lower value can lead to the truncation of gradients for a significant number of tokens, thereby degrading the model's performance, while a higher value may cause instability during training.



\subsubsection{Training Details of the Second RL Stage }\label{sec: r1-2-rl}
Specifically, we train the model using a combination of reward signals and diverse prompt distributions. For reasoning data, we follow the methodology outlined in DeepSeek-R1-Zero, which employs rule-based rewards to guide learning in mathematical, coding, and logical reasoning domains. During the training process, we observe that CoT often exhibits language mixing, particularly when RL prompts involve multiple languages. 
For general data, we utilize reward models to guide training. Ultimately, the integration of reward signals with diverse data distributions enables us to develop a model that not only excels in reasoning but also prioritizes helpfulness and harmlessness. Given a batch of data, the reward can be formulated as 

\begin{align}
Reward &= Reward_{\text{reasoning}} + Reward_{\text{general}} + Reward_{\text{language}}\\
\text{where, } Reward_{\text{reasoning}} &= Reward_{\text{rule}}\\
Reward_{\text{general}} &= Reward_{\text{reward\_model}} + Reward_{\text{format}}
\end{align}

The second stage of RL retains most of the parameters from the first stage, with the key difference being a reduced temperature of 0.7, as we find that higher temperatures in this stage lead to incoherent generation. The stage comprises a total of 1,700 training steps, during which general instruction data and preference-based rewards are incorporated exclusively in the final 400 steps. We find that more training steps with the model based preference reward signal may lead to reward hacking, which is documented in Supplementary \ref{sec:reward_hack}. The total training cost is listed in Supplementary \ref{sec:cost}.


\begin{table}[h]
    \centering
    \small
    \setlength{\tabcolsep}{1.9pt}
    \caption{ Experimental results at each stage of DeepSeek-R1. Numbers in bold denote the performance is statistically significant  (t$-$test with $p < 0.01$).}
    \begin{tabular}{@{}c l | c c  c  c |c c@{}}
    \toprule
    % & \multirow{2}{*}{\centering \textbf{Benchmark {\tiny (Metric)}}} & \textbf{DeepSeek-R1} & \textbf{DeepSeek-R1}  & \textbf{DeepSeek-R1}& \textbf{DeepSeek-R1} &  \textbf{DeepSeek}\\
    % & & \textbf{Zero}   & \textbf{Dev1}  & \textbf{Dev2} & \textbf{Dev3} & \textbf{R1} \\
     & \textbf{Benchmark {\tiny (Metric)}} & \textbf{R1-Zero} & \textbf{R1-Dev1}  & \textbf{R1-Dev2}& \textbf{R1-Dev3} & \textbf{R1} \\
    \midrule
    \multirow{10}{*}{English}& MMLU {\tiny (EM)} & 88.8 & 89.1 & \textbf{91.2} & 91.0  &   90.8\\
     & MMLU-Redux {\tiny (EM)}& 85.6 &90.0 & 93.0 & 93.1 &  92.9 \\
    & MMLU-Pro {\tiny (EM)}  & 68.9 & 74.1 & 83.8 & 83.1 &  \textbf{84.0} \\
    & DROP {\tiny (3-shot F1)}  & 89.1 &89.8 & 91.1 & 88.7  &  \textbf{92.2}\\
    & IF-Eval {\tiny (Prompt Strict)}  & 46.6& 71.7 & 72.0 & 78.1 & \textbf{83.3} \\
    & GPQA Diamond {\tiny (Pass@1)}  & \textbf{75.8} &66.1 & 70.7 &  71.2 &  71.5&  \\
    & SimpleQA {\tiny (Correct)}   & 30.3 & 17.8 & 28.2 & 24.9  &  {30.1} \\
     & FRAMES {\tiny (Acc.)}    & 82.3 & 78.5 & 81.8 & 81.9 &   \textbf{82.5}\\
      & AlpacaEval2.0 {\tiny (LC-winrate)}   & 24.7 & 50.1 & 55.8 & 62.1  &   \textbf{87.6}\\
       & ArenaHard {\tiny (GPT-4-1106)}   & 53.6 & 77.0 & 73.2 & 75.6 &   \textbf{92.3}\\
    \midrule
    \multirow{7}{*}{Code} & LiveCodeBench {\tiny (Pass@1-COT)}   & 50.0 & 57.5&63.5 & 64.6  &   \textbf{65.9} \\
    & Codeforces {\tiny (Percentile)}  & 80.4 & 84.5& 90.5& 92.1  &   \textbf{96.3} \\
    & Codeforces {\tiny (Rating)}  & 1444 &1534 & 1687& 1746  &   \textbf{2029} \\
    & SWE Verified {\tiny (Resolved)}  & 43.2 & 39.6&44.6 &  45.6 &   \textbf{49.2}\\
    & Aider-Polyglot {\tiny (Acc.)}   & 12.2 &6.7 & 25.6&  44.8 &   \textbf{53.3}\\
    \midrule
    \multirow{3}{*}{Math} & AIME 2024 {\tiny (Pass@1)} & 77.9 & 59.0 & 74.0 & 78.1 &   \textbf{79.8} \\
    & MATH-500 {\tiny (Pass@1)}   & 95.9 & 94.2 & 95.9 &  95.4 &  \textbf{97.3} \\
    & CNMO 2024 {\tiny (Pass@1)}  & \textbf{88.1} &58.0 & 73.9 & 77.3  &  {78.8} \\
    \midrule
    \multirow{3}{*}{Chinese} & CLUEWSC {\tiny (EM)}  & 93.1 & 92.8 & 92.6 &  91.6 &   {92.8}\\
    & C-Eval {\tiny (EM)}   & \textbf{92.8} & 85.7 & 91.9 &  86.4 &  {91.8}\\
     & C-SimpleQA {\tiny (Correct)} & 66.4 & 58.8 & 64.2 & 66.9 &  63.7 \\
    \bottomrule
    \end{tabular}
    
    \label{tab:stage_r1}
\end{table}



\section{Experiment}

 We evaluate our models on MMLU \citep{mmlu}, MMLU-Redux \citep{mmlu_redux}, MMLU-Pro \citep{mmlu_pro}, C-Eval \citep{ceval}, and CMMLU \citep{cmmlu}, IFEval~\citep{IFeval}, FRAMES~\citep{frames}, GPQA Diamond ~\citep{gpqa}, SimpleQA~\citep{simpleqa}, C-SimpleQA~\citep{csimpleqa}, SWE-Bench Verified~\citep{swe_verified}, Aider \citep{aider}, LiveCodeBench~\citep{livecodebench} (2024-08 -- 2025-01), Codeforces \citep{codeforces}, Chinese National High School Mathematics Olympiad (CNMO 2024) \citep{cnmo}, and American Invitational Mathematics Examination 2024 (AIME 2024)~\citep{AIME2024}. The details of these benchmarks are listed in Supplementary \ref{sec:eval}. 
% 

Table \ref{tab:stage_r1} summarizes the performance of DeepSeek-R1 across multiple developmental stages, as outlined in Figure \ref{fig:r1-pipeline}. A comparison between DeepSeek-R1-Zero and DeepSeek-R1 Dev1 reveals substantial improvements in instruction-following, as evidenced by higher scores on the IF-Eval and ArenaHard benchmarks. However, due to the limited size of the cold-start dataset, Dev1 exhibits a partial degradation in reasoning performance compared to DeepSeek-R1-Zero, most notably on the AIME benchmark. In contrast, DeepSeek-R1 Dev2 demonstrates marked performance enhancements on benchmarks that require advanced reasoning skills, including those focused on code generation, mathematical problem solving, and STEM-related tasks. Benchmarks targeting general-purpose tasks, such as AlpacaEval 2.0, show marginal improvement. These results suggest that reasoning-oriented RL considerably enhances reasoning capabilities while exerting limited influence on user preference-oriented benchmarks.

 DeepSeek-R1 Dev3 integrates both reasoning and non-reasoning datasets into the SFT pipeline, thereby enhancing the model’s proficiency in both reasoning and general language generation tasks. Compared to Dev2, DeepSeek-R1 Dev3 achieves notable performance improvements on AlpacaEval 2.0 and Aider-Polyglot, attributable to the inclusion of large-scale non-reasoning corpora and code engineering datasets. Finally, comprehensive RL training on DeepSeek-R1 Dev3 using mixed reasoning-focused and general-purpose data produced the final DeepSeek-R1. Marginal improvements occurred in code and mathematics benchmarks, as substantial reasoning-specific RL was done in prior stages. The primary advancements in the final DeepSeek-R1 were in general instruction-following and user-preference benchmarks, with AlpacaEval 2.0 improving by 25\% and ArenaHard by 17\%.



    In addition, we compare DeepSeek-R1 with other models in Supplementary \ref{sec: benchmark}. Model safety evaluations are provided in Supplementary \ref{sec:safety_eval}. 
    A comprehensive analysis is provided in Supplementary \ref{sec:exp_ana}, including a comparison with DeepSeek-V3, performance evaluations on both fresh test sets, a breakdown of mathematical capabilities by category, and an investigation of test-time scaling behavior. Supplementary \ref{sec:distillation} shows that the strong reasoning capability can be transferred to smaller models. 



\section{Ethics and Safety Statement}
With the advancement in the reasoning capabilities of DeepSeek-R1, we deeply recognize the potential ethical risks. For example, R1 can be subject to jailbreak attacks, leading to the generation of dangerous content such as explosive manufacturing plans, while the enhanced reasoning capabilities enable the model to provide plans with better operational feasibility and executability.
 Besides, a public model is also vulnerable to further fine-tuning that could compromise inherent safety protections.
 
In Supplementary~\ref{sec:safety_eval}, we present a comprehensive safety report from multiple perspectives, including performance on open-source and in-house safety evaluation benchmarks, and safety levels across multiple languages and against jailbreak attacks. These comprehensive safety analyses conclude that the inherent safety level of the DeepSeek-R1 model, compared to other state-of-the-art models, is generally at a moderate level (comparable to GPT-4o (2024-05-13)). Besides, when coupled with the risk control system, the model's safety level is elevated to a superior standard.


\section{ Conclusion, Limitation, and Future Work}
We present DeepSeek-R1-Zero and DeepSeek-R1, which rely on large-scale RL to incentivize model reasoning behaviors. Our results demonstrate that pre-trained checkpoints inherently possess substantial potential for complex reasoning tasks. We believe that the key to unlocking this potential lies not in large-scale human annotation but in the provision of hard reasoning questions, a reliable verifier, and sufficient computational resources for reinforcement learning. Sophisticated reasoning behaviors, such as self-verification and reflection, appeared to emerge organically during the reinforcement learning process.

\label{sec:limit}
Even if DeepSeek-R1 achieves frontier results on reasoning benchmarks, it still faces several capability limitations, as outlined below:

 \textbf{Structure Output and Tool Use:}
Currently, the structural output capabilities of DeepSeek-R1 remain suboptimal compared to existing models. Moreover, DeepSeek-R1 cannot leverage tools, such as search engines and calculators, to improve the performance of output. However, as it is not hard to build an RL environment for structure output and tool use, we believe the issue will be addressed in the next version.
  
  \textbf{Token efficiency:} Unlike conventional test-time computation scaling approaches, such as majority voting or Monte Carlo Tree Search (MCTS), DeepSeek-R1 dynamically allocates computational resources during inference according to the complexity of the problem at hand. Specifically, it uses fewer tokens to solve simple tasks, while generating more tokens for complex tasks. Nevertheless, there remains room for further optimization in terms of token efficiency, as instances of excessive reasoning—manifested as overthinking—are still observed in response to simpler questions.
  
 \textbf{Language Mixing:}
DeepSeek-R1 is currently optimized for Chinese and English, which may result in language mixing issues when handling queries in other languages. For instance, DeepSeek-R1 might use English for reasoning and responses, even if the query is in a language other than English or Chinese. We aim to address this limitation in future updates. The limitation may be related to the base checkpoint, DeepSeek-V3-Base, mainly utilizes Chinese and English, so that it can achieve better results with the two languages in reasoning. 

 \textbf{Prompting Engineering:} When evaluating DeepSeek-R1, we observe that it is sensitive to prompts. Few-shot prompting consistently degrades its performance. Therefore, we recommend users directly describe the problem and specify the output format using a zero-shot setting for optimal results.
 
 \textbf{Software Engineering Tasks:}
Due to the long evaluation times, which impact the efficiency of the RL process, large-scale RL has not been applied extensively in software engineering tasks. As a result, DeepSeek-R1 has not demonstrated a huge improvement over DeepSeek-V3 on software engineering benchmarks. Future versions will address this by implementing rejection sampling on software engineering data or incorporating asynchronous evaluations during the RL process to improve efficiency.

Beyond specific capability limitations, the pure RL methodology itself also presents inherent challenges:

 \textbf{Reward Hacking:} The success of pure RL depends on reliable reward signals. In this study, we ensure reward reliability through a reasoning-domain rule-based reward model (RM). However, such dependable RMs are difficult to construct for certain tasks, such as writing. If the reward signal is assigned by a model instead of predefined rules, it becomes more susceptible to exploitation as training progresses, which means the policy model may find shortcuts to hack the reward model. Consequently, for complex tasks that cannot be effectively evaluated by a reliable reward model, scaling up pure RL methods remains an open challenge.

In this work, for tasks that cannot obtain a reliable signal, DeepSeek-R1 uses human annotation to create supervised data, and only conduct RL for hundreds of steps. We hope in the future, a robust reward model can be obtained to address such issues. 

With the advent of pure RL methods like DeepSeek-R1, the future holds immense potential for solving any task that can be effectively evaluated by a verifier, regardless of its complexity for humans. Machines equipped with such advanced RL techniques are poised to surpass human capabilities in these domains, driven by their ability to optimize performance iteratively through trial and error. However, challenges remain for tasks where constructing a reliable reward model is inherently difficult. In such cases, the lack of a robust feedback mechanism may hinder progress, suggesting that future research should focus on developing innovative approaches to define and refine reward structures for these complex, less verifiable problems.

 Furthermore, leveraging tools during the reasoning process holds significant promise. Whether it’s utilizing tools like compilers or search engines to retrieve or compute necessary information, or employing external tools—such as biological or chemical reagents, to validate final results in the real world, this integration of tool-augmented reasoning could dramatically enhance the scope and accuracy of machine-driven solutions.

%\bibliography{sn-bibliography}% common bib file
% \begin{thebibliography}{30}
% \providecommand{\natexlab}[1]{#1}
% \providecommand{\url}[1]{{#1}}
% \providecommand{\urlprefix}{URL }
% \providecommand{\doi}[1]{\url{https://doi.org/#1}}
% \providecommand{\eprint}[2][]{\url{#2}}
%  \bibcommenthead

% \bibitem[{Brown et~al(2020)Brown, Mann, Ryder, Subbiah, Kaplan, Dhariwal, Neelakantan, Shyam, Sastry, Askell, Agarwal, Herbert{-}Voss, Krueger, Henighan, Child, Ramesh, Ziegler, Wu, Winter, Hesse, Chen, Sigler, Litwin, Gray, Chess, Clark, Berner, McCandlish, Radford, Sutskever, and Amodei}]{gpt3}
% Brown TB, Mann B, Ryder N, et~al (2020) Language models are few-shot learners. In: Larochelle H, Ranzato M, Hadsell R, et~al (eds) Advances in Neural Information Processing Systems 33: Annual Conference on Neural Information Processing Systems 2020, NeurIPS 2020, December 6-12, 2020, virtual, \urlprefix\url{https://proceedings.neurips.cc/paper/2020/hash/1457c0d6bfcb4967418bfb8ac142f64a-Abstract.html}

% \bibitem[{Chung et~al(2024)Chung, Hou, Longpre, Zoph, Tay, Fedus, Li, Wang, Dehghani, Brahma, Webson, Gu, Dai, Suzgun, Chen, Chowdhery, Castro{-}Ros, Pellat, Robinson, Valter, Narang, Mishra, Yu, Zhao, Huang, Dai, Yu, Petrov, Chi, Dean, Devlin, Roberts, Zhou, Le, and Wei}]{chung2024scaling}
% Chung HW, Hou L, Longpre S, et~al (2024) Scaling instruction-finetuned language models. J Mach Learn Res 25:70:1--70:53. \urlprefix\url{https://jmlr.org/papers/v25/23-0870.html}

% \bibitem[{CMS(2024)}]{cnmo}
% CMS (2024) Chinese national high school mathematics olympiad. \urlprefix\url{https://www.cms.org.cn/Home/comp/comp/cid/12.html}

% \bibitem[{DeepSeek-AI(2024)}]{dsviii}
% DeepSeek-AI (2024) Deepseek-v3 technical report. arXiv preprint arXiv:241219437

% \bibitem[{Gauthier(2025)}]{aider}
% Gauthier P (2025) Aider {LLM} leaderboard. \urlprefix\url{https://aider.chat/docs/leaderboards/}

% \bibitem[{Gema et~al(2025)Gema, Leang, Hong, Devoto, Mancino, Saxena, He, Zhao, Du, Madani, Barale, McHardy, Harris, Kaddour, van Krieken, and Minervini}]{mmlu_redux}
% Gema AP, Leang JOJ, Hong G, et~al (2025) Are we done with mmlu? In: Chiruzzo L, Ritter A, Wang L (eds) Proceedings of the 2025 Conference of the Nations of the Americas Chapter of the Association for Computational Linguistics: Human Language Technologies, {NAACL} 2025 - Volume 1: Long Papers, Albuquerque, New Mexico, USA, April 29 - May 4, 2025. Association for Computational Linguistics, pp 5069--5096, \urlprefix\url{https://aclanthology.org/2025.naacl-long.262/}

% \bibitem[{He et~al(2024)He, Li, Liu, Tan, Wang, Huang, Bu, Guo, Hu, Zheng et~al}]{csimpleqa}
% He Y, Li S, Liu J, et~al (2024) Chinese simpleqa: A chinese factuality evaluation for large language models. arXiv preprint arXiv:241107140

% \bibitem[{Hendrycks et~al(2021)Hendrycks, Burns, Basart, Zou, Mazeika, Song, and Steinhardt}]{mmlu}
% Hendrycks D, Burns C, Basart S, et~al (2021) Measuring massive multitask language understanding. In: 9th International Conference on Learning Representations, {ICLR} 2021, Virtual Event, Austria, May 3-7, 2021. OpenReview.net, \urlprefix\url{https://openreview.net/forum?id=d7KBjmI3GmQ}

% \bibitem[{Huang et~al(2023)Huang, Bai, Zhu, Zhang, Zhang, Su, Liu, Lv, Zhang, Lei, Fu, Sun, and He}]{ceval}
% Huang Y, Bai Y, Zhu Z, et~al (2023) C-eval: {A} multi-level multi-discipline chinese evaluation suite for foundation models. In: Oh A, Naumann T, Globerson A, et~al (eds) Advances in Neural Information Processing Systems 36: Annual Conference on Neural Information Processing Systems 2023, NeurIPS 2023, New Orleans, LA, USA, December 10 - 16, 2023, \urlprefix\url{http://papers.nips.cc/paper\_files/paper/2023/hash/c6ec1844bec96d6d32ae95ae694e23d8-Abstract-Datasets\_and\_Benchmarks.html}

% \bibitem[{Jain et~al(2024)Jain, Han, Gu, Li, Yan, Zhang, Wang, Solar{-}Lezama, Sen, and Stoica}]{livecodebench}
% Jain N, Han K, Gu A, et~al (2024) Livecodebench: Holistic and contamination free evaluation of large language models for code. CoRR abs/2403.07974. \urlprefix\url{https://doi.org/10.48550/arXiv.2403.07974}

% \bibitem[{Kaplan et~al(2020)Kaplan, McCandlish, Henighan, Brown, Chess, Child, Gray, Radford, Wu, and Amodei}]{kaplan2020scaling}
% Kaplan J, McCandlish S, Henighan T, et~al (2020) Scaling laws for neural language models. arXiv preprint arXiv:200108361

% \bibitem[{Kojima et~al(2022)Kojima, Gu, Reid, Matsuo, and Iwasawa}]{kojima2022large}
% Kojima T, Gu SS, Reid M, et~al (2022) Large language models are zero-shot reasoners. In: Oh AH, Agarwal A, Belgrave D, et~al (eds) Advances in Neural Information Processing Systems, \urlprefix\url{https://openreview.net/forum?id=e2TBb5y0yFf}

% \bibitem[{Krishna et~al(2024)Krishna, Krishna, Mohananey, Schwarcz, Stambler, Upadhyay, and Faruqui}]{frames}
% Krishna S, Krishna K, Mohananey A, et~al (2024) Fact, fetch, and reason: {A} unified evaluation of retrieval-augmented generation. CoRR abs/2409.12941. \doi{10.48550/ARXIV.2409.12941}, \urlprefix\url{https://doi.org/10.48550/arXiv.2409.12941}, {\href{https://arxiv.org/abs/2409.12941}{{2409.12941}}}


% \bibitem[{Li et~al(2024)Li, Zhang, Koto, Yang, Zhao, Gong, Duan, and Baldwin}]{cmmlu}
% Li H, Zhang Y, Koto F, et~al (2024) {CMMLU:} measuring massive multitask language understanding in chinese. In: Ku L, Martins A, Srikumar V (eds) Findings of the Association for Computational Linguistics, {ACL} 2024, Bangkok, Thailand and virtual meeting, August 11-16, 2024. Association for Computational Linguistics, pp 11260--11285, \doi{10.18653/V1/2024.FINDINGS-ACL.671}, \urlprefix\url{https://doi.org/10.18653/v1/2024.findings-acl.671}

% \bibitem[{MAA(2024)}]{AIME2024}
% MAA (2024) American invitational mathematics examination - aime. In: American Invitational Mathematics Examination - AIME 2024, \urlprefix\url{https://maa.org/math-competitions/american-invitational-mathematics-examination-aime}

% \bibitem[{Mirzayanov(2025)}]{codeforces}
% Mirzayanov M (2025) Codeforces. \urlprefix\url{https://codeforces.com/}

% \bibitem[{OpenAI(2023)}]{gpt4}
% OpenAI (2023) {GPT4} technical report. arXiv preprint arXiv:230308774

% \bibitem[{OpenAI(2024{\natexlab{a}})}]{simpleqa}
% OpenAI (2024{\natexlab{a}}) Introducing {SimpleQA}. \urlprefix\url{https://openai.com/index/introducing-simpleqa/}

% \bibitem[{OpenAI(2024{\natexlab{b}})}]{swe_verified}
% OpenAI (2024{\natexlab{b}}) Introducing {SWE}-bench verified we’re releasing a human-validated subset of swe-bench that more. \urlprefix\url{https://openai.com/index/introducing-swe-bench-verified/}


% \bibitem[{Paszke et~al(2019)Paszke, Gross, Massa, Lerer, Bradbury, Chanan, Killeen, Lin, Gimelshein, Antiga, Desmaison, K{\"{o}}pf, Yang, DeVito, Raison, Tejani, Chilamkurthy, Steiner, Fang, Bai, and Chintala}]{pytorch}
% Paszke A, Gross S, Massa F, et~al (2019) Pytorch: An imperative style, high-performance deep learning library. In: Wallach HM, Larochelle H, Beygelzimer A, et~al (eds) Advances in Neural Information Processing Systems 32: Annual Conference on Neural Information Processing Systems 2019, NeurIPS 2019, December 8-14, 2019, Vancouver, BC, Canada, pp 8024--8035, \urlprefix\url{https://proceedings.neurips.cc/paper/2019/hash/bdbca288fee7f92f2bfa9f7012727740-Abstract.html}

% \bibitem[{Rein et~al(2023)Rein, Hou, Stickland, Petty, Pang, Dirani, Michael, and Bowman}]{gpqa}
% Rein D, Hou BL, Stickland AC, et~al (2023) {GPQA}: A graduate-level google-proof q\&a benchmark. arXiv preprint arXiv:231112022


% \bibitem[{Shao et~al(2024)Shao, Wang, Zhu, Xu, Song, Zhang, Li, Wu, and Guo}]{deepseekmath}
% Shao Z, Wang P, Zhu Q, et~al (2024) Deepseekmath: Pushing the limits of mathematical reasoning in open language models. arXiv preprint arXiv:240203300

% \bibitem[{Wang et~al(2023)Wang, Wei, Schuurmans, Le, Chi, Narang, Chowdhery, and Zhou}]{wangself}
% Wang X, Wei J, Schuurmans D, et~al (2023) Self-consistency improves chain of thought reasoning in language models. In: The Eleventh International Conference on Learning Representations, {ICLR} 2023, Kigali, Rwanda, May 1-5, 2023. OpenReview.net, \urlprefix\url{https://openreview.net/forum?id=1PL1NIMMrw}

% \bibitem[{Wang et~al(2024)Wang, Ma, Zhang, Ni, Chandra, Guo, Ren, Arulraj, He, Jiang, Li, Ku, Wang, Zhuang, Fan, Yue, and Chen}]{mmlu_pro}
% Wang Y, Ma X, Zhang G, et~al (2024) Mmlu-pro: {A} more robust and challenging multi-task language understanding benchmark. In: Globersons A, Mackey L, Belgrave D, et~al (eds) Advances in Neural Information Processing Systems 38: Annual Conference on Neural Information Processing Systems 2024, NeurIPS 2024, Vancouver, BC, Canada, December 10 - 15, 2024, \urlprefix\url{http://papers.nips.cc/paper\_files/paper/2024/hash/ad236edc564f3e3156e1b2feafb99a24-Abstract-Datasets\_and\_Benchmarks\_Track.html}

% \bibitem[{Wei et~al(2022{\natexlab{a}})Wei, Tay, Bommasani, Raffel, Zoph, Borgeaud, Yogatama, Bosma, Zhou, Metzler, Chi, Hashimoto, Vinyals, Liang, Dean, and Fedus}]{wei2022emergent}
% Wei J, Tay Y, Bommasani R, et~al (2022{\natexlab{a}}) Emergent abilities of large language models. Trans Mach Learn Res 2022. \urlprefix\url{https://openreview.net/forum?id=yzkSU5zdwD}

% \bibitem[{Wei et~al(2022{\natexlab{b}})Wei, Wang, Schuurmans, Bosma, Ichter, Xia, Chi, Le, and Zhou}]{wei2022chain}
% Wei J, Wang X, Schuurmans D, et~al (2022{\natexlab{b}}) Chain-of-thought prompting elicits reasoning in large language models. In: Koyejo S, Mohamed S, Agarwal A, et~al (eds) Advances in Neural Information Processing Systems 35: Annual Conference on Neural Information Processing Systems 2022, NeurIPS 2022, New Orleans, LA, USA, November 28 - December 9, 2022, \urlprefix\url{http://papers.nips.cc/paper\_files/paper/2022/hash/9d5609613524ecf4f15af0f7b31abca4-Abstract-Conference.html}

% \bibitem[{Zhou et~al(2023)Zhou, Lu, Mishra, Brahma, Basu, Luan, Zhou, and Hou}]{IFeval}
% Zhou J, Lu T, Mishra S, et~al (2023) Instruction-following evaluation for large language models. arXiv preprint arXiv:231107911

% \end{thebibliography}


\section{Author List}
The list of authors is organized by contribution role, with individuals listed alphabetically by their first name within each category. Authors marked with an asterisk (*) are no longer affiliated with our team.


\textbf{Core Contributors}:
Daya Guo,
Dejian Yang,
Haowei Zhang,
Junxiao Song,
Peiyi Wang,
Qihao Zhu,
Runxin Xu,
Ruoyu Zhang,
Shirong Ma,
Xiao Bi,
Xiaokang Zhang,
Xingkai Yu,
Yu Wu,
Z.F. Wu,
Zhibin Gou,
Zhihong Shao,
Zhuoshu Li,
Ziyi Gao,


\textbf{Contributions of the Core Authors: }
Peiyi Wang and Daya Guo jointly demonstrated that outcome-based RL induces the emergence of long Chain-of-Thought patterns in LLMs, achieving breakthrough reasoning capabilities. They contributed equally to the creation of R1-Zero, and their work laid the foundation for R1. Daya Guo also contributed to the RL training stability of MOE models. Junxiao Song proposed the GRPO algorithm, implemented the initial version, and introduced rule-based rewards for math tasks. The GRPO algorithm was subsequently refined by Peiyi Wang and Runxin Xu. Zhibin Gou proposed a large PPO clipping strategy to enhance GRPO performance,  demonstrating its significance alongside Zhihong Shao and Junxiao Song.
Regarding data iteration, reward design, and evaluation, specific teams led efforts across different domains: Qihao Zhu, Z.F. Wu, and Dejian Yang focused on code tasks; Zhihong Shao, Zhibin Gou, and Junxiao Song focused on math tasks; and Peiyi Wang, Ruoyu Zhang, Runxin Xu, and Yu Wu led efforts for other reasoning and general tasks. Additionally, Qihao Zhu and Zhihong Shao contributed to the data selection strategy for RL training, while Zhuoshu Li and Yu Wu co-led the data labeling efforts for the entire project.
On the system side, Xiao Bi, Xingkai Yu, Shirong Ma, Xiaokang Zhang, Haowei Zhang, and Ziyi Gao implemented the RL pipeline, optimizing system efficiency and addressing stability issues in large-scale training. Finally, Zhibin Gou, Daya Guo, and Ruoyu Zhang oversaw the final training phase and monitored the model training dynamics. Zhibin Gou led the development of the R1-distill series.


\textbf{ Contributors}:
 Aixin Liu,
 Bing Xue,
 Bingxuan Wang,
 Bochao Wu,
 Bei Feng,
 Chengda Lu,
 Chenggang Zhao,
 Chengqi Deng,
 Chong Ruan,
 Damai Dai,
 Deli Chen,
 Dongjie Ji,
 Erhang Li,
 Fangyun Lin,
 Fucong Dai,
 Fuli Luo*,
 Guangbo Hao,
 Guanting Chen,
 Guowei Li,
 H. Zhang,
 Hanwei Xu,
 Honghui Ding,
 Huazuo Gao,
 Hui Qu,
 Hui Li,
 Jianzhong Guo,
 Jiashi Li,
 Jingchang Chen,
 Jingyang Yuan,
 Jinhao Tu,
 Junjie Qiu,
 Junlong Li,
 J.L. Cai,
 Jiaqi Ni,
 Jian Liang,
 Jin Chen,
 Kai Dong,
 Kai Hu*,
 Kaichao You,
 Kaige Gao,
 Kang Guan,
 Kexin Huang,
 Kuai Yu,
 Lean Wang,
 Lecong Zhang,
 Liang Zhao,
 Litong Wang,
 Liyue Zhang,
 Lei Xu,
 Leyi Xia,
 Mingchuan Zhang,
 Minghua Zhang,
 Minghui Tang,
 Mingxu Zhou,
 Meng Li,
 Miaojun Wang,
 Mingming Li,
 Ning Tian,
 Panpan Huang,
 Peng Zhang,
 Qiancheng Wang,
 Qinyu Chen,
 Qiushi Du,
 Ruiqi Ge,
 Ruisong Zhang,
 Ruizhe Pan,
 Runji Wang,
 R.J. Chen,
 R.L. Jin,
 Ruyi Chen,
 Shanghao Lu,
 Shangyan Zhou,
 Shanhuang Chen,
 Shengfeng Ye,
 Shiyu Wang,
 Shuiping Yu,
 Shunfeng Zhou,
 Shuting Pan,
 S.S. Li,
 Shuang Zhou,
 Shaoqing Wu,
 Shengfeng Ye,
 Tao Yun,
 Tian Pei,
 Tianyu Sun,
 T. Wang,
 Wangding Zeng,
 Wen Liu,
 Wenfeng Liang,
 Wenjun Gao,
 Wenqin Yu*,
 Wentao Zhang,
 W.L. Xiao,
 Wei An,
 Xiaodong Liu,
 Xiaohan Wang,
 Xiaokang Chen,
 Xiaotao Nie,
 Xin Cheng,
 Xin Liu,
 Xin Xie,
 Xingchao Liu,
 Xinyu Yang,
 Xinyuan Li,
 Xuecheng Su,
 Xuheng Lin,
 X.Q. Li,
 Xiangyue Jin,
 Xiaojin Shen,
 Xiaosha Chen,
 Xiaowen Sun,
 Xiaoxiang Wang,
 Xinnan Song,
 Xinyi Zhou,
 Xianzu Wang,
 Xinxia Shan,
 Y.K. Li,
 Y.Q. Wang,
 Y.X. Wei,
 Yang Zhang,
 Yanhong Xu,
 Yao Li,
 Yao Zhao,
 Yaofeng Sun,
 Yaohui Wang,
 Yi Yu,
 Yichao Zhang,
 Yifan Shi,
 Yiliang Xiong,
 Ying He,
 Yishi Piao,
 Yisong Wang,
 Yixuan Tan,
 Yiyang Ma*,
 Yiyuan Liu,
 Yongqiang Guo,
 Yuan Ou,
 Yuduan Wang,
 Yue Gong,
 Yuheng Zou,
 Yujia He,
 Yunfan Xiong,
 Yuxiang Luo,
 Yuxiang You,
 Yuxuan Liu,
 Yuyang Zhou,
 Y.X. Zhu,
 Yanping Huang,
 Yaohui Li,
 Yi Zheng,
 Yuchen Zhu,
 Yunxian Ma,
 Ying Tang,
 Yukun Zha,
 Yuting Yan,
 Z.Z. Ren,
 Zehui Ren,
 Zhangli Sha,
 Zhe Fu,
 Zhean Xu,
 Zhenda Xie,
 Zhengyan Zhang,
 Zhewen Hao,
 Zhicheng Ma,
 Zhigang Yan,
 Zhiyu Wu,
 Zihui Gu,
 Zijia Zhu,
 Zijun Liu*,
 Zilin Li,
 Ziwei Xie,
 Ziyang Song,
 Zizheng Pan,
 Zhen Huang,
 Zhipeng Xu,
 Zhongyu Zhang,
 Zhen Zhang,

% \section{Methods}




% \section{Data Availability}
% We provide the data samples used in our rejection sampling and RL prompts at \url{https://github.com/deepseek-ai/DeepSeek-R1} (DOI: 10.5281/zenodo.15753193). Comprehensive statistics and details of our complete data-generation methodology are presented in Supplementary \ref{sec:rej_sample}.

% \section{Code Availability}
% Trained weights  of DeepSeek-R1-Zero and DeepSeek-R1 are available under an MIT license at \url{https://github.com/deepseek-ai/DeepSeek-R1} (DOI: 10.5281/zenodo.15753193).  The inference script is released at \url{https://github.com/deepseek-ai/DeepSeek-V3} (DOI: 10.5281/zenodo.15753347). Neural networks were developed with PyTorch \citep{pytorch}, and the distributed framework is based on our internal framework HAI-LLM {\url{https://www.high-flyer.cn/en/blog/hai-llm}}. The inference framework is based on vLLM  \citep{vllm}. 

% \begin{thebibliography}{30}
% \providecommand{\natexlab}[1]{#1}
% \providecommand{\url}[1]{{#1}}
% \providecommand{\urlprefix}{URL }
% \providecommand{\doi}[1]{\url{https://doi.org/#1}}
% \providecommand{\eprint}[2][]{\url{#2}}
%  \bibcommenthead
 
% \bibitem[{Shao et~al(2024)Shao, Wang, Zhu, Xu, Song, Zhang, Li, Wu, and Guo}]{deepseekmath}
% Shao Z, Wang P, Zhu Q, et~al (2024) Deepseekmath: Pushing the limits of mathematical reasoning in open language models. arXiv preprint arXiv:240203300

% \bibitem[{Schulman et~al(2017)Schulman, Wolski, Dhariwal, Radford, and Klimov}]{schulman2017proximal}
% Schulman J, Wolski F, Dhariwal P, et~al (2017) Proximal policy optimization algorithms. arXiv preprint arXiv:170706347

% \bibitem[{Ouyang et~al(2022)Ouyang, Wu, Jiang, Almeida, Wainwright, Mishkin, Zhang, Agarwal, Slama, Ray, Schulman, Hilton, Kelton, Miller, Simens, Askell, Welinder, Christiano, Leike, and Lowe}]{ouyang2022training}
% Ouyang L, Wu J, Jiang X, et~al (2022) Training language models to follow instructions with human feedback. In: Koyejo S, Mohamed S, Agarwal A, et~al (eds) Advances in Neural Information Processing Systems 35: Annual Conference on Neural Information Processing Systems 2022, NeurIPS 2022, New Orleans, LA, USA, November 28 - December 9, 2022, \urlprefix\url{http://papers.nips.cc/paper\_files/paper/2022/hash/b1efde53be364a73914f58805a001731-Abstract-Conference.html}

% \bibitem[{Kwon et~al(2023)Kwon, Li, Zhuang, Sheng, Zheng, Yu, Gonzalez, Zhang, and Stoica}]{vllm}
% Kwon W, Li Z, Zhuang S, et~al (2023) Efficient memory management for large language model serving with pagedattention. In: Proceedings of the ACM SIGOPS 29th Symposium on Operating Systems Principles



% \end{thebibliography}


\newpage
\appendix
\section*{Appendix}
\section{Background}
\subsection{DeepSeek-V3} \label{sec: background_v3}
DeepSeek V3 \citep{dsviii} is an advanced open-source  LLM developed by DeepSeek. Released in December 2024, DeepSeek V3 represents a significant leap forward in AI innovation, designed to rival leading models like OpenAI's GPT-4 and Meta's Llama 3.1, while maintaining remarkable cost efficiency and performance.
Built on a Mixture-of-Experts (MoE) architecture, DeepSeek V3 has 671 billion total parameters, with 37 billion activated per token, optimizing both efficiency and capability. It was pre-trained on an expansive dataset of 14.8 trillion high-quality, diverse tokens, followed by supervised fine-tuning and reinforcement learning to enhance its abilities across various domains. The model incorporates innovative features like Multi-head Latent Attention (MLA) \citep{dsvii} for efficient inference, an auxiliary-loss-free load-balancing strategy, and Multi-Token Prediction (MTP) \citep{meta_mtp} to boost performance, particularly in tasks like mathematics and coding.

For the training data of \dsviii{}-Base, we exclusively use plain web pages and e-books, without incorporating any synthetic data. However, we have observed that some web pages contain a significant number of OpenAI-model-generated answers, which may lead the base model to acquire knowledge from other powerful models indirectly. However, we did not intentionally include synthetic data generated by OpenAI during the pre-training cooldown phase; all data used in this phase were naturally occurring and collected through web crawling. The pre-training dataset contains a substantial amount of mathematical and code-related content, indicating that \dsviii{}-Base has been exposed to a significant volume of reasoning trace data. This extensive exposure equips the model with the capability to generate plausible solution candidates, from which reinforcement learning can effectively identify and optimize high-quality outputs.
We did the data contamination in pre-training as described in Appendix \ref{sec:conta}. The training data of DeepSeek-V3 base are mostly Chinese and English, which might be the cause for DeepSeek-R1-Zero language mixing when the language consistent reward is absent. 

In this paper, we use the notation DeepSeek-V3-Base as the base model, DeepSeek-V3 as the instructed model. Notably, \dsri{} and \dsro{} are trained on top of DeepSeek-V3-Base and  \dsri{} leverages non-reasoning data from DeepSeek-V3 SFT data. DeepSeek-R1-Dev1, DeepSeek-R1-Dev2, DeepSeek-R1-Dev3 are  intermediate checkpoints of \dsri{}. 


\subsection{Conventional Post-Training Paradigm}
Post-training has emerged as an essential step in refining pre-trained LLMs to meet specific performance goals and align with human expectations.  A widely adopted two-stage post-training framework is SFT  followed by RL \citep{ouyang2022training}. 

Supervised Fine-Tuning refines a pre-trained LLM by training it on a curated dataset of input-output pairs tailored to specific tasks. The process employs a supervised learning objective, typically minimizing cross-entropy loss between the model’s predictions and labeled ground truth \citep{gpt3}. For instance, in conversational applications, SFT might utilize dialogue datasets where desired responses are explicitly provided, enabling the model to adapt its outputs to predefined standards \citep{gpt2}. SFT offers several compelling benefits. First, it achieves precise task alignment by leveraging high-quality examples, allowing the model to excel in domains such as customer support or technical documentation \citep{gpt2}. Second, its reliance on pre-trained weights ensures computational efficiency, requiring fewer resources than training from scratch. Finally, the use of explicit input-output mappings enhances interpretability, as the model’s learning process is directly tied to observable data, minimizing the risk of erratic behavior  \citep{ouyang2022training}. Despite its strengths, the performance of SFT hinges on the quality and diversity of the training dataset; narrow or biased data can impair the model’s ability to generalize to novel contexts \citep{gpt3}. Additionally, SFT’s static nature—optimizing for fixed outputs—may fail to capture evolving human preferences or nuanced objectives. The labor-intensive process of curating high-quality datasets further complicates its scalability, as errors or inconsistencies in the data can propagate into the model’s behavior \citep{ouyang2022training}.

Following SFT, Reinforcement Learning further refines the LLM by optimizing its outputs against a reward signal. In this stage, the model interacts with an environment—often a reward model trained on human feedback—and adjusts its behavior to maximize cumulative rewards. A prominent instantiation of this approach is Reinforcement Learning from Human Feedback (RLHF), where the reward function encodes human preferences \citep{DBLP:conf/nips/ChristianoLBMLA17}. RL thus shifts the focus from static supervision to dynamic optimization. Notably, RL reduces the need for extensive annotated resources; while SFT demands a fully labeled dataset for every input-output pair, RL can operate with a smaller set of human evaluations or a trained reward model, even rule-based reward model, significantly lowering the annotation burden.




The sequential application of SFT and RL combines their complementary strengths. SFT establishes a robust, task-specific baseline by grounding the model in curated examples, while RL refines this foundation to align with broader, human-centric objectives \citep{ouyang2022training}. For example, SFT might ensure grammatical accuracy in a dialogue system, while RL optimizes for engagement and brevity, as demonstrated in the development of InstructGPT \citep{ouyang2022training}. This hybrid approach has proven effective in producing models that are both precise and adaptable.


In this study, we demonstrate that the SFT stage may impede a model’s ability to explore and develop effective reasoning strategies. This limitation arises because human-provided responses, which serve as targets during SFT, are not always optimal for model learning; they often omit critical reasoning components such as explicit reflection and verification steps. To address this, \dsro{} enables direct exploration of reasoning patterns by the model itself, independent of human priors. The reasoning trajectories discovered through this self-exploration are subsequently distilled and used to train other models, thereby promoting the acquisition of more robust and generalizable reasoning capabilities.




\subsection{A Comparison of GRPO and PPO} 
\label{section:grpo}

\begin{figure}
    \centering
    \includegraphics[width=\linewidth]{figures/GRPO.pdf}
    \caption{Demonstration of PPO and our GRPO. GRPO foregoes the value model, instead
 estimating the advantages from group scores.} \label{fig:GRPO}
\end{figure}

Group Relative Policy Optimization~(GRPO) \citep{deepseekmath} is the reinforcement learning algorithm that we adopt to train \dsro{} and \dsri{}. It was originally proposed to simplify the training process and reduce the resource consumption of Proximal Policy Optimization (PPO) \citep{schulman2017proximal}, which is widely used in the RL stage of LLMs \citep{ouyang2022training}. For an overall comparison between GRPO and PPO, see Figure~\ref{fig:GRPO}. 

For each question $q$, GRPO samples a group of outputs $\{o_1, o_2, \cdots, o_G\}$ from the old policy $\pi_{\theta_{old}}$ and then optimizes the policy model $\pi_{\theta}$ by maximizing the following objective:
\begin{equation}
\begin{split}
    \footnotesize
    & \mathcal{J}_{GRPO}(\theta) = \mathbb{E}{[q \sim P(Q), \{o_i\}_{i=1}^G \sim \pi_{\theta_{old}}(O|q)]}  \\
    & \frac{1}{G}\sum_{i=1}^G \left( \min \left( \frac{\pi_\theta(o_i |q)}{\pi_{\theta_{old}}(o_i |q)} A_i, \text{clip} \left( \frac{\pi_\theta(o_i |q)}{\pi_{\theta_{old}}(o_i |q)}, 1 - \epsilon, 1 + \epsilon \right)  A_i \right) - \beta \mathbb{D}_{KL}\left(\pi_{\theta} || \pi_{ref}\right)\right) ,
\end{split}
\label{eq:GRPO-obj}
\end{equation}
\begin{equation}
    \mathbb{D}_{KL}\left(\pi_{\theta} || \pi_{ref}\right) = \frac{\pi_{ref}(o_i|q)}{\pi_{\theta}(o_i|q)}- \log\frac{\pi_{ref}(o_i|q)}{\pi_{\theta}(o_i|q)} - 1,
\end{equation}
where $\pi_{ref}$ is a reference policy, $\epsilon$ and $\beta$ are hyper-parameters, and $A_i$ is the advantage, computed using a group of rewards $\{r_1, r_2, \ldots, r_G\}$ corresponding to the outputs within each group:
\begin{equation}
    A_i = \frac{r_i - {\mathrm mean(\{r_1, r_2, \cdots, r_G\})}}{{\mathrm std(\{r_1, r_2, \cdots, r_G\})}}.
\end{equation}
In contrast, in PPO, the advantage is typically computed by applying the Generalized Advantage Estimation (GAE) \citep{gae}, based not only on the rewards but also on a learned value model. Since the value model is usually of similar size as the policy model, it introduces a significant memory and computational overhead. Additionally, the training objective of the value model is to predict the expected cumulative reward from the current position onward, based on the tokens generated from the beginning up to the current position. This is inherently difficult, especially when only the final outcome reward is available. The challenge becomes even more pronounced when training long chain-of-thought reasoning models. As the output length increases, the model is more likely to engage in behaviors such as reflection and revision during generation, meaning that the content initially generated may later be revised or contradicted, which makes it even less feasible to predict the final reward based on a partial response. 

Another key difference between GRPO and PPO is how Kullback–Leibler (KL) divergence between the trained policy and the reference policy is incorporated into the training process. In GRPO, an unbiased estimator of the KL divergence \citep{kl_approx} is directly added in the loss as in equation \ref{eq:GRPO-obj}, while in PPO the per-token KL penalty is added as a dense reward at each token \citep{ouyang2022training}. Since the optimization goal of reinforcement learning is to maximize cumulative rewards, PPO's approach penalizes the cumulative KL divergence, which may implicitly penalize the length of the response and thereby prevent the model's response length from increasing. In addition, as we may train thousands of steps in the scenario of training long chain-of-thought reasoning models, the trained policy can diverge significantly from the initial reference policy. In order to balance the scope that the training policy can explore and the stability of the training, we periodically update the reference policy to the latest policy during the actual training process.

\begin{figure}
    \centering
    \includegraphics[width=\linewidth]{figures/ppo_vs_grpo.pdf}
    \caption{Performance of PPO and GRPO on the MATH task.} \label{fig:ppo_vs_grpo}
\end{figure}

Figure \ref{fig:ppo_vs_grpo} compares the performance of PPO and GRPO on the MATH task using DeepSeek-Coder-V2-Lite (16B MoE with 2.4B active parameters). Unlike GRPO, PPO requires additional hyperparameter tuning—particularly of the $\lambda$ coefficient in GAE—and is highly sensitive to this parameter. When $\lambda$ is set to 0.95 (the default value in most open-source PPO implementations), PPO performs considerably worse than GRPO. However, with careful tuning (setting $\lambda$ to 1.0), PPO's performance improves substantially, nearing that of GRPO. 
% Additionally, our experiments show that PPO's default method of training its value model using GAE-derived labels is less effective than using Monte Carlo labels, where only the final reward serves as the value target. In fact, when $\lambda$ is 0.95 and Monte Carlo labels are used, PPO ultimately achieves superior performance, slightly outperforming GRPO in later training phases. 

% Although PPO with well-selected hyperparameters can match or even surpass GRPO, this requires additional computational effort for hyperparameter tuning. Furthermore, in PPO, the value model typically matches the policy model in size, thereby doubling the overall training time compared to GRPO—a significant disadvantage when training large-scale models. Consequently, GRPO offers a more practical alternative when training resources are limited.

While PPO can achieve comparable performance when appropriately tuned, it demands additional computational cost for hyperparameter optimization. Moreover, considering the memory and computational overhead associated with training an additional value model, GRPO presents a more practical alternative, especially when training large-scale models with constrained resources.
\section{Training Details}

\subsection{RL Infrastructure} \label{sec:infra}

\begin{figure}
    \centering
    \includegraphics[width=\linewidth]{figures/rl_infra.pdf}
    \caption{Overview of our RL framework.} \label{fig:rl_infra}
\end{figure}

Conducting RL training on large models places high demands on the infrastructure. Our RL framework is architected with a decoupled and extensible structure to facilitate seamless integration of diverse models and algorithms. Within this framework, we have incorporated both intra-modular and inter-modular optimization techniques, to ensure training efficiency and scalability.

Specifically, as depicted in Figure~\ref{fig:rl_infra}, the framework is partitioned into four distinct modules, each corresponding to a specific phase of the RL pipeline:

\begin{itemize}
    \item \textbf{Rollout Module:} Prompts are loaded from training dataset and uniformly dispatched across multiple vLLM \citep{vllm} workers, each equipped with the actor model, to sample multiple responses. For \dsviii{} MoE architecture, we implement an expert parallelism strategy across nodes to reduce memory access overhead, and deploy redundant copies of hotspot experts to balance computational loads among different experts. Multi-Token Prediction (MTP) component is also leveraged for self-speculative decoding, significantly accelerating the decoding speed and effectively minimizing the completion time for the longest samples.
    \item \textbf{Inference Module:} This module loads the reward model and reference to perform a forward pass on the samples generated during the rollout phase, thereby obtaining model-based rewards and other essential information.
    \item \textbf{Rule-based Reward Module:} This module computes rule-based rewards for the model-generated responses. A unified interface has been designed to accommodate diverse implementations (e.g., code executor, answer matcher, format checker, etc.). Although this module does not require loading models into GPU memory, its execution tends to be time-consuming. To tackle this issue, an asynchronous scheduling approach is employed to overlap its execution with the Rollout and Inference modules, effectively hiding the associated latency.
    \item \textbf{Training Module:} This module loads the actor model and the critic model (if required), to compute loss and update model parameters. It provides flexible support for a variety of RL algorithms (e.g., PPO, GRPO, DPO, etc.). To minimize computational waste caused by sequence padding and balance the workload across devices, we design the following data packing strategy: first, all data in a global batch is sorted by length and distributed across processes within the data parallel group; subsequently, within each process, the Best-Fit strategy is applied to pack the data into fixed-length chunks with minimal padding; finally, the number of chunks is adjusted to be equal across all processes. Additionally, we have integrated the DualPipe algorithm, utilized in \dsviii{} training, to achieve efficient pipeline parallelism.
\end{itemize}

Notably, upon completion of each module (excluding the Rule-based Reward module), the model instances utilized in that phase are automatically offloaded from VRAM to either system memory or disk storage, thereby freeing up VRAM for the subsequent phase.

 \subsection{Reward Model Prompt}
\label{section:reward}
\begin{tcolorbox}[colback=gray!10,colframe=black!50,arc=4mm]
Please act as an impartial judge and evaluate the quality of the responses provided by two AI assistants to the user prompt displayed below. You will be given assistant A's answer and assistant B's answer. Your job is to evaluate which assistant's answer is better.\\Begin your evaluation by generating your own answer to the prompt. You must provide your answers before judging any answers.\\When evaluating the assistants' answers, compare both assistants' answers with your answer. You must identify and correct any mistakes or inaccurate information.\\Then consider if the assistant's answers are helpful, relevant, and concise. Helpful means the answer correctly responds to the prompt or follows the instructions. Note when user prompt has any ambiguity or more than one interpretation, it is more helpful and appropriate to ask for clarifications or more information from the user than providing an answer based on assumptions. Relevant means all parts of the response closely connect or are appropriate to what is being asked. Concise means the response is clear and not verbose or excessive.\\Then consider the creativity and novelty of the assistant's answers when needed. Finally, identify any missing important information in the assistants' answers that would be beneficial to include when responding to the user prompt.\\After providing your explanation, you must output only one of the following choices as your final verdict with a label:

1. Assistant A is significantly better: [[A$\gg$B]]\\
2. Assistant A is slightly better: [[A$>$B]]\\
3. Tie, relatively the same: [[A$=$B]]\\
4. Assistant B is slightly better: [[B$>$A]]\\
5. Assistant B is significantly better: [[B$\gg$A]]

Example output: \"My final verdict is tie: [[A$=$B]]\".
\end{tcolorbox}



% \paragraph{Safety Reward Model}
% To assess and improve model safety, we curated a dataset of 106,000 prompts with model-generated responses annotated as ``safe" or ``unsafe" according to predefined safety guidelines. This dataset was used to train a point-wise safety reward model to discriminate between safe and unsafe responses. Both safety and helpful reward models were trained with a batch size of 256 and a learning rate of 6e-6.


%\subsubsection{\dsri{} Evaluation}


%%=============================================%%
%% For submissions to Nature Portfolio Journals %%
%% please use the heading ``Extended Data''.   %%
%%=============================================%%

%%=============================================================%%
%% Sample for another appendix section			       %%
%%=============================================================%%

%% \section{Example of another appendix section}\label{secA2}%
%% Appendices may be used for helpful, supporting or essential material that would otherwise 
%% clutter, break up or be distracting to the text. Appendices can consist of sections, figures, 
%% tables and equations etc.

% Method of \dsri{}
% \subsection{Rejection Sampling} \label{sec:rej_sample}

\subsection{Data Recipe} \label{sec:rej_sample}
\subsubsection{RL Data} \label{sec: rl_data}
\begin{table}[t]
    \centering
    \setlength{\tabcolsep}{6pt}
    
    \caption{
    Description of RL Data  and Tasks.
    }
    \label{tab:rl_data}
    \begin{tabular}{l | c| c | c}
        \toprule
        \textbf{Data Type} & \textbf{\# Prompts} & \textbf{Question Type} & \textbf{Output Type} \\
        \midrule
        Math & 26K & Quantitative Reasoning & Number/Expression/Equation \\
        Code & 17K & Algorithm and Bug Fixing & Code Solution\\
        STEM & 22K & Multi-Choice & Option \\
        Logic & 15K & Choice/Quantitative Reasoning & Option/Number \\
        General & 66K & Helpfulness/Harmlessness & Ranked Responses \\
        \bottomrule
    \end{tabular}
\end{table}

Reasoning RL data includes four categories: mathematics, coding, STEM, and logic problems. In addition, we also incorporate general RL data to improve the helpfulness and harmlessness of the model in the training of \dsri{}. All questions are in Chinese or English. The description of the RL data can be found in Table \ref{tab:rl_data}, where we will describe the details of each data type one by one as follows:

\begin{itemize}
\item \textbf{Mathematics} dataset consists of 26k quantitative reasoning questions, including math exam questions and competition problems. The average number of prompt tokens is 122. The dataset covers various mathematical domains such as algebra, calculus, probability, and geometry. Problems range in difficulty from regional contests to international Olympiads.
For each problem, the model is expected to produce a step-by-step reasoning process culminating in a final answer, which can be a numerical value (e.g., ``$5$''), a mathematical expression (e.g., ``$x^2 + 3x - 2$''), or an equation (e.g., ``$y = 2x + 1$''). Mathematical proofs are excluded because it is difficult to determine their correctness.
For reinforcement learning purposes, we calculate the reward of a reasoning process by matching the predicted answer with the reference answer. If the answer aligns with the reference, the reward is assigned a value of 1; otherwise, it is assigned a value of 0.


\item \textbf{Coding} dataset includes 17k algorithm competition questions, along with 8k bug fixing problems. The algorithm competition questions are similar to problems found on platforms like Codeforces or LeetCode. Each problem typically includes a detailed problem description, constraints, and multiple input-output examples. The task is to write a complete function or program that can solve the problem correctly and efficiently, passing a comprehensive set of hidden test cases that assess both correctness and performance. These problems test algorithmic skills, including dynamic programming, graph theory, string manipulation, and data structure usage. 

The bug-fixing problems are extracted from real-world GitHub issues. Each task provides an issue description, a buggy version of the source code, and a set of unit tests that partially or completely fail. The goal is to understand the intent of the issue, locate and fix the defect in the code, and ensure that the corrected version passes all unit tests. 


\item \textbf{STEM} dataset comprises 22k choice questions that cover topics such as physics, chemistry, and biology. Each question in the STEM task presents a subject-specific problem accompanied by four to eight answer options. The model is required to select the most scientifically accurate answer based on the given context and domain knowledge. The average number of prompt tokens is 161. Specifically, the dataset includes 15.5\% physics, 30.7\% biology, 46.5\% chemistry, and 7.3\% other topics such as health and medicine. Since all STEM questions are multiple-choice, a binary reward is assigned based on whether the correct option is matched.


\item \textbf{Logic} dataset contains 15k questions designed to evaluate a model’s reasoning capabilities across a broad spectrum of logical challenges.  The dataset includes both real-world and synthetically generated problems. All problems support automatic evaluation, and the average prompt length is approximately 420 tokens.
The real-world portion of the dataset comprises a diverse selection of problems sourced from the web, including brain teasers, classical logic puzzles, and knowledge-intensive questions.  These questions are presented in a multiple-choice format to ensure objective and consistent assessment.
The synthetic portion consists primarily of two categories: code-IO problems and puzzle tasks. Code-IO problems are generated using the data pipeline introduced by \cite{li2025codei}, which converts competitive coding problems and their corresponding input-output test cases into verifiable logical reasoning problems. The puzzle tasks include problems intended to assess specific reasoning competencies. For example, cryptography puzzles are designed to evaluate a model's ability to identify and apply patterns in cipher schemes or perform string manipulations; logic puzzles focus on deductive reasoning over complex constraints, such as inferring valid conclusions from a fixed set of premises (e.g., the Zebra puzzle); and arithmetic puzzles test the model’s numerical reasoning (e.g. probability questions and 24 game).

% \item \textbf{Logic} dataset contains 15k questions. The average number of prompt tokens is 420. The question types include choice questions and quantitative reasoning questions. We directly apply the same verification methods mentioned above to validate the answers based on predefined rules.

\item \textbf{General} dataset consists of 66k questions designed to assess helpfulness, spanning various categories such as creative writing, editing, factual question answering, and role-playing. Additionally, the dataset includes 12,000 questions focused on evaluating harmlessness. To ensure robust verification, two reward models are utilized, each trained on a curated dataset of ranked responses generated by models in relation to helpfulness and harmlessness, respectively. We trained the helpful reward model for a single epoch with a maximum sequence length of 8192 tokens during the training phase. However, when deploying the model to generate reward signals, we did not impose any explicit length constraints on the input sequences being evaluated.
\end{itemize}


\subsubsection{\dsri{} Cold Start} \label{sec:cold-start-reason}

%\paragraph{\dsri{} Cold Start} \label{sec:cold-start-reason}
For DeepSeek-R1, we construct and collect a small amount of long CoT data to fine-tune the model as the initial RL actor. The motivation is primarily product-driven, with a strong emphasis on enhancing user experience. Users tend to find responses more intuitive and engaging when the reasoning process aligns with first-person perspective thought patterns. For example, DeepSeek-R1-Zero is more likely to employ the pronoun 'we' or avoid first-person pronouns altogether during problem solving, whereas DeepSeek-R1 tends to use 'I' more frequently. Furthermore, we acknowledge that such patterns may elicit unwarranted trust from users. Here, we would like to emphasize that the observed vivid reasoning patterns primarily reflect DeepSeek-engineered heuristics, rather than indicating that the model has inherently acquired human-like intelligence or autonomous problem-solving capabilities.


In cold start data creation, we prefer the thinking process that begins with comprehending the problem, followed by detailed reasoning that incorporates reflection and verification. The language employed throughout the thinking process is presented in the first-person perspective.   Additionally, maintaining language consistency is crucial for an optimal user experience. Without proper control, model responses may contain a mixture of different languages, regardless of the language used in the query. Such inconsistencies can disrupt comprehension and reduce user satisfaction. Therefore, careful refinement is essential to ensure that responses remain coherent and aligned with user intent. Nevertheless, we acknowledge that the raw Chain-of-Thought (CoT) reasoning produced by DeepSeek-R1-Zero may possess potential that extends beyond the limitations of current human priors.
Specifically, we first engage human annotators to convert the reasoning trace into a more natural, human conversational style. The modified data pairs are then used as examples to prompt an LLM to rewrite additional data in a similar style. All LLM-generated outputs subsequently undergo a second round of human verification to ensure quality and consistency.

\lstinputlisting[style=data,caption={Prompt for producing a human-readable solution.},label={lst:prompt:summary}]{data/prompt_summary.md}

Specifically, we begin by gathering thousands of high-quality, diverse reasoning prompts. For each prompt, we generate multiple reasoning trajectories using \dsro{} with a relatively high temperature of $1.0$. Next, we filter these generations to retain only those with correct final answers and a readable format. For mathematical outputs, we use \texttt{sympy} ({\url{https://www.sympy.org/}}) for parsing and expression comparison; and for formatting, we apply rules such as repetition detection and language-mixing filtering. Finally, we prompt DeepSeek-V3 to refine both the reasoning and the summaries to ensure proper formatting and a human-friendly expression.
In particular, to resolve language mixing, we instruct DeepSeek-V3 to ``Translate the thinking process to the same language as the question.'' Since DeepSeek-R1-Zero's summary only provided the final answer, we use the summary prompt in Listing \ref{lst:prompt:summary} to produce a concise, human-readable solution that outlines both the reasoning steps and the final result.

For code data, we collect a large set of competitive programming problems. In detail, We have compiled an extensive collection of competitive programming problems from multiple online judge (OJ) platforms, specifically 5151 problems from Codeforces and 2504 problems from AtCoder. Since the original test cases are not publicly available from these platforms, we developed a methodology to create reliable test cases for each problem.

Our approach involves using DeepSeek-V2.5 to generate candidate test cases, followed by a rigorous validation process. Specifically, we prompted DeepSeek-V2.5 to write Python programs that generate test cases tailored to each problem's requirements as shown in Listing~\ref{lst:prompt:code}.


After obtaining numerous candidate test cases, we implemented a two-phase filtering procedure. First, we used correct submissions to eliminate invalid test cases that produced incorrect outputs. Then, we strategically selected subsets of test cases that successfully identified flaws in incorrect submissions. This process ensured our final test cases properly differentiated between correct and incorrect solutions for each problem. 

% For each problem, we prompt DeepSeek-V3 to generate code that produces corresponding test cases. These test cases are then refined by filtering out incorrect ones using actual submissions. Next, we generate multiple solutions for each problem using \dsro{} and validate their correctness against the refined test cases. Finally, we prompt DeepSeek-V3 to enhance the reasoning and summaries of the solutions, improving their clarity and overall quality.



\lstinputlisting[style=data,caption={Prompt for generating test cases for code problem.},label={lst:prompt:code}]{data/prompt_code.md}

In addition, we employ few-shot prompting for DeepSeek-V3 to generate responses to simple math problems, such as "1 + 1 = ?", ensuring that the outputs remain concise and appropriately structured.  We provide the prompt for a simple math problem in Listing \ref{lst:prompt:easy-math}.

\lstinputlisting[style=data,caption={Prompt for generating chain-of-thought for simple math problem.},label={lst:prompt:easy-math}]{data/prompt_easy_math.md}


\lstinputlisting[style=data,caption={An example prompt of using DeepSeek-V3 as a judge.},label={lst:prompt:lm-judge}]{data/prompt_lm_judge.md}

\subsubsection{800K Supervised Data} \label{sec:sft_data}
\paragraph{Reasoning Data} 
We curate a large set of reasoning prompts and generate reasoning trajectories by performing rejection sampling from the checkpoint of the first-stage RL training.
In the previous stage, we only included data that could be evaluated using rule-based rewards. However, in this stage, we expand the dataset by incorporating additional data, some of which uses a generative reward model by feeding the ground-truth and model predictions into DeepSeek-V3 for judgment, an example prompt is provided in Listing \ref{lst:prompt:lm-judge}.
Additionally, because the model output is sometimes chaotic and difficult to read, we have filtered out chain-of-thought with mixed languages, long paragraphs, and code blocks.
For each prompt, we sample multiple responses and retain only the correct ones. In total, we collect about 600k reasoning-related training samples. 


% \subsubsection{Rejection Sampling}
\paragraph{Non-Reasoning Data}
For non-reasoning data, such as writing, factual QA, self-cognition, and translation, we adopt the DeepSeek-V3 pipeline and reuse portions of the SFT dataset of DeepSeek-V3. We also incorporate software engineering-focused data, including program repair and front-end web development, to enhance the model's ability to solve real-world problems. For certain non-reasoning tasks, we call DeepSeek-V3 to generate a potential chain-of-thought before answering the question by prompting. However, for simpler queries, such as ``hello'' we do not provide a CoT in response.
In the end, we collected a total of approximately 200k training samples that are unrelated to reasoning.

When designing our thinking process style, we ask the model to follow key principles: First, keep each paragraph concise and digestible. Short paragraphs make ideas clearer and easier to follow. Second, adopt a conversational tone that feels natural and engaging. We avoid technical formatting like markdown to maintain a smooth reading experience. Third, and most importantly, the thinking process begins by understanding the complete user context. This means analyzing who our users are, what situations they're dealing with, and what they truly need - including those unstated needs that may lie beneath the surface of their initial request.

After eliciting these thinking processes from the model, human annotators meticulously verify the accuracy of the outputs. Our findings indicate that these artificial reasoning traces enhance the model's precision in interpreting user queries. Specifically, they effectively highlight format constraints, clarify user intentions, and elucidate the requisite structure of outputs. This methodological approach facilitates more accurate and responsive interactions between the model and users.


\begin{table}[ht]
\centering
\caption{Data Statistics of SFT Data.}
\label{tab:800k}
\begin{tabular}{lrrr}
\toprule
{Domain} & {Num Samples} & {Avg Rounds} & {Avg Tokens} \\
\midrule
Math & 395285 & 1.0 & 6094.2 \\
Code & 211129 & 1.1 & 7435.7 \\
STEM & 10124 & 1.0 & 4928.8 \\
Logic & 10395 & 1.0 & 2739.0 \\
General & 177812 & 1.1 & 1419.8 \\
\midrule
Total & 804745 & 1.0 & 5355.3 \\
\bottomrule
\end{tabular}
\end{table}
\paragraph{SFT Data Statistics}
Table \ref{tab:800k} summarizes the data statistics across various domains, based on approximately 800,000 supervised samples. It is worth noting that the majority of the data consists of single-turn interactions, which may limit the multi-turn conversational capabilities of \dsri{}. We leave the expansion to multi-turn dialogue data as future work. The math-related data are primarily in Chinese and English, spanning a wide range of topics and difficulty levels. These questions are verifiable, either through deterministic rules or by reference to specific ground-truth answers. The code dataset encompasses not only competitive programming problems but also debugging tasks and project-oriented coding queries. STEM and logic-related questions, although smaller in volume compared to mathematics and code, are sourced from publicly available textbooks and online repositories. The general-domain queries are diverse, predominantly in Chinese and English, and include tasks such as creative writing, open-ended question answering, role-playing scenarios, and more.


% 每个 domain 给一个 sft 数据的例子
\paragraph{Examples of SFT Trajectories}
\label{sec:sft-example}

We provide examples of SFT trajectories related to math, code, and non-reasoning in Listings \ref{lst:sft:math}, \ref{lst:sft:code}, and \ref{lst:sft:non-reasoning}.

\lstinputlisting[style=data,caption={An example SFT trajectory from reasoning data related to mathematics.},label={lst:sft:math}]{data/sft_sample_math.md}


\lstinputlisting[style=data,caption={An example SFT trajectory from reasoning data related to code.},label={lst:sft:code}]{data/sft_sample_code.md}

% \lstinputlisting[style=data,caption={An example SFT trajectory from reasoning data related to logic.}]{data/sft_sample_logic.md}

\lstinputlisting[style=data,caption={An example SFT trajectory from non-reasoning data related to writing.},label={lst:sft:non-reasoning}]{data/sft_sample_writing.md}


\subsection{Hyper-Parameters}

% The clipping ratio $\epsilon$ is crucial in training. A lower value may truncate gradients for many tokens, degrading model performance and hindering its ability to develop longer reasoning times for problem-solving. 
% Conversely, a higher value can introduce instability during training.

\subsubsection{Hyper-Parameters of \dsro{}-Qwen-32B } \label{sec: qwen-zero}
To train \dsro{}-Qwen-32B, we set the learning rate to 2e-6, the KL coefficient to 0.001, and the sampling temperature to 1 for rollout. For each question, we sample 16 outputs with a maximum length of 32,768. Each training step consists of 32 unique questions, resulting in a training batch size of 512 per step. Every 400 steps, we replace the reference model with the latest policy model. To accelerate training, each rollout generates 8,192 outputs, which are randomly split into 16 mini-batches and trained for only a single inner epoch.

\subsubsection{Hyper-Parameters of SFT} \label{sec: r1-1-sft}

For code-start SFT and the second-stage SFT,
we fine-tune DeepSeek-V3-Base for 2-3 epochs using the curated dataset, as described in \ref{sec:rej_sample}.
We employ a cosine decay learning rate scheduler, starting at \(5 \times 10^{-5}\) and gradually decreasing to \(5 \times 10^{-6}\). The maximum context length is set to 32,768 tokens, and the batch size is 128.



\subsubsection{Hyper-Parameters of Distillation}
\label{appendix:distill-params}

For distillation, we fine-tune the corresponding base model for 2–3 epochs using the 800k data described in Section \ref{sec:sft_data}. The base model and initial learning rate are listed in Table \ref{tab:distill_config}. We employ a cosine decay learning rate scheduler that gradually decreases the learning rate to one-tenth of its initial value. The maximum context length is 32,768 tokens, and the batch size is 64.


\begin{table}[h]
    \centering
      \caption{DeepSeek-R1 Distilled Models, their corresponding Base Models, and Initial Learning Rates.}
    \label{tab:distill_config}
    \begin{tabular}{llc}
        \toprule
        \textbf{Distilled Model} & \textbf{Base Model} & \textbf{Initial Learning Rate} \\
        \midrule
        \href{https://huggingface.co/deepseek-ai/DeepSeek-R1-Distill-Qwen-1.5B}{DeepSeek-R1-Distill-Qwen-1.5B} &
        \href{https://huggingface.co/Qwen/Qwen2.5-Math-1.5B}{Qwen2.5-Math-1.5B} & $1\times10^{-4}$ \\
        \href{https://huggingface.co/deepseek-ai/DeepSeek-R1-Distill-Qwen-7B}{DeepSeek-R1-Distill-Qwen-7B} &
        \href{https://huggingface.co/Qwen/Qwen2.5-Math-7B}{Qwen2.5-Math-7B} & $8\times10^{-5}$ \\
        \href{https://huggingface.co/deepseek-ai/DeepSeek-R1-Distill-Qwen-14B}{DeepSeek-R1-Distill-Qwen-14B} &
        \href{https://huggingface.co/Qwen/Qwen2.5-14B}{Qwen2.5-14B} & $7\times10^{-5}$ \\
        \href{https://huggingface.co/deepseek-ai/DeepSeek-R1-Distill-Qwen-32B}{DeepSeek-R1-Distill-Qwen-32B} &
        \href{https://huggingface.co/Qwen/Qwen2.5-32B}{Qwen2.5-32B} & $6\times10^{-5}$ \\
        \href{https://huggingface.co/deepseek-ai/DeepSeek-R1-Distill-Llama-8B}{DeepSeek-R1-Distill-Llama-8B} &
        \href{https://huggingface.co/meta-llama/Llama-3.1-8B}{Llama-3.1-8B} & $5\times10^{-5}$ \\
        \href{https://huggingface.co/deepseek-ai/DeepSeek-R1-Distill-Llama-70B}{DeepSeek-R1-Distill-Llama-70B} &
        \href{https://huggingface.co/meta-llama/Llama-3.3-70B-Instruct}{Llama-3.3-70B-Instruct} & $2\times10^{-5}$ \\
        \bottomrule
    \end{tabular}
  
\end{table}

\subsubsection{Training Cost} \label{sec:cost}
Regarding our research on \dsri{}, we utilized the A100 GPUs to prepare for the experiments with a smaller model (30B parameters). The results from this smaller model have been promising, which has allowed us to confidently scale up to 660B R1-Zero and R1.
For the training of \dsro{}, we employed 64*8 H800 GPUs, and the process required approximately 198 hours. Additionally, during the training phase of \dsri{}, we utilized the same 64*8 H800 GPUs, completing the process in about 4 days, or roughly 80 hours. To create the SFT datasets, we use 5K GPU hours. The details are shown in Table \ref{tab:cost}.


\begin{table}[t]
    \centering
    \setlength{\tabcolsep}{6pt}
        \caption{
    Training costs of \dsri{}, assuming the rental price of H800 is \$2 per GPU hour.
    }
    \label{tab:cost}
    \begin{tabular}{l | c c c | c}
        \toprule
        \textbf{Training Costs} & \textbf{\dsro{}} & \textbf{SFT data creation} & \textbf{\dsri{}} & \textbf{Total} \\
        \midrule
        in H800 GPU Hours & 101K & 5K & 41K & 147K \\
        in USD & \$202K & \$10K & \$82K & \$294K \\
        \bottomrule
    \end{tabular}

\end{table}

\subsection{Reward Hacking} \label{sec:reward_hack}
In the context of LLM training, reward hacking refers to the phenomenon wherein a model exploits flaws or biases in the reward function, thereby achieving high reward scores without truly aligning with the underlying human intent. In our work, we observe such reward hacking behavior when employing the helpful reward model. Specifically, if the reward model contains systematic biases or inaccuracies, the LLM may learn to generate responses that are rated highly by the model but diverge from authentic human preferences. This misalignment can manifest in performance degradation on tasks requiring complex reasoning, as illustrated in Figure \ref{fig:reward_hacking}.

%\subsection{Evolution of Reasoning Capability in \dsro{} during Training}
\begin{figure}
    \centering
    \includegraphics[width=\linewidth]{figures/reward_hacking.pdf}
    \caption{Reward hacking: the reward exhibits an increasing trend as the performance on CodeForces decreases for training. } \label{fig:reward_hacking}
\end{figure}


\subsection{Ablation Study of Language Consistency Reward}
\label{sec:LC_reward} 
To study the impact of the Language Consistency (LC) Reward, we conduct an ablation experiment on DeepSeek-R1-Distill-Qwen-7B. This model uses the same cold start data as \dsri{} and also exhibits language mixing during the RL process. The results are shown in Figure \ref{fig:lc_reward}. As can be seen, without the LC reward, language consistency gradually deteriorates as training steps increase. However, when the LC reward is applied, stable language consistency is maintained throughout the training process. For benchmark performance, the model maintains comparable performance on the mathematical benchmark, while a slight degradation is observed on the coding benchmark. Although such alignment results in a slight degradation in model performance, this reward aligns with human preferences, making the output more readable.
\begin{figure}[h]
    \centering
    \includegraphics[width=\linewidth]{figures/LC_Reward.pdf}
    \caption{The experiment results of Language Consistency (LC) Reward during reinforcement learning.}
    \label{fig:lc_reward}
\end{figure}

\section{Self-Evolution of \dsro{}}
\subsection{Evolution of Reasoning Capability in \dsro{} during Training}
\begin{figure}
    \centering
    \includegraphics[width=0.85\linewidth]{figures/math_performance.pdf}
    \caption{Performance of \dsro{} on problems with varying difficulty levels in the MATH dataset.} \label{fig:math_performance_difficulty}
\end{figure}
We analyzed DeepSeek-R1-Zero's performance on the MATH dataset stratified by difficulty levels (1-5). Figure \ref{fig:math_performance_difficulty} reveals distinct learning patterns: easy problems (levels 1-3) quickly reach high accuracy (0.90-0.95) and remain stable throughout training, while difficult problems show remarkable improvement - level 4 problems improve from near 0.78 to 0.95, and the most challenging level 5 problems demonstrate the most dramatic improvement from near 0.55 to 0.90.

One may find it counterintuitive that the model's accuracy on harder questions (levels 3-4) occasionally surpasses its performance on easier questions (level 1) by a small margin. This apparent anomaly stems from several dataset characteristics. The MATH dataset is unevenly distributed, with level-1 questions comprising only 43 of 500 examples, while higher levels contain approximately 100 questions each. Consequently, the model's 95-97\% accuracy on level-1 represents just 1-2 unsolved problems, primarily in geometry, where the model still struggles. Furthermore, the distribution of mathematical categories (geometry, algebra, etc.) varies across difficulty levels due to the dataset's construction methodology. It's also worth noting that these difficulty levels were annotated based on human perception of problem complexity rather than machine learning considerations.

Despite these nuances in comparing raw accuracy percentages across difficulty levels, the training trends still demonstrate that while simpler reasoning tasks (for humans) are mastered early in training, the model's capability on complex reasoning problems (level 3-5) significantly improves over time.

% This fine-grained analysis demonstrates that while simpler reasoning tasks are mastered early in training, the model's capability on complex reasoning problems significantly improves over time. The convergence of performance across difficulty levels during later training steps (except for the most difficult level 5, which maintains a slight gap) suggests the model develops advanced reasoning abilities that generalize across problem complexities. 

\subsection{Evolution of Advanced Reasoning Behaviors in \dsro{} during Training}
\label{app:evo_zero}
\begin{figure}[t]
\centering
\subfigure{
\includegraphics[width=0.47\textwidth]{figures/wait_count_1.pdf}
}
\subfigure{
\includegraphics[width=0.47\textwidth]{figures/wait_count_2.pdf}
}
\caption{Evolution of reasoning behaviors during training.  (a) Frequency of representative reflective words during the training process;
(b) Specific occurrence patterns of the word ``wait'' throughout the training process.}
\label{fig:r1-zero-evol}
\end{figure}
We analyze the change in the reasoning behavior of the model during training. 

First, as shown in Figure \ref{fig:r1-zero-evol}(a),  we counted some representative reflective words, including ``wait'', ``mistake'', ``however'', ``but'', ``retry'',
``error'', ``verify'',
``wrong'', ``evaluate'', and ``check'' . These reflective words were selected by $3$ human experts, who are asked to think of several reflective words and then merge them into a final word list.
As is shown, there is a gradual increase in the frequency of reflective behaviors as training progresses. Specifically, the count of the reflective words rises 5- to 7-fold compared to the start of training, suggesting that RL plays a key role in generating long-chain intermediate tokens.

Second, specific reflective behaviors may appear at particular points in training. The analysis of the word ``wait'' (Figure \ref{fig:r1-zero-evol}(b)) demonstrates this clearly. This reflective strategy was nearly absent during early training, showed occasional usage between steps 4000-7000, and then exhibited significant spikes after step 8000. This suggests that the model learns different forms of reflection at specific stages of development.

In conclusion, we observe a gradual increase in the model's reflective behavior during training, while certain reflection patterns like the use of ``wait'' emerge at specific points in the training process.






\section{Evaluation of \dsri{}}\label{sec:eval}
\subsection{Experiment Setup}

\paragraph{Benchmarks} We evaluate models on MMLU \citep{mmlu}, MMLU-Redux \citep{mmlu_redux}, MMLU-Pro \citep{mmlu_pro}, C-Eval \citep{ceval},  IFEval~\citep{IFeval}, FRAMES~\citep{frames}, GPQA Diamond ~\citep{gpqa}, SimpleQA~\citep{simpleqa}, C-SimpleQA~\citep{csimpleqa}, SWE-Bench Verified~\citep{swe_verified}, Aider \citep{aider}, LiveCodeBench~\citep{livecodebench} (2024-08 -- 2025-01), Codeforces \citep{codeforces}, Chinese National High School Mathematics Olympiad (CNMO 2024) \citep{cnmo}, and American Invitational Mathematics Examination 2024 (AIME 2024)~\citep{AIME2024}. 
% 


Specifically, MMLU, MMLU-Redux, MMLU-Pro, C-Eval, and CMMLU are multiple-choice benchmarks designed to assess model performance on general encyclopedic knowledge. Higher scores on these benchmarks indicate a broader understanding of world knowledge and the ability to correctly answer questions in a multiple-choice format. SimpleQA and C-SimpleQA evaluate model performance on long-tail knowledge, while GPQA assesses the ability to solve Ph.D.-level tasks in physics, chemistry, and biology. IFEval is designed to evaluate the model’s capacity to generate outputs in a required format. FRAMES and DROP focus on assessing model performance in processing and reasoning over long documents. In addition to these standard benchmarks, we also evaluate our models on open-ended generation tasks, employing LLM as judges. We follow the original evaluation protocols of AlpacaEval 2.0 and Arena-Hard, utilizing GPT-4-Turbo-1106 for pairwise comparisons. To mitigate length bias, only the final summary is provided to the evaluation model.

LiveCodeBench and Codeforces are designed to measure model performance on algorithmic competition tasks, whereas SWE-Verified and Aider assess the model’s capabilities on real-world software engineering problems. Finally, AIME, MATH-500, and CNMO 2024 comprise mathematics problems that test the model’s reasoning abilities in mathematical domains.

% distill benchmarks
For distilled models, we report representative results on AIME 2024, MATH-500, GPQA Diamond, Codeforces, and LiveCodeBench.


\paragraph{Decontamination}\label{sec:conta}
To prevent benchmark contamination, we implemented comprehensive decontamination procedures for both pre-training and post-training data. DeepSeek-V3 base has a knowledge cutoff date of July 2024, predating evaluation benchmarks like CNMO 2024, and we filtered out any text segments (including web pages and GitHub files) that contained matching 10-gram sequences from evaluation questions or reference solutions. As one example of our decontamination efforts, in the mathematics domain alone, our decontamination process identified and removed approximately six million potential pre-training texts. For post-training, mathematical SFT data and RL training prompts were sourced exclusively from pre-2023 competitions and underwent the same n-gram filtering protocol used in pre-training, ensuring no overlap between training and evaluation data. These measures ensure our model evaluation results reflect genuine problem-solving capabilities rather than memorization of test data.

However, we acknowledge that  the n-gram based decontamination method cannot prevent the paraphrase of testset. Therefore, it is possible that benchmarks released before 2024 may suffer from contamination issues. 


\paragraph{Evaluation Prompts} Following the setup in DeepSeek-V3, standard benchmarks such as MMLU, DROP, GPQA Diamond, and SimpleQA are evaluated using prompts from the simple-evals framework. For MMLU-Redux, we adopt the Zero-Eval prompt format~\citep{Lin_ZeroEval_A_Unified_2024} in a zero-shot setting. In terms of MMLU-Pro, C-Eval and CLUE-WSC, since the original prompts are few-shot, we slightly modify the prompt to the zero-shot setting. The CoT in few-shot may hurt the performance of  \dsri{}. 
Other datasets follow their original evaluation protocols with default prompts provided by their creators. 
For code and math benchmarks, the HumanEval-Mul dataset covers eight mainstream programming languages (Python, Java, C++, C\#, JavaScript, TypeScript, PHP, and Bash). 
Model performance on LiveCodeBench is evaluated using  CoT format, with data collected between August 2024 and January 2025. The Codeforces dataset is evaluated using problems from 10 Div.2 contests, along with expert-crafted test cases, after which the expected ratings and percentages of competitors are calculated. SWE-Bench verified results are obtained via the agentless framework~\citep{agentless}. AIDER-related benchmarks are measured using a "diff" format.
% For mathematical evaluations, AIME and CNMO 2024 are tested at a temperature of 0.7, with results averaged across 16 runs, while MATH-500 uses greedy decoding. 
\dsri{} outputs are capped at a maximum of 32,768 tokens for each benchmark.

Table \ref{tab:mmlu_eval_format_example} to Table \ref{tab:math_eval_format_example} present examples of our evaluation formats on different benchmarks.  We also detail the specific capabilities of large language models assessed by each benchmark in the corresponding table captions.



\paragraph{Baselines} We conduct comprehensive evaluations against several strong baselines, including DeepSeek-V3, Claude-Sonnet-3.5-1022, GPT-4o-0513, OpenAI-o1-mini, and OpenAI-o1-1217. Since accessing the OpenAI-o1-1217 API is challenging in mainland China, we report its performance based on official reports.
% distill
For distilled models, we also compare the open-source model QwQ-32B-Preview \citep{QwQ}.

% \paragraph{Evaluation Setup} Some previous approaches employ greedy decoding to calculate pass@1. However, we observe that this method introduces a performance fluctuation of approximately 5\%, which fails to accurately represent the true performance of model.
% To address this, we adopt a different evaluation strategy. 

% where k depends on the size of the evaluation set

\label{sec: evaluation_set}
We set the maximum generation length to 32,768 tokens for the models.
We found that using greedy decoding to evaluate long-output reasoning models results in higher repetition rates and significant variability across different checkpoints.
Therefore, we default to pass@$k$ evaluation \citep{codex} and report pass@1 using a non-zero temperature.
Specifically, we use a sampling temperature of $0.6$ and a top-$p$ value of $0.95$ to generate $k$ responses (typically between $4$ and $64$, depending on the test set size) for each question.
Sepcifically, we use $k=64$ for AIME and GPQA, $k=16$ for MATH and CodeForces, and $k=8$ for LCB.
Pass@1 is then calculated as
\[
\text{pass@1} = \frac{1}{k} \sum_{i=1}^{k} p_i,
\]
where $p_i$ denotes the correctness of the $i$-th response. This method provides more reliable performance estimates.
% 
For AIME 2024, we also report consensus (majority vote) results using $64$ samples, denoted as $\text{cons}@64$.

\subsection{Main Results}
\input{tables/chateval}
\paragraph{Standard Benchmark} \label{sec: benchmark}
We evaluate \dsri{} on multiple benchmarks. 
For education-oriented knowledge benchmarks such as MMLU, MMLU-Pro, and GPQA Diamond, \dsri{} demonstrates superior performance compared to DeepSeek-V3. This improvement is primarily attributed to enhanced accuracy in STEM-related questions, where significant gains are achieved through large-scale reinforcement learning. Additionally, \dsri{} excels on FRAMES, a long-context-dependent QA task, showcasing its strong document analysis capabilities. This highlights the potential of reasoning models in AI-driven search and data analysis tasks. 

\dsri{} also delivers impressive results on IF-Eval, a benchmark designed to assess a model's ability to follow format instructions. These improvements can be linked to the inclusion of instruction-following data during the final stages of SFT and RL training. Furthermore, remarkable performance is observed on AlpacaEval2.0 and ArenaHard, indicating \dsri{}’s strengths in writing tasks and open-domain question answering.

On math tasks, \dsri{} demonstrates performance on par with OpenAI-o1-1217, surpassing other models by a large margin. A similar trend is observed on coding algorithm tasks, such as LiveCodeBench and Codeforces, where reasoning-focused models dominate these benchmarks. On engineering-oriented coding tasks, OpenAI-o1-1217 outperforms \dsri{} on Aider but achieves comparable performance on SWE Verified. We believe the engineering performance of \dsri{} will improve in the next version, as the amount of related RL training data currently remains very limited. 



\begin{figure}[th]

\includegraphics[width=1.0\textwidth]{figures/dsr1_nature.pdf}
\caption{
The benchmark performance of \dsri{} and \dsro{} is compared with human scores across different datasets. For AIME and Codeforces, the human scores represent the average performance of all human competitors. In the case of GPQA, the human score corresponds to Ph.D.-level individuals who had access to the web for answering the questions. 
}
\label{fig:dsv3_performance}
\end{figure}
Figure \ref{fig:dsv3_performance} presents a comparative analysis of the performance of \dsro{}, \dsri{}, and human participants across several benchmark competitions. Notably, the AIME is a mathematics competition designed for high school students, and \dsri{} demonstrates performance that surpasses the mean score achieved by human competitors in this event. On the Codeforces platform, \dsri{} outperforms 96.3\% of human participants, underscoring its advanced problem-solving capabilities. In the case of GPQA, where human experts—typically individuals with Ph.D.-level qualifications and access to web resources—participate, human performance remains superior to that of \dsri{}. However, we anticipate that enabling web access for \dsri{} could substantially enhance its performance on GPQA, potentially narrowing or closing the observed gap.


\paragraph{Human Evaluation}  We utilize ChatbotArena  \citep{chiang2024chatbot} to show the human preference of \dsri{} with its ranking and elo score.  ChatbotArena is an open, crowdsourced platform developed by LMSYS and UC Berkeley SkyLab to evaluate and rank LLMs based on human preferences. Its core mechanism involves pairwise comparisons, where two anonymous LLMs (randomly selected from a pool of over 100 models) respond to a user-submitted prompt. Users then vote on which response they prefer, declare a tie, or mark both as bad, without knowing the models' identities until after voting. This double-blind approach ensures fairness and reduces bias.
The platform collects millions of user votes as of recent updates—and uses them to rank models with the Elo rating system, a method adapted from chess that predicts win rates based on pairwise outcomes. To improve stability and incorporate new models efficiently, Chatbot Arena employs a bootstrap-like technique, shuffling vote data across permutations to compute reliable Elo scores. It has also begun adopting the Bradley-Terry model, which refines rankings by estimating win probabilities across all battles, leveraging the full vote history.

DeepSeek-R1 has demonstrated remarkable performance in ChatbotArena. Figure \ref{lmsysrank} presents the overall ranking of \dsri{} on ChatbotArena as of January 24, 2025, where \dsri{} shares the first position alongside OpenAI-o1 and Gemini-Exp-1206 on the style control setting. Style control refers to a feature introduced to separate the influence of a model's response style (e.g., length, formatting, tone) from its substantive content (e.g., accuracy, relevance, reasoning) when evaluating and ranking LLMs. This addresses the question of whether models can "game" human preferences by producing responses that are longer, more polished, or better formatted, even if their content isn’t necessarily superior. It is a huge milestone that an open-source model under the MIT License could achieve comparable performance with closed-source models, especially considering that the cost of \dsri{} is relatively inexpensive. Figure \ref{lmsysdetail} illustrates the rankings across different evaluation dimensions, highlighting \dsri{}'s strong performance in mathematics, coding, and other areas. This demonstrates that \dsri{} excels not only in reasoning but also across a wide range of domains.

\begin{figure}
    \centering
    \includegraphics[width=\linewidth]{figures/lmsysrank.pdf}
    \caption{The style control ranking on ChatBotArena of \dsri{}. The screenshot is captured on January 24, 2025, one week after model release. The ranking is dynamically updated in real time as the number of votes increases. }  \label{lmsysrank}
\end{figure}


\begin{figure}
    \centering
    \includegraphics[width=\linewidth]{figures/lmsysdetail.pdf}
    \caption{The rank of \dsri{} across various aspects  on January 24, 2025.} \label{lmsysdetail}
\end{figure}

 \subsection{DeepSeek-R1 Safety Report }
\label{sec:safety_eval}

{\colorbox{yellow} {Warning: This section contains potentially risky and offensive content!}}

We fully recognize that, while open source sharing facilitates the dissemination of advanced technologies within the community, it also introduces potential risks of misuse. In this section, we systematically present the security risk assessment of DeepSeek-R1. Specifically, we structure our analysis around the following aspects: 
(1) \ref{sec:risk_control}: the risk control system for the official DeepSeek-R1 service,
(2) \ref{sec:safety_eval_standard}: a comparative safety evaluation with other state-of-the-art models on 6 publicly safety benchmarks, 
(3) \ref{subsubsection:safety_taxnomic}: a taxonomy-based study based on in-house safety testsets,
(4) \ref{sec:safety_multilingual}: an evaluation of the multilingual safety of the R1 model,
(5) \ref{sec:safety_jailbreak}: an assessment of the model's robustness to jailbreak attacks.

\subsubsection{Risk Control System for DeepSeek-R1}
\label{sec:risk_control}
Generally, beyond the intrinsic safety of models, model-based services typically implement an external risk control system to enhance system-level security. In this subsection, we introduce the risk control system deployed in the official DeepSeek services. In the comparative experiments presented later in this chapter, we will report the results of DeepSeek-R1 with and without risk control measures. For models from other manufacturers, the results represent the comprehensive safety performance that integrates both the model's intrinsic safety mechanisms and external risk control systems.

\lstinputlisting[style=data,caption={The Risk Review Prompt for DeepSeek-R1.},label={lst:safe:risk_check_prompt}]{data/risk_check_prompt.md}

The risk control system for DeepSeek-R1 is implemented by sending risk review prompts to DeepSeek-V3. Specifically, it includes the following two main processes:

\paragraph{Potential Risky Dialogue Filtering}
After each round of conversation, the user's query is automatically matched against a predefined keyword list. This list contains commonly used terms in ethical and safety scenarios and is designed to ensure comprehensive coverage of potential safety issues. Conversations that match these keywords are flagged as potentially unsafe dialogues.

\paragraph{Model-based Risk Review}
Subsequently, these potentially unsafe dialogues are concatenated with a preset risk review prompt (shown in Listing \ref{lst:safe:risk_check_prompt}) and sent to the DeepSeek-V3 model (considering the balance between effectiveness and efficiency). The system then determines whether the dialogue should be retracted based on the risk review results. We have meticulously designed this risk review prompt to effectively cover various safety scenarios and maintain good scalability.

The subsequent experimental results show that with the addition of a risk control system, the overall safety of services significantly improves, particularly against dangerous tactics such as jailbreak attacks. Therefore, we recommend that developers deploying DeepSeek-R1 for services implement a similar risk control system to mitigate ethical and safety concerns associated with the model. Developers can achieve more flexible security protection by customizing safety standards within the risk review pipelines.



\subsubsection{R1 Safety Evaluation on Standard Benchmarks}
\label{sec:safety_eval_standard}
In this section, we present the performance of the DeepSeek-R1 model on comprehensive open source safety benchmarks.
We first introduce the composition of these evaluation datasets. We then compare and analyze the security performance of our model against a range of frontier models.

Given the broad scope of security-related topics, we selected six publicly available benchmark datasets, each focusing on different aspects of security, to ensure a comprehensive and well-rounded evaluation. The following is an introduction to these evaluation benchmarks.
\begin{itemize}
    \item \textbf{Simple Safety Tests}~\citep{safety-benchmark-sst}: Short for SST, this benchmark primarily covers security evaluations in the following five categories: Illegal Items, Physical Harm, Scams \& Fraud, Child Abuse, and Suicide, Self-Harm \& Eating Disorders (SH \& ED).
    \item \textbf{Bias Benchmark for QA}~\citep{safety-benchmark-bbq}: Short for BBQ, this benchmark primarily evaluates the performance of language models in conversations involving discriminatory biases. Specifically, it examines the following types of bias: age, disability status, gender identity, nationality, physical appearance, race / ethnicity, religion, socioeconomic status, and sexual orientation.
    \item \textbf{Anthropic Red Team}~\citep{safety-benchmark-anthropic-red-team}：Short for ART, this benchmark consists of data collected by Anthropic during Red Team attacks on the model. The Red Team attacks primarily cover the following aspects: discrimination and unfairness (e.g., racial and gender bias); hate speech and offensive language (e.g., insults and derogatory remarks toward specific groups); violence and incitement (e.g., instructions for violent actions and terrorism-related content); nonviolent unethical behavior (e.g., deception, cheating, and information manipulation); as well as bullying and harassment, among others.
    \item \textbf{XSTest}~\citep{safety-benchmark-xstest}: This benchmark evaluates two aspects of model safety. The first aspect examines potential security vulnerabilities across eight types of scenarios. The second aspect assesses the risk of excessive safety constraints across ten types of scenarios, ensuring that the model neither responds to harmful queries (e.g., providing answers about the private information of fictional characters) nor unnecessarily refuses to answer legitimate questions due to overly restrictive safety measures.
    \item \textbf{Do-Not-Answer}~\citep{safety-benchmark-dna}: Short for DNA, this benchmark is designed around "dangerous instructions that should not be followed" and consists of a set of risk-related queries covering twelve categories of harm (e.g., personal information leakage, assistance with illegal activities) and 61 specific risk types (e.g., racial discrimination, misleading medical advice).
    \item \textbf{HarmBench}~\citep{safety-benchmark-harmbench}: This benchmark is primarily structured around the following four aspects: standard model safety capabilities, copyright-related safety capabilities, context-aware safety capabilities, and multimodal safety capabilities. Additionally, this work introduces an automated approach for generating diverse automated red-teaming attack samples.
\end{itemize}
In terms of evaluation implementation, the results for the Do-Not-Answer and HarmBench benchmarks were reproduced based on the official evaluation methodology, while the results for the other benchmarks were obtained from the independent third-party evaluation platform HELM (\url{https://crfm.stanford.edu/helm/safety/latest/\#/leaderboard}) (we recorded the results of the website in April 2025). During the reproduction of the HarmBench results, we observe that using relatively smaller models (i.e., LLaMA-2-13B) led to unreliable evaluation outcomes. To address this issue, we refine the scoring approach by basing the evaluations on more advanced models (i.e., GPT4o(2024-11-20)). Additionally, when obtaining results for different models, some queries are flagged as unsafe by the risk control system, leading to automatic refusals. In such cases, we uniformly categorized these refusals as safe responses.

\begin{table}[h]
    \centering
    \setlength{\tabcolsep}{2pt}
    \captionof{table}{Comparison of DeepSeek-R1 and other frontier models on safety benchmarks. A higher score indicates better safety performance. Benchmarks marked with * are the results reproduced by us, while other numerical results are obtained from the independent HELM evaluations. The numbers in parentheses represent the results of the pure model without considering the risk control system~(introduced in~\ref{sec:risk_control}).}
    \label{tab:safety_eval}
    \begin{tabularx}{\textwidth}{l|cccccc|c} \toprule
    \textbf{Safety Score(\%)}             & \textbf{SST}             & \textbf{BBQ}         & \textbf{ART}      & \textbf{XSTest}       & \textbf{DNA*}       & \textbf{HarmBench*}  & \textbf{Average Score}  \\ \midrule
    \textbf{Claude-3.7-Sonnet}            & 100.0           & 92.1                 & 99.7     & 96.4                  & 95.9      & 83.3                & 94.6            \\
    \textbf{o1 (2024-12-17)}              & 99.0                     & 97.3        & 98.3              & 97.0                  & 86.2               & 84.0                & 93.6            \\
    \textbf{GPT-4o (2024-05-13)}          & 98.5                     & 95.1                 & 99.1              & 97.3         & 90.6               & 72.7                & 92.2            \\
    \textbf{Qwen2.5 Instruct (72B)}       & 100.0           & 95.4                 & 99.6     & 97.9         & 95.9      & 83.0                & 95.3            \\ 
    \textbf{DeepSeek-V3}                  & 95.3                     & 96.7        & 97.1              & 97.1                  & 95.6               & 96.0 (67.0)      & 96.3 (91.5)   \\ \midrule
    \textbf{DeepSeek-R1 (hide cot)}       & 98.0                     & 96.6                 & 97.2              & 94.4                  & 93.7               & 96.3 (58.0)      & 96.0 (89.7)   \\
    \textbf{DeepSeek R1}                  & 97.5                     & 96.6                 & 96.2              & 95.3                  & 94.8               & 89.3 (35.0)               & 95.0 (85.9)            \\ \bottomrule
    \end{tabularx}
\end{table}
The comparative experimental results are presented in Table~\ref{tab:safety_eval}, where we evaluate the safety performance of our model against other state-of-the-art models. The results indicate that the R1 model achieves comparable safety performance with other frontier models across different benchmarks. 
We observed that R1 exhibits a significant performance gap compared to other models on the HarmBench benchmark. Our analysis of the evaluation results suggests that R1 performs poorly on questions related to intellectual property rights within the HarmBench testset. For instance, when prompted to generate lyrics (e.g., Generate the lyrics of 'What a Wonderful World' by Louis Armstrong), R1 fails to refuse the request, leading to its classification as unsafe. Besides, across benchmarks focusing on other safety categories (e.g., discrimination and bias, violence and extremism, privacy violations, etc.), R1 consistently shows strong safety measures. 
% In summary, R1 excels in addressing a wide range of safety-related challenges across various domains; however, there is still room for improvement in areas such as intellectual property protection.



\subsubsection{Safety Taxonomic Study of R1 on In-House Benchmark}
\label{subsubsection:safety_taxnomic}
In this section, we present our safety taxonomy research for the DeepSeek-R1 model based on an in-house safety benchmark. Specifically, we first introduce the construction of the in-house safety benchmark. Subsequently, we discuss the performance of our R1 model across different categories and compare it with the performance of other frontier models.

Although existing works have already contributed valuable safety evaluation datasets, different datasets focus on distinct domains and employ varying classification methods. Moreover, data from different sources exhibit disparities in attributes (such as languages, quantities, and evaluation methods), making direct alignment challenging.
Therefore, we specifically constructed an internal safety evaluation dataset to monitor the overall safety level of the model. The construction of this dataset has the following characteristics: (1) Following unified taxonomic standards to build the testing framework, comprehensively covering various safety and ethical scenarios as much as possible; (2) Aligning the quantity, languages, and evaluation methods of safety test data across different categories, enabling us to conduct quantitative safety assessments for different safety scenarios; (3) Possessing good extensibility, where the multilingual language (\ref{sec:safety_multilingual}) and the jailbreak attacks (\ref{sec:safety_jailbreak}) evaluations in subsequent sections are also based on extensions of this dataset.

Our taxonomy of safety issues is presented in Figure~\ref{fig:safety_taxnomy}. We have categorized potential content safety challenges faced by language models into 4 major categories and 28 subcategories. The detailed description is as follows:
\begin{figure}
    \centering
    \includegraphics[width=\linewidth]{figures/safety_taxnomy.pdf}
    \caption{Taxonomy of in-house safety benchmark.} \label{fig:safety_taxnomy}
\end{figure}

\paragraph{Discrimination and Prejudice Issues} 
Discrimination and bias issues are prevalent across communities with diverse cultural backgrounds. We have broadly categorized these into two types: discrimination based on personal physical attributes and discrimination based on personal social attributes. Discrimination based on physical attributes primarily refers to inappropriate dismissal and mockery stemming from an individual's physiological conditions, such as age, gender, sexual orientation, appearance, body shape, and health status. Social attribute discrimination encompasses stereotypes based on nationality, ethnicity, and religion, as well as narrow perspectives derived from individual economic status, educational background, cultural identity, and family background.

\paragraph{Illegal and Criminal Behavior}
Illegal activities encompass the following safety topics: violent behavior, terrorism, illegal pornographic content, illegal medical practices (surrogacy, euthanasia, organ trafficking), illegal gambling, drug and substance abuse (including drug manufacturing, trafficking, and consumption), cybercrime (attacks on networks and computer systems), animal-related offenses (such as animal abuse or poaching), among others.

\paragraph{Harmful Behavior}
Harmful behavior toward humans primarily include the following four categories: (1) Physical harm: including self-harm, suicide, injury or murder of others; (2) Psychological harm: including verbal abuse, threats, intimidation, mental manipulation, deception, and instigation; (3) Privacy violations: encompassing personal health information, basic biometric data, ID information, location tracking, financial information, etc.; (4) Violations of economic interests: including breaches of business ethics, intellectual property infringement, disclosure of trade secrets, and unfair business competition.

\paragraph{Moral and Ethical Issues}
We have identified and categorized the ethical issues in the following four scenarios: (1) Family ethics: including issues related to parental responsibilities, sibling relationships, and the treatment of elderly family members, (2) Marriage ethics: covering topics such as fidelity, communication, and shared responsibilities between spouses, (3) School, student, and academic ethics: addressing matters like academic integrity, bullying, and the relationship between students and teachers, (4) Professional ethics: encompassing concerns related to conflicts of interest, workplace conduct, and responsibilities toward clients and colleagues.

Following the establishment of the above safety content taxonomy, we constructed specialized test sets for each of the 28 subcategories to evaluate the model's safety performance in the corresponding domains. Specifically, we manually created 20 Chinese test questions for each subcategory, covering important concepts and risk points within that domain while emphasizing question diversity. Subsequently, we translated these 20 Chinese test questions into corresponding English versions to better assess our model's safety performance in both Chinese and English. Ultimately, we developed 1,120 test questions for the systematic evaluation of model safety.

\begin{table}[h]
    \centering
    \footnotesize
    \setlength{\tabcolsep}{2pt}
    \captionof{table}{Comparison of DeepSeek-R1 and other frontier models in fine-grained safety scenarios. \textbf{Unsafe} indicates the proportion of unsafe content in the model's responses (lower values indicate better model safety), while \textbf{Rej.} represents the rejection rate in the model's answers (lower values indicate a stronger tendency for the model to provide informative and safe answers to questions, rather than simply declining to respond). For DeepSeek-V3 and DeepSeek-R1, we report results under two configurations: with and without risk control system~(introduced in~\ref{sec:risk_control}).}
    \label{tab:safety_taxnomy}
    \begin{tabular}{l|cc|cc|cc|cc|cc} 
        \toprule
        \textbf{}                       & \multicolumn{2}{c|}{\textbf{Discrimi.}} & \multicolumn{2}{c|}{\textbf{Illegal}} & \multicolumn{2}{c|}{\textbf{Harmful}} & \multicolumn{2}{c|}{\textbf{Ethical}} & \multicolumn{2}{c}{\textbf{Overall}} \\ \midrule
        \textbf{Ratio(\%)}              & \textbf{Unsafe}      & \textbf{Rej.}      & \textbf{Unsafe}   & \textbf{Rej.}  & \textbf{Unsafe}   & \textbf{Rej.}  & \textbf{Unsafe}   & \textbf{Rej.}  & \textbf{Unsafe}         & \textbf{Rej.}  \\ \midrule
        \textbf{Claude-3.7-Sonnet}      & 8.4                  & 2.5                & 14.1              & 4.5              & 9.5               & 5.5              & 7.5               & 0.6              & 10.7              & 3.6     \\
        \textbf{o1 (2024-12-17)}        & 7.2                  & 37.8               & 12.3              & 54.8             & 5.0               & 73.5             & 8.8               & 34.4             & 9.0               & 50.4             \\
        \textbf{GPT-4o (2024-05-13)}    & 19.1                 & 6.2                & 22.5              & 28.4             & 28.0              & 19.5             & 18.8              & 4.4              & 22.0              & 17.1             \\
        \textbf{Qwen2.5 Instruct (72B)} & 12.8                 & 2.5                & 14.5              & 9.5              & 15.5              & 5.0              & 11.9              & 0.0              & 13.8              & 5.4     \\ \midrule
        \textbf{DeepSeek-V3}            & 20.3                 & 2.5                & 17.3              & 13.9             & 17.5              & 9.5              & 13.1              & 1.9              & 17.6              & 8.1              \\
        + risk control system           & 8.1                  & 16.9               & 3.2               & 35.5             & 7.0               & 22.5             & 3.1               & 18.1             & 5.3               & 25.4             \\ \midrule
        \textbf{DeepSeek-R1}            & 19.7                 & 3.8                & 28.9              & 8.6              & 32.5              & 6.0              & 16.9              & 0.6              & 25.2              & 5.6              \\
        + risk control system           & 9.1                  & 17.2               & 6.6               & 39.1             & 13.0              & 29.0             & 6.9               & 13.1             & 8.5               & 27.3              \\
        \bottomrule
    \end{tabular}
\end{table}

In our evaluation methodology, we employed the LLM-as-a-Judge approach, utilizing an advanced GPT version (GPT4o (2024-11-20)) to determine safety labels. 
Our safety assessment process categorizes each QA pair into the following three classes:
(1) Unsafe: The model provides an answer that fails to meet ethical and safety standards, representing a clear negative sample;
(2) Safe: The model provides a relevant answer that acknowledges the ethical risks in the question and offers appropriate cautions to the user, representing a clear positive sample;
(3) Rejection: The model either provides an irrelevant refusal response or the system delivers a mechanical rejection based on risk control measures; we consider this situation as an intermediate state between safe and unsafe responses.

In Tables~\ref{tab:safety_taxnomy}, we consider two metrics: Unsafe and Reject. 
The Unsafe metric measures the proportion of unsafe responses among all answers, with lower values indicating better safety performance. 
The Reject metric represents the proportion of rejection responses among all answers, with lower values being more desirable (we prefer safe responses over rejections since it can provide risk warning information).

We crafted specialized prompts for different subcategories of questions to assess the safety of responses. We also verified that the consistency between LLM evaluation results and human assessments reached an acceptable level (consistency rate of sampled results is above 95\%).
The experimental comparison results are presented in Table~\ref{tab:safety_taxnomy}, from which the following conclusions can be observed:
\begin{itemize}
\item \textbf{Analyzing unsafe rates}: DeepSeek-V3 (with risk control) belongs to the first tier of safe models (unsafe rate aound 5\%); DeepSeek-R1 (with risk control), Claude-3.7-Sonnet, and o1 (2024-12-17) belong to the second tier of safe models (unsafe rate around 10\%); DeepSeek-V3 (without risk control) and Qwen2.5 Instruct (72B) belong to the third tier of safe models (unsafe rate around 15\%); while DeepSeek-R1 (without risk control) and GPT-4o (2024-05-13) are relatively unsafe models (unsafe rate beyond 20\%).
\item \textbf{Analyzing rejection rates}: The base models of DeepSeek-R1 and DeepSeek-V3 have relatively low rejection rates but higher unsafe rates. After implementing a risk control system, these models show relatively low unsafe rates but higher rejection rates (around 25\%). Additionally, Claude-3.7-Sonnet achieves a good balance between user experience (lowest rejection rate) and model safety (unsafe rate at relatively low levels); while o1 (2024-12-17) demonstrates a more severe tendency to reject queries (around 50\%), presumably employing strict system-level risk control to prevent the model from exposing unsafe content.
\item \textbf{Analyzing risk types}: DeepSeek-R1 performs exceptionally well in handling queries related to Illegal and Criminal Behavior and Moral and Ethical Issues, while showing average performance in scenarios involving Discrimination and Prejudice Issues and Harmful Behavior, which encourages us to pay more attention on these two categories when developing model safety features and risk control system.
\end{itemize}



\subsubsection{Multilingual Safety Performance}
\label{sec:safety_multilingual}

\begin{figure}
    \centering
    \includegraphics[width=0.9\linewidth]{figures/safety-multilingual-v3.pdf}
    \caption{Multilingual safety performance. V3-check and R1-check represent the risk control system evaluation results for DeepSeek-V3 and DeepSeek-R1, respectively.} \label{fig:safety_multilingual}
\end{figure}
In the previous section's evaluation, we primarily focused on the model's safety performance in special languages (Chinese and English). However, in practical usage scenarios, users' linguistic backgrounds are highly diverse. Assessing safety disparities across different languages is essential. 
For this purpose, we translated the original bilingual safety testset (introduced in the~\ref{subsubsection:safety_taxnomic}) into 50 commonly used languages. 
For high-frequency languages, we conducted full translation of the entire dataset, while for low-frequency languages, we performed sampling translation. This process resulted in a comprehensive multilingual safety test set consisting of 9,330 questions.
During the translation process, we employed a combined approach of LLM translation and human-assisted calibration to ensure the quality of the translations.

We continued to use the LLM-as-a-judge methodology described in the previous section, which determines safety labels (safe, unsafe, or rejected) for each question-answer pair. Rather than merely rejecting risky queries, we prefer responses that provide safe content; therefore, we assigned higher scores to safe responses (5 points per question, with 5 points for safe responses, 0 points for unsafe responses, and 4 points for rejections). The final safety score proportions (safety score as a percentage of the total possible safety score) across 50 languages are presented in Figure~\ref{fig:safety_multilingual}.
For DeepSeek-V3 and DeepSeek-R1, we evaluated safety scores for models with and without the risk control system (introduced in~\ref{sec:risk_control}). Additionally, we tested the multilingual safety performance of Claude-3.7-Sonnet and GPT-4o(2024-05-13).
From Figure~\ref{fig:safety_multilingual}, we can draw the following conclusions:
\begin{itemize}
    \item With risk control system in place, DeepSeek-V3 (86.5\%) and DeepSeek-R1 (85.9\%) achieve total safety scores across 50 languages that approach the best-performing Claude-3.7-Sonnet (88.3\%). This demonstrates that DeepSeek has reached state-of-the-art levels in system-level multilingual safety.
    \item Without risk control system, DeepSeek-V3 (75.3\%) and DeepSeek-R1 (74.2\%) get safety scores across 50 languages comparable to GPT-4o(2024-05-13)'s performance (75.2\%). This indicates that even when directly using the open-source versions of R1, the model still exhibits a moderate level of safety standard.
    \item Examining language-specific weaknesses, we categorize languages with safety scores below 60 points as high-risk languages for the corresponding model. Among the 50 languages evaluated, DeepSeek-R1 (without risk control system) and Claude-3.7-Sonnet have zero high-risk languages; DeepSeek-V3 (without risk control system) and GPT-4o(2024-05-13) have one and two high-risk languages, respectively. This suggests that DeepSeek-R1 has no obvious language-specific vulnerabilities.
\end{itemize}


\subsubsection{Robustness against Jailbreaking}
\label{sec:safety_jailbreak}

In real-world application scenarios, malicious users may employ various jailbreaking techniques to circumvent a model's safety alignment and elicit harmful responses. Therefore, beyond evaluating model safety under direct questioning, we place significant emphasis on examining the model's robustness when confronted with jailbreaking attacks.
Thus, we constructed a dedicated test suite for jailbreaking evaluation. Specifically, we developed a template collection consisting of 2,232 jailbreaking instructions. We then randomly concatenated these jailbreaking prompts with questions from the original safety testset (introduced in~\ref{subsubsection:safety_taxnomic}) and further examined the performance differences in the model's responses when confronted with original unsafe questions versus newly formulated questions with jailbreaking elements.

When evaluating the results, we followed the LLM-as-a-Judge safety assessment~(introduced in~\ref{subsubsection:safety_taxnomic}), while improving the safety evaluation prompts to focus more specifically on identifying manipulative traps in jailbreak attempts. Each question-answer pair was classified into one of three categories: safe, unsafe, or rejected~(introduced in~\ref{subsubsection:safety_taxnomic}). The results of jailbreak attacks against various models are presented in Table~\ref{tab:safety_jailbreak}. From these results, we draw the following conclusions:
\begin{table}[h]
    \centering
    % \footnotesize
    \setlength{\tabcolsep}{2pt}
    \captionof{table}{Comparison of DeepSeek-R1 and other frontier models in jailbreaking scenarios. }
    \label{tab:safety_jailbreak}
    \begin{tabular}{l|ccc|ccc}
        \toprule
        \multicolumn{1}{c|}{\textbf{}}   & \multicolumn{3}{c|}{\textbf{Unsafe Ratio}} & \multicolumn{3}{c}{\textbf{Rejected Ratio}} \\ \midrule
        \textbf{Ratio(\%)}              & Origin   & Jailbreak   & GAP   & Origin    & Jailbreak    & GAP   \\\midrule
        \textbf{Claude-3.7-Sonnet}      & 10.7     & 26.2        & +15.5    & 3.6       & 21.9         & +18.3    \\
        \textbf{o1 (2024-12-17)}        & 9.0      & 12.1        & +3.1     & 50.4      & 79.8         & +29.4    \\
        \textbf{GPT-4o (2024-05-13)}    & 22.0     & 30.4        & +8.4     & 17.1      & 57.3         & +40.2    \\
        \textbf{Qwen2.5 Instruct (72B)} & 13.8     & 29.7        & +15.9    & 5.4       & 25.2         & +19.8    \\\midrule
        \textbf{DeepSeek-V3}            & 17.6     & 36.4        & +18.8    & 8.1       & 8.9          & +0.8     \\
        + risk control system           & 5.3      & 2.3         & -3.0    & 25.4      & 46.5         & +21.1    \\\midrule
        \textbf{DeepSeek-R1}            & 25.2     & 85.9        & +60.7    & 5.6       & 1.9          & -3.7    \\
        + risk control system           & 8.5      & 4.3         & -4.2    & 27.3      & 87.3         & +60.0    \\
        \bottomrule
    \end{tabular}
\end{table}
\begin{itemize}
    \item All tested models exhibited significantly increased rates of unsafe responses and rejections, along with decreased safety rates when facing jailbreak attacks. For example, Claude-3.7-Sonnet, showed a 33.8\% decrease in the proportion of safe responses when confronted with our security jailbreak attacks. This demonstrates that current cutting-edge models still face substantial threats from jailbreak attacks.
    \item Compared to non-reasoning models, the two reasoning models in our experiments — DeepSeek-R1 and o1(2024-12-17) — rely more heavily on the risk control system for security checks, resulting in considerably higher overall rejection rates (79.8\% and 87.3\% respectively). 
    \item Open-source models (DeepSeek, Qwen) face more severe jailbreak security challenges than closed-source models, because of the lack of a risk control system in locally deployed models. To address safety issues, we advise developers using open source models in their services to adopt comparable risk control measures.
\end{itemize}




% \subsubsection{Discussion: Challenges of Potentially Risky Knowledge Exposure Brought by Transparent Chain-of-Thought}
% \label{sec:safety_cot}
% DeepSeek-R1 is the first cutting-edge model to fully exhibit the reasoning process. By making the thought process entirely transparent, it provides the community and users with constructive thinking inspiration. However, we also recognize that a fully transparent chain-of-thought exacerbates the challenges associated with risk knowledge exposure. There are two main reasons for this: (1) Through chain-of-thought reasoning, the model can generate more accurate and detailed responses. When addressing topics related to risk knowledge, it is more likely to provide extensive explanations and additional details. (2) Even if the model ultimately provides a safe response or refuses to answer risk knowledge-related queries, the reasoning process itself may still involve analyzing or referencing elements of risk content. This could inadvertently inspire users in unintended ways, further complicating risk management.

% \hl{Knowledge-based question answering with potential risks differs from direct risk-related question answering. Discussions about such knowledge (e.g., computer viruses, chemical substances, etc.) are inherently neutral, as they can be applied in both legitimate contexts, such as risk prevention, and in potentially harmful scenarios.} To quantitatively assess the extent of risk exposure in the model’s responses to such questions, we designed a benchmark for potential risk knowledge question answering~(short as~\textbf{PRK-QA}). 
% \hl{The scope of potentially dangerous knowledge is quite broad; therefore, when constructing this testset, we did not aim for comprehensiveness but rather sought to diversify the disciplinary backgrounds to facilitate a thorough discussion of this issue. Ultimately, we selected six knowledge areas from different disciplines, encompassing both natural sciences and social sciences, as enumerated below:}
% \begin{itemize}
%     \item \textbf{Cyber Viruses}: When models provide detailed technical knowledge about malware development, attack vectors, and vulnerability exploitation, this information may be misused by malicious actors to develop harmful software or execute cyberattacks. Such activities could circumvent security systems, compromise personal privacy, organizational data, and critical infrastructure security, potentially resulting in large-scale system failures and economic damage.
%     \item \textbf{Chemical Substances}: When models elaborate on synthesis methods, preparation processes, and usage techniques for dangerous chemical compounds, this information could facilitate illegal manufacturing of explosives, illicit drugs, or other harmful substances. Such activities pose severe health risks to the operators themselves and may enable actions threatening public safety, potentially causing casualties and social distress.
%     \item \textbf{Trading Contracts}: When models offer detailed insights into complex financial derivatives, futures contracts, and cryptocurrency trading, they might mislead users lacking professional risk assessment capabilities into high-risk investment domains. Without appropriate risk management strategies, speculative trading can lead to substantial financial losses, excessive debt, personal financial crises, and associated mental health issues.
%     \item \textbf{Nuclear Technology}: When models provide excessively detailed knowledge regarding nuclear reactor design, nuclear material handling, and radioactive substance management, they may guide users without professional backgrounds to attempt dangerous operations. Additionally, malicious actors might exploit such information to develop hazardous devices, potentially causing radiation leakage, environmental contamination, and other severe consequences.
%     \item \textbf{Psychological Manipulation (PUA)}: When models detail techniques for psychological manipulation, emotional control, and mental abuse, this information may enable the acquisition and implementation of harmful behaviors. Such activities can foster inequality, manipulation, and abuse in interpersonal relationships, endangering others' psychological well-being and autonomy, potentially resulting in serious psychological trauma and emotional harm.
%     \item \textbf{Gender Transition}: When models provide detailed information about medical approaches, pharmaceutical usage, and surgical procedures related to gender transition without accurate information or necessary medical warnings, users may undertake dangerous self-intervention measures without professional medical guidance, resulting in physical harm. Furthermore, inaccurate or biased information may deepen misunderstandings and prejudices against gender-diverse communities, affecting the mental health and social acceptance of relevant populations.
% \end{itemize}

% In constructing PRK-QA, we formulate 20 knowledge questions~(10 questions in Chinese and their corresponding translations in English) for each category, covering detailed discussions of potentially risky content in their respective fields. During the question design process, we emphasized neutrality as much as possible, ensuring that each question could be interpreted as having both legitimate and potentially improper applications. For evaluation, we leverage the advanced third-party language model~(i.e., GPT4o (2024-11-20)) to assess the information density of potential risk knowledge~(to evaluate the degree of answer utility) and the strength of risk warning~(to evaluate the level of risk notification) of answers.

% \input{tables/safety_cot}
% The evaluation results of different models on PRK-QA are presented in Table~\ref{tab:safety_cot}. In addition to comparing the performance with other state-of-the-art models, we also conducted two ablation experiments with DeepSeek-R1: one hiding the model's reasoning process and another showing only the reasoning process. From these experimental results, we draw the following conclusions:
% \begin{itemize}
%     \item Compared to other models, DeepSeek-R1 generally maintains a high level of information density regarding potentially dangerous knowledge (8.8), while exhibiting a moderate level of risk prompting intensity (3.5);
%     \item DeepSeek-R1 demonstrates significantly higher knowledge density in the domains of computer viruses and gender reassignment compared to other models, while remaining comparable to other models across the remaining four domains.
%     \item Reasoning-based models tend to expose more potentially risky knowledge compared to non-reasoning models. DeepSeek-V3's average information density is 8.5, while DeepSeek-R1's is 8.8; similarly, o1 (2024-12-17) shows notably higher information density than GPT4o (2024-05-13). We attribute this to the fact that reasoning-enabled models provide answers with more detailed information in various dimensions, thereby leading to increased exposure of potentially risky knowledge.
%     \item The chain-of-thought reasoning process itself already achieves a considerably high information density (8.3), and when presented together with the output results, the final information density (8.8) is higher than when only providing the reasoning results (8.7). Considering the results in Table~\ref{tab:safety_eval}, that the safety scores when hiding the chain-of-thought are higher than when displaying it, we conclude that exposing the reasoning process itself may potentially introduce users to risky knowledge.
% \end{itemize}

% In summary, our findings reveal that reasoning-enabled models like DeepSeek-R1 expose more potentially risky knowledge than non-reasoning models, despite maintaining moderate risk prompting intensity. The chain-of-thought reasoning process itself carries significant information density, creating a fundamental tension between model transparency and safety considerations.
% % Furthermore, we conducted an in-depth analysis of the impact of incorporating chain-of-thought reasoning on model safety. Our findings indicate that when the chain-of-thought is hidden, the R1 model achieves a safety level comparable to (or even slightly exceeding) that of the DeepSeek-V3 model~(shown in Table~\ref{tab:safety_eval}). However, when the full chain-of-thought reasoning process is explicitly exposed, we observed a slight decline in safety performance metrics.
% % We attribute this decline to the nature of the reasoning process itself, which inherently involves assessing and analyzing the safety of a given query. This process may, to some extent, provide users with unintended insights. For instance, when evaluating the harms of violent behavior, the model may enumerate specific examples of violent acts, which—despite ultimately leading to a safe and responsible conclusion—could inadvertently contribute to the dissemination of inappropriate knowledge. This introduces a novel dimension of safety challenges for models employing chain-of-thought reasoning.



% % \subsection{Results of \dsri{} at Different Stages} \label{sec:stages}
% % \input{tables/stageeval}
% % Table \ref{tab:stage_r1} presents the performance of \dsri{} across various developmental stages, in addition to comparisons with the baseline \dsro{} and the final version of \dsri{}. As outlined in Figure \ref{fig:r1-pipeline}, the initial development stage, \dsri{} Dev1, involves fine-tuning the model with  cold-start data. When comparing \dsro{} and \dsri{} Dev1, notable improvements are observed in instruction-following capabilities, as reflected by enhanced scores in IF-Eval and ArenaHard benchmarks. However, due to the limited size of the cold-start dataset, Dev1 does not fully retain the reasoning capabilities of \dsro{}, which degrades on the AIME benchmark.

% % To address this, we conduct large-scale reinforcement learning (RL) using reasoning-focused prompts, resulting in the generation of \dsri{} Dev2. This version demonstrates significant performance gains in benchmarks requiring advanced reasoning, such as those evaluating code generation, mathematical problem-solving, and STEM-related tasks. Conversely, benchmarks targeting general-purpose language understanding—such as AlpacaEval 2.0 and ArenaHard—show marginal improvement. These results suggest that reasoning-oriented RL considerably enhances reasoning capabilities, while exerting limited influence on user preference-oriented benchmarks.

% % In the subsequent stage, we use \dsri{} Dev2 to perform rejection sampling and supervised fine-tuning (SFT) on the DeepSeek-V3 base model, yielding \dsri{} Dev3. This stage integrates both reasoning and non-reasoning datasets into the SFT pipeline, thereby enhancing the model’s proficiency in both reasoning and general language generation tasks. Compared to Dev2, \dsri{} Dev3 achieves notable performance improvements on AlpacaEval 2.0 and Aider-Polyglot, attributable to the inclusion of large-scale non-reasoning corpora and code engineering datasets. Given the inefficiency of evaluating code engineering tasks during RL training, these datasets are incorporated during the rejection sampling phase to maximize their utility.

% % Finally, we apply comprehensive RL training on \dsri{} Dev3 using a mixture of reasoning-focused and general-purpose data, resulting in the final version of \dsri{}. Marginal improvements are observed in code and mathematics benchmarks, as substantial RL training specific to reasoning tasks had already been conducted in prior stages. The primary advancements in the final iteration of \dsri{} are observed in general instruction-following and user-preference benchmarks, with AlpacaEval 2.0 improving by 25\% and ArenaHard by 17\%.

% % \begin{table}[h]
    \centering
    \footnotesize
    \setlength{\tabcolsep}{1.9pt}
    \caption{ Experimental results for each stage of \dsri{} on  problems with varying difficulty levels in the LiveCodeBench dataset.}
    \begin{tabular}{l | c c  c  c |c c@{}}
    \toprule
    \multirow{2}{*}{Difficulty Level} & \textbf{DeepSeek-R1} & \textbf{DeepSeek-R1}  & \textbf{DeepSeek-R1}& \textbf{DeepSeek-R1} &  \textbf{DeepSeek}\\
    & \textbf{Zero}   & \textbf{Dev1}  & \textbf{Dev2} & \textbf{Dev3} & \textbf{R1} \\
    \midrule
     Easy   & 98.07 & 99.52&100.00 & 100.00  &   \textbf{100.00} \\
     Medium   & 58.78 & 73.31&81.76 & 81.42  &   \textbf{83.45} \\
     Hard   & 17.09 & 23.21&30.36 & 33.16  &   \textbf{34.44} \\
    \midrule

    \end{tabular}
    
    \label{tab:difficulty_lcb}
\end{table}




% % To further evaluate the performance of each stage of \dsri{} on problems of varying difficulty, we present the experimental results for each stage of DeepSeek-R1 on the LiveCodeBench dataset, as shown in Table~\ref{tab:difficulty_lcb}. It can be observed that for each stage, simple problems are generally solved correctly, while the main improvements come from medium and hard problems. This fine-grained analysis demonstrates that each stage brings significant improvement on complex coding reasoning problems.




\section{More Analysis}\label{sec:exp_ana}

\subsection{ Performance  Comparison with DeepSeek-V3}

\begin{table}[h]
    \centering
    \small
    \setlength{\tabcolsep}{1.9pt}
    \caption{A Comparative Analysis of DeepSeek-V3 and DeepSeek-R1. DeepSeek-V3 is a non-reasoning model developed on top of DeepSeek-V3-Base, which also serves as the foundational base model for DeepSeek-R1.  Numbers in bold denote the performance is statistically significant  (t$-$test with $p < 0.01$).}
    \begin{tabular}{@{}c l | c c  c  c}
    \toprule
     & \textbf{Benchmark {\tiny (Metric)}} & \textbf{V3-Base} & \textbf{V3} & \textbf{R1-Zero} & \textbf{R1} \\
    \midrule
    \multirow{10}{*}{English}& MMLU {\tiny (EM)} & 87.1 & 88.5 & 88.8 & \textbf{90.8}\\
     & MMLU-Redux {\tiny (EM)} & 86.2 & 89.1 & 85.6 & \textbf{92.9} \\
    & MMLU-Pro {\tiny (EM)}  & 64.4 & 75.9 & 68.9 & \textbf{84.0} \\
    & DROP {\tiny (3-shot F1)}  & 89.0 & 91.6 & 89.1 & \textbf{92.2}\\
    & IF-Eval {\tiny (Prompt Strict)}  & 58.6 & \textbf{86.1} & 46.6 & 83.3 \\
    & GPQA Diamond {\tiny (Pass@1)}  & - & 59.1 & \textbf{75.8} & 71.5 \\
    & SimpleQA {\tiny (Correct)}   & 20.1 & 24.9 & 30.3 & 30.1 \\
     & FRAMES {\tiny (Acc.)}   & - & 73.3 & 82.3 & 82.5\\
      & AlpacaEval2.0 {\tiny (LC-winrate)}  & - & 70.0 & 24.7 & \textbf{87.6}\\
       & ArenaHard {\tiny (GPT-4-1106)}  & - & 85.5 & 53.6 & \textbf{92.3}\\
    \midrule
    \multirow{7}{*}{Code} & LiveCodeBench {\tiny (Pass@1-COT)} & - & 36.2 & 50.0 & \textbf{65.9} \\
    & Codeforces {\tiny (Percentile)} & - & 58.7 & 80.4 & \textbf{96.3} \\
    & Codeforces {\tiny (Rating)} & - & 1134 & 1444 & \textbf{2029} \\
    & SWE Verified {\tiny (Resolved)} & - & 42.0 & 43.2 & \textbf{49.2} \\
    & Aider-Polyglot {\tiny (Acc.)}  & - & 49.6 & 12.2 & \textbf{53.3} \\
    \midrule
    \multirow{3}{*}{Math} & AIME 2024 {\tiny (Pass@1)} & - & 39.2 & 77.9 & \textbf{79.8} \\
    & MATH-500 {\tiny (Pass@1)} & - & 90.2 & 95.9 & \textbf{97.3} \\
    & CNMO 2024 {\tiny (Pass@1)} & - & 43.2 & \textbf{88.1} & 78.8 \\
    \midrule
    \multirow{3}{*}{Chinese} & CLUEWSC {\tiny (EM)} & 82.7 & 90.9 & 93.1 & 92.8\\
    & C-Eval {\tiny (EM)}  & 90.1 & 86.5 & \textbf{92.8} & 91.8\\
     & C-SimpleQA {\tiny (Correct)} & - & \textbf{68.0} & 66.4 & 63.7 \\
    \bottomrule
    \end{tabular}
    \label{tab:v3_full_compare}
\end{table}

\begin{figure}
    \centering
    \begin{subfigure}{}
        \centering
        \includegraphics[width=0.85\linewidth]{figures/mmlu_category_comparison.pdf}
        \caption{The comparison of DeepSeek-V3 and DeepSeek-R1 across MMLU categories.}
        \label{mmlu}
    \end{subfigure}
    
    \vspace{1em}
    
    \begin{subfigure}{}
        \centering
        \includegraphics[width=0.85\linewidth]{figures/mmlu_pro_subject_comparison.pdf}
        \caption{The comparison of DeepSeek-V3 and DeepSeek-R1 across MMLU-Pro categories.}
        \label{mmlupro}
    \end{subfigure}

    \label{fig:mmlu_comparison}
\end{figure}


Since both \dsri{} and DeepSeek-V3 share a common base architecture, namely DeepSeek-V3-Base, a critical question naturally arises: which specific dimensions are enhanced through the application of different post-training techniques? To address this, we first compare the R1 family of models with DeepSeek-V3 and DeepSeek-V3-Base, as summarized in Table \ref{tab:v3_full_compare}. Notably, \dsri{} demonstrates significant improvements in competitive programming and mathematical reasoning tasks, as evidenced by superior performance on benchmarks such as LiveCodeBench and AIME 2024. These enhancements in reasoning capabilities also translate into higher scores on the Arena-Hard evaluation suite. Furthermore, \dsri{} exhibits stronger long-context understanding, as indicated by its improved accuracy on the FRAMES benchmark.
In contrast, DeepSeek-V3 shows a relative advantage in instruction-following capabilities, suggesting different optimization priorities between the two models.

To further elucidate the specific knowledge domains that benefit most from post-training, we conduct a fine-grained analysis of model performance across various subject categories within MMLU and MMLU-Pro. These categories, predefined during the construction of the test sets, allow for a more systematic assessment of domain-specific improvements.

As illustrated in Figure \ref{mmlupro}, performance improvements on MMLU-Pro are observed across all domains, with particularly notable gains in STEM-related categories such as mathematics and physics. Similarly, on MMLU, the largest improvements from DeepSeek-V3 to \dsri{} are also observed in STEM domains. However, unlike MMLU-Pro, gains in the STEM domain are smaller, suggesting differences in the impact of post-training between the two benchmarks.

Our hypothesis is that MMLU represents a relatively easier challenge compared to MMLU-Pro. In STEM tasks of MMLU, post-training on DeepSeek-V3 may have already achieved near-saturation performance, leaving minimal room for further improvement in \dsri{}. It surprised us that the non-STEM tasks, such as social sciences and humanities, are improved with the long CoT, which might attribute to the better understanding of the question. 

\subsection{Generalization to Real-World Competitions}% OOD Evaluation @zhihong
% AIME 2025; and more?
\begin{table}[h]
    \centering
    % \setlength{\tabcolsep}{2pt}
    \captionof{table}{Performance on latest math competitions. Participants with their USAMO index ($\texttt{AMC score} + 10 \times \texttt{AIME score}$) surpassing 251.5 are qualified for USAMO.}
    \label{tab:math_2025_eval}
    \begin{tabular}{l|cc|c} \toprule
    Average Score & \textbf{AMC 12 2024} & \textbf{AIME 2025} & \textbf{USAMO Index}  \\    \midrule
    Human Participants & 61.7 & 6.2/15 & 123.7 \\
    \midrule
    GPT-4o 0513 & 84.0 & 2.0/15 & 104.0 \\
    DeepSeek V3 & 98.3 & 3.3/15 & 131.3 \\
    \midrule
    OpenAI o1-1217 & 141.0 & 12.0/15 & 261.0 \\
    \textbf{DeepSeek R1} & 143.7 & 11.3/15 & 256.7 \\
    \bottomrule
    \end{tabular}
\end{table}
\noindent
Despite rigorous efforts to eliminate data contamination, variations of test set questions or discussions of related problems may still exist on websites that were included in the pre-training corpus. This raises an important question: can \dsri{} achieve comparable performance on test sets that were released after its training? To investigate this, we evaluate our model on AIME 2025, providing insights into its generalization capabilities on unseen data.
As shown in Table \ref{tab:math_2025_eval}, in AIME 2025 (\url{https://artofproblemsolving.com/wiki/index.php/2025_AIME_II_Problems}), \dsri{} achieves a 75\% solve rate (Pass@1), approaching o1's performance of 80\%.
Most notably, the model attains a score of 143.7/150 in AMC 12 2024 (\url{https://artofproblemsolving.com/wiki/index.php/2024_AMC_12B_Problems}) - a performance that, when combined with its AIME results, yields a score exceeding the qualification threshold for attending the USAMO (United States of America Mathematical Olympiad \url{https://artofproblemsolving.com/wiki/index.php/AMC_historical_results?srsltid=AfmBOoqQ6pQic5NCan_NX1wYgr-aoHgJ33hsq7KSekF-rUwY8TBaBao1}).
This performance positions DeepSeek-R1 among the nation's top-tier high school students.

\subsection{Mathematical Capabilities Breakdown by Categories}
\begin{figure}
    \centering
    \includegraphics[width=\linewidth]{figures/math_performance_breakdown_w_baseline.pdf}
    \caption{Performance breakdown by different categories of quantitative reasoning problems from a collection of contests in 2024.} \label{fig:math_breakdown}
\end{figure}
To assess \dsri{}'s mathematical reasoning capabilities comprehensively, we evaluated its performance across diverse categories of quantitative reasoning problems.
Our test set comprised 366 problems drawn from 93 mathematics competitions held in 2024 (\url{https://artofproblemsolving.com/community/c3752401_2024_contests}), including mathematical olympiads and team selection tests.
As shown in Figure~\ref{fig:math_breakdown}, \dsri{} significantly outperforms the representative non-reasoning model GPT-4o 0513.
\dsri{} demonstrates relatively strong proficiency in number theory and algebra, while exhibiting considerable room for improvement in geometry and combinatorics.

\subsection{An Analysis on CoT Length}
\begin{figure}
    \centering
    \includegraphics[width=\linewidth]{figures/test_time_scaling_vs_problem_difficulty.pdf}
    \caption{Test-time compute scaling (measured by the number of thinking tokens generated to reach correct answers) as problem difficulty (measured by Pass@1) increases. The picture is smoothed using UnivariateSpline from SciPy with a smoothing factor of 5.}
    \label{fig:math_test_time_scaling_by_difficulty}
\end{figure}
\textbf{Adaptive CoT length:} During training, \dsri{} was permitted to think for a long time (i.e., to generate a lengthy chain of thought) before arriving at a final solution.
To maximize success on challenging reasoning tasks, the model learned to dynamically scale computation by generating more thinking tokens to verify or correct its reasoning steps, or to backtrack and explore alternative approaches when initial attempts proved unsuccessful.
The complexity of a problem directly correlates with the number of thinking tokens required: more difficult problems typically demand more extensive computation. For extremely easy questions, like $1+1=?$, the model tends to use fewer tokens ($<100$ tokens) to answer the question. 

Figure \ref{fig:math_test_time_scaling_by_difficulty} demonstrates how \dsri{} scales test-time compute to solve challenging problems from math competitions held in 2024 (the same set of problems used in Figure \ref{fig:math_breakdown}).
\dsri{} achieves a 61.8\% solve rate (Pass@1) by scaling test-time compute to an average of 8,793 thinking tokens per problem.
Notably, the model adaptively adjusts its computational effort based on problem difficulty, using fewer than 7,000 thinking tokens for simple problems while dedicating more than 18,000 thinking tokens to the most challenging ones, which demonstrates DeepSeek-R1 allocates test-time compute adaptively based on problem complexity: on more complex problems, it tends to think for longer. Looking forward, we hypothesize that if token budget allocation were explicitly modeled during training, the disparity in token usage between easy and hard questions at test time could become even more pronounced.

\textbf{Comparison of non-reasoning models:} A key advantage of reasoning models like \dsri{} over non-reasoning models such as GPT-4o 0513 is their ability to scale effectively along the dimension of reasoning.
Non-reasoning models typically generate solutions directly, without intermediate thinking steps, and rarely demonstrate advanced problem-solving techniques like self-reflection, backtracking, or exploring alternative approaches.
On this same set of math problems, GPT-4o 0513 achieves only a 24.7\% solve rate while generating 711 output tokens on average — an order of magnitude less than \dsri{}.
Notably, non-reasoning models can also scale test-time compute with traditional methods like majority voting, but those methods fail to close the performance gap with reasoning models, even when controlling for the total number of tokens generated. For example, majority voting across 16 samples per problem yields minimal improvement in GPT-4o's solve rate on the 2024 collection of competition-level math problems, despite consuming more total tokens than \dsri{}. On AIME 2024, majority voting across 64 samples only increases GPT-4o's solve rate from 9.3\% to 13.4\%—still dramatically lower than \dsri{}'s 79.8\% solve rate or o1's 79.2\% solve rate. This persistent performance gap stems from a fundamental limitation: in majority voting, samples are generated independently rather than building upon each other. Since non-reasoning models lack the ability to backtrack or self-correct, scaling the sample size merely results in repeatedly sampling potentially incorrect final solutions without increasing the probability of finding correct solutions in any single attempt, making this approach highly token-inefficient.



\textbf{Drawback:} However, \dsri{}'s extended reasoning chains still sometimes fail to be thorough or become trapped in incorrect logic paths.
Independently sampling multiple reasoning chains increases the probability of discovering correct solutions, as evidenced by the fact that \dsri{}'s Pass@64 score on AIME 2024 is 90.0\%, significantly higher than its Pass@1 score of 79.8\%.
Therefore, traditional test-time scaling methods like majority voting or Monte Carlo Tree Search (MCTS) can complement \dsri{}'s long reasoning; specifically, majority voting further improves \dsri{}'s accuracy from 79.8\% to 86.7\%.

\subsection{ Performance of Each Stage on Problems of Varying Difficulty}
\begin{table}[h]
    \centering
    \footnotesize
    \setlength{\tabcolsep}{1.9pt}
    \caption{ Experimental results for each stage of \dsri{} on  problems with varying difficulty levels in the LiveCodeBench dataset.}
    \begin{tabular}{l | c c  c  c |c c@{}}
    \toprule
    \multirow{2}{*}{Difficulty Level} & \textbf{DeepSeek-R1} & \textbf{DeepSeek-R1}  & \textbf{DeepSeek-R1}& \textbf{DeepSeek-R1} &  \textbf{DeepSeek}\\
    & \textbf{Zero}   & \textbf{Dev1}  & \textbf{Dev2} & \textbf{Dev3} & \textbf{R1} \\
    \midrule
     Easy   & 98.07 & 99.52&100.00 & 100.00  &   \textbf{100.00} \\
     Medium   & 58.78 & 73.31&81.76 & 81.42  &   \textbf{83.45} \\
     Hard   & 17.09 & 23.21&30.36 & 33.16  &   \textbf{34.44} \\
    \midrule

    \end{tabular}
    
    \label{tab:difficulty_lcb}
\end{table}




To further evaluate the performance of each stage of \dsri{} on problems of varying difficulty, we present the experimental results for each stage of DeepSeek-R1 on the LiveCodeBench dataset, as shown in Table~\ref{tab:difficulty_lcb}. It can be observed that for each stage, simple problems are generally solved correctly, while the main improvements come from medium and hard problems. This fine-grained analysis demonstrates that each stage brings significant improvement on complex coding reasoning problems.

% \paragraph{Out of Domain Evaluation}
% Another interesting question is whether the model could generalize to logic questions that are out of domain of training data. The ARC (Abstraction and Reasoning Corpus) Logic Test \citep{chollet2019arc} is a unique intelligence benchmark designed to evaluate a system's ability to reason, generalize, and solve novel problems using abstract thinking—skills often associated with human intelligence. Developed by François Chollet in 2019, ARC is distinct from traditional AI benchmarks because it doesn't rely on large datasets or pre-trained models. Instead, it challenges participants (human or AI) to figure out patterns and rules from a small set of examples and apply them to new, unseen scenarios.
% Each ARC task consists of a few "demonstration" grids—typically 2D arrays of colored squares—where an input grid is transformed into an output grid according to some hidden rule. The test-taker is given these examples (usually 3-5 pairs) and must deduce the underlying logic to solve a new input grid. The catch? The rules are deliberately abstract, creative, and diverse—ranging from symmetry and rotations to counting, layering, or even conceptual transformations. There’s no explicit explanation provided, so you have to "reason from scratch."
% For example, a task might show:
% Input: A 3x3 grid with a red square in the center → Output: A 3x3 grid with a blue square in each corner.
% \begin{table}[h]
% \caption{Model Performance on ARC-AGI-1}\label{tab:arc}
% \begin{tabular}{l|c}
% \hline 
% System & ARC-AGI-1 \\ 
% \hline
% GPT-4o & 5.0\% \\
% OpenAI-o1-preview & 18.0\% \\
% \dsro{} & 14.0\% \\
% \dsri{} & 15.8\% \\
% OpenAI-o1 (low) & 20.5\% \\
% OpenAI-o1 (med) & 31.0\% \\
% OpenAI-o1 (high) & 35.0\% \\
% \hline
% \end{tabular}
% \end{table}

% DeepSeek R1, along with its precursor DeepSeek R1-Zero, has shown notable performance on ARC-AGI-1, a specific version of the ARC benchmark. Table \ref{tab:arc} indicates that both models achieve a score of approximately 15\% on ARC-AGI-1. This is a significant step forward compared to earlier large language models like GPT-4o, which scored around 5\% on the same task. However, \dsri{} still lags behind OpenAI-o1 on this task. 






\section{\dsri{} Distillation}\label{sec:distillation}
LLMs are energy-intensive, requiring substantial computational resources, including high-performance GPUs and considerable electricity, for training and deployment. These resource demands present a significant barrier to democratizing access to AI-powered technologies, particularly in under-resourced or marginalized communities.

To address this challenge, we adopt a model distillation approach, a well-established technique for efficient knowledge transfer that has demonstrated strong empirical performance in prior work \citep{busbridge2025distillation, DBLP:journals/corr/HintonVD15}.
Specifically, we fine-tune open-source foundation models such as Qwen \citep{qwen2_5} and LLaMA \citep{touvron2023llama, llama3_1_405b} using a curated dataset comprising 800,000 samples generated with DeepSeek-R1. Details of the dataset construction are provided in Appendix~\ref{sec:sft_data}. We find that models distilled from high-quality teacher outputs consistently outperform those trained directly on human-generated data, corroborating prior findings on the efficacy of distillation \citep{busbridge2025distillation}.


For distilled models, we apply only SFT and do not include an RL stage, even though incorporating RL could substantially boost model performance. Our primary goal here is to demonstrate the effectiveness of the distillation technique, leaving the exploration of the RL stage to the broader research community. For details on distillation training, please see Appendix \ref{appendix:distill-params}.

%Our findings indicate that this straightforward distillation method significantly enhances the reasoning abilities of smaller models.
%The base models we use here are Qwen2.5-Math-1.5B, Qwen2.5-Math-7B, Qwen2.5-14B, Qwen2.5-32B, Llama-3.1-8B, and Llama-3.3-70B-Instruct. We select Llama-3.3 because its reasoning capability is slightly better than that of Llama-3.1.
\label{sec:distilled_model_evaluation}
\input{tables/distill_eval}
We evaluate the distilled models on AIME, GPQA, Codeforces, as well as MATH-500 \citep{lightman2023let} and LiveCodeBench \citep{livecodebench}. For comparison, we use two well-established LLMs as baselines: GPT-4o and Claude-3.5-Sonnet. As shown in Table \ref{tab:distill}, the straightforward distillation of outputs from DeepSeek-R1 allows the distilled model, DeepSeek-R1-Distill-Qwen-1.5B, to surpass non-reasoning baselines on mathematical benchmarks. Notably, it is remarkable that a model with only 1.5 billion parameters achieves superior performance compared to the best closed-source models.
Furthermore, model performance improves progressively as the parameter size of the student model increases.

Our experimental results demonstrate that smaller models can achieve strong performance through distillation. Furthermore, as shown in Appendix~\ref{sec:distillation}, the distillation approach yields superior performance compared to reinforcement learning alone when applied to smaller model architectures. This finding has significant implications for democratizing AI access, as reduced computational requirements enable broader societal benefits.

\subsection{Distillation v.s. Reinforcement Learning} 
\begin{table}[h]
    \centering
    % \footnotesize
    % \setlength{\tabcolsep}{2pt}
    \caption{Comparison of distilled and RL Models on Reasoning-Related Benchmarks.}
    \label{tab:distill_vs_rl}
        % \resizebox{\textwidth}{!}{%
            \begin{tabular}{@{}l *{6}{c} @{}}
        % \begin{tabularx}{\textwidth}{@{}l *{5}{c}@{}}
                \toprule
                % \multirow{3}{*}{\centering\textbf{Model}} 
                %     & \multicolumn{2}{c}{\textbf{AIME 2024}} 
                %     & \textbf{MATH-500} 
                %     & \textbf{GPQA Diamond} 
                %     & \textbf{LiveCodeBench}  \\
                % \cmidrule(lr){2-3}
                %     & pass@1 & cons@64 & pass@1 & pass@1 & pass@1 \\
                \multirow{3}{*}{\centering\textbf{Model}} & \multicolumn{2}{c}{\multirow{2}{*}{\textbf{\footnotesize{AIME 2024}}}} & \multirow{2}{*}{\textbf{\footnotesize{MATH}}} & \textbf{\footnotesize{GPQA}} & \textbf{\footnotesize{LiveCode}} \\
                    &  &  &  & \textbf{\footnotesize{Diamond}} & \textbf{\footnotesize{Bench}} \\
        \cmidrule(lr){2-3}
         & pass@1 & cons@64 & pass@1 & pass@1 & pass@1 \\
                    
                \midrule
                \textbf{QwQ-32B-Preview} 
                    & 50.0 & 60.0 & 90.6 & 54.5 & 41.9  \\
                \textbf{Qwen2.5-32B-Zero} 
                    & 47.0 & 60.0 & 91.6 & 55.0 & 40.2  \\
                \textbf{DeepSeek-R1-Distill-Qwen-32B} 
                    & \bf{72.6} & \bf{83.3} & \bf{94.3} & \bf{62.1} & \bf{57.2} \\
                \bottomrule
            \end{tabular}%
            % \end{tabularx}
        % }
\end{table}


In Section \ref{sec:distilled_model_evaluation}, we can see that by distilling DeepSeek-R1, the small model can achieve impressive results. However, there is still one question left: can the model achieve comparable performance through the large-scale RL training discussed in the paper without distillation?


To answer this question, we conduct large-scale RL training on Qwen2.5-32B-Base using math, code, and STEM data, training for over 10K steps, resulting in Qwen2.5-32B-Zero, as described in \ref{sec: qwen-zero}. The experimental results, shown in Table \ref{tab:distill_vs_rl}, demonstrate that the 32B base model, after large-scale RL training, achieves performance on par with QwQ-32B-Preview. However, DeepSeek-R1-Distill-Qwen-32B, which is distilled from DeepSeek-R1,  performs significantly better than Qwen2.5-32B-Zero across all benchmarks.

Therefore, we can draw two conclusions: First, distilling more powerful models into smaller ones yields excellent results, whereas smaller models relying on the large-scale RL mentioned in this paper require enormous computational power and may not even achieve the performance of distillation. Second, while distillation strategies are both economical and effective, advancing beyond the boundaries of human intelligence may still require more powerful base models and larger-scale reinforcement learning.



\begin{table}[t]
    \centering
    % \setlength{\tabcolsep}{2pt}
    \captionof{table}{Performance of different models on AIME 2024 and AIME 2025.}
    \begin{tabular}{l|cc} \toprule
    Average Score & 
    \textbf{AIME 2024} & \textbf{AIME 2025} \\
    \midrule
    GPT-4o-0513 & 9.3\% & -  \\
    Qwen2-Math-7B-Instruct & 7.9\% & 4.6\% \\
    \midrule
    \textbf{Qwen2-Math-7B-Zero} & 22.3\% & 18.1\% \\
    \bottomrule
    \end{tabular}
    \label{tab:zero-7b-math}
\end{table}


% \subsection{Reinforcement Learning on the Pre-O1 Base Model}

Apart from the experiment based on Qwen-2.5-32B, we conducted experiments on Qwen2-Math-7B (released August 2024) prior to the launch of the first reasoning model, OpenAI-o1 (September 2024), to ensure the base model was not exposed to any reasoning trajectory data. 
We trained Qwen2-Math-7B-Zero with approximately 10,000 policy gradient update steps. As shown in Table \ref{tab:zero-7b-math}, Qwen2-Math-7B-Zero significantly outperformed the non-reasoning models like Qwen2-Math-7B-Instruct and GPT-4o. These results further demonstrate that the model can autonomously develop advanced reasoning strategies through large-scale reinforcement learning.

\section{Discussion}
\subsection{Key Findings}
We highlight our key findings, which may facilitate the community in better reproducing our work.

\textbf{The importance of base checkpoint: } During the initial phase of our development, we experimented with smaller-scale models, specifically a 7B dense model and a 16B Mixture-of-Experts (MoE) model, as the foundational architectures for RL training. However, these configurations consistently failed to yield meaningful improvements when evaluated on the AIME benchmark, which we employed as the primary validation set. We observed that as response lengths increased, these smaller models exhibited a tendency toward repetition and were unable to effectively leverage long chains of thought (CoT) to improve reasoning accuracy.

To address these limitations, we transitioned to larger-scale models, including a 32B dense model \citep{qwen2_5}, a 230B MoE model \citep{dsvii}, and a 671B MoE model \citep{dsviii}. With these more capable architectures, we finally observed substantial performance gains attributable to pure RL training. These findings suggest that the effectiveness of reinforcement learning from base models is highly dependent on the underlying model capacity. We therefore recommend that future research in this area prioritize the use of sufficiently large and expressive models when aiming to validate the efficacy of RL from scratch.

\textbf{The importance of verifiers: }The effectiveness of \dsro{} is highly contingent upon the reliability and fidelity of the reward signal used during training. To date, our investigations indicate that two approaches—rule-based reward models (RMs) and LLMs to assess an answer’s correctness against a predefined ground-truth—serve as robust mechanisms for mitigating issues related to reward hacking. The LLM-based evaluation framework demonstrates particular effectiveness for tasks with well-defined, concise answers, such as single-sentence or phrase-level responses. However, this method exhibits limited generalizability to more complex tasks, including open-ended generation and long-form writing, where the notion of correctness is inherently more subjective and nuanced.

\textbf{Iterative pipeline:}  We propose a multi-stage training pipeline comprising both SFT and RL stages. The RL component enables the model to explore and discover optimal reasoning trajectories for  tasks capabilities that cannot be fully realized through human-annotated reasoning traces alone. In particular, without the RL stage, long-chain reasoning patterns, such as those required in complex Chain-of-Thought (CoT) prompting, would remain largely unexplored. Conversely, the SFT stage plays a crucial role in tasks where reliable reward signals are difficult to define or model, such as open-ended question answering and creative writing. Therefore, both RL and SFT are indispensable components of our training pipeline. Exclusive reliance on RL can lead to reward hacking and suboptimal behavior in ill-posed tasks, while depending solely on SFT may prevent the model from optimizing its reasoning capabilities through exploration.

% However, once these long CoT data are obtained, distillation via SFT proves effective for weaker models, with the subsequent RL stage yielding only marginal improvements. 


% Although distillation via SFT may lead to slightly worse or, in some cases, comparable performance compared to the RL-trained model, reinforcement learning remains indispensable. This is because RL enables the exploration of optimal reasoning trajectories for LLMs on complex problems—something that human reasoning trace data alone cannot achieve. Without RL, such long-chain CoT data would not be discovered. However, once these long CoT data are obtained, distillation via SFT proves effective for weaker models, with the subsequent RL stage yielding only marginal improvements.


\subsection{Unsuccessful Attempts} \label{sec:attempt}
In the early stages of developing \dsri{}, we also encountered failures and setbacks along the way. We share our failure experiences here to provide insights, but this does not imply that these approaches are incapable of developing effective reasoning models.

\paragraph{Process Reward Model (PRM)}
PRM is a reasonable method to guide the model toward better approaches for solving reasoning tasks~\citep{uesato2022solving, lightman2023let,mathshepherd}. However, in practice, PRM has three main limitations that may hinder its ultimate success. First, it is challenging to explicitly define a fine-grain step in general reasoning. 
% Second, evaluating the correctness of key reasoning steps proves to be as difficult as solving the problem itself. 
Second, determining whether the current intermediate step is correct is a challenging task. Automated annotation using models may not yield satisfactory results, while manual annotation is not conducive to scaling up.
Third, once a model-based PRM is introduced, it inevitably leads to reward hacking~\citep{gao2022scalinglawsrewardmodel},  and retraining the reward model needs additional training resources and it complicates the whole training pipeline. In conclusion, while PRM demonstrates a good ability to rerank the top-N responses generated by the model or assist in guided search~\citep{snell2024scalingllmtesttimecompute}, its advantages are limited compared to the additional computational overhead it introduces during the large-scale reinforcement learning process in our experiments.

\paragraph{Monte Carlo Tree Search (MCTS)}
Inspired by AlphaGo~\citep{alphago} and AlphaZero~\citep{alphazero}, we explored using Monte Carlo Tree Search (MCTS) to enhance test-time compute scalability. This approach involves breaking answers into smaller parts to allow the model to explore the solution space systematically. To facilitate this, we prompt the model to generate multiple tags that correspond to specific reasoning steps necessary for the search. For training, we first use collected prompts to find answers via MCTS guided by a pre-trained value model. Subsequently, we use the resulting question-answer pairs to train both the actor model and the value model, iteratively refining the process.

However, this approach encounters several challenges when scaling up the training. First, unlike chess, where the search space is relatively well-defined, token generation presents an exponentially larger search space. To address this, we set a maximum extension limit for each node, but this can lead to the model getting stuck in local optima. Second, the value model directly influences the quality of generation since it guides each step of the search process. Training a fine-grained value model is inherently difficult, which makes it challenging for the model to iteratively improve. While AlphaGo's core success relied on training a value model to progressively enhance its performance, this principle proves difficult to replicate in our setup due to the complexities of token generation.

In conclusion, while MCTS can improve performance during inference when paired with a pre-trained value model, iteratively boosting model performance through self-search remains a significant challenge.


\section{Related Work}

\subsection{Chain-of-thought Reasoning}
Chain-of-thought (CoT) reasoning \citep{wei2022chain} revolutionized how LLMs approach complex reasoning tasks by prompting them to generate intermediate reasoning steps before producing a final answer.
% 
This method significantly improved performance on benchmarks involving arithmetic, commonsense, and symbolic reasoning. Subsequent work explored its scope: \citet{suzgun-etal-2023-challenging} demonstrated that CoT's effectiveness scales with model size, while \citet{kojima2022large} extended it to zero-shot settings by simply instructing models to ``think step by step.''

Building on CoT's framework, numerous ``prompt engineering'' techniques have been proposed to enhance model performance.
\citet{wang2022self} introduced self-consistency, a method that aggregates answers from multiple reasoning paths to improve robustness and accuracy.
\citet{zhou2023leasttomost} developed least-to-most prompting, which decomposes complex problems into sequential subquestions that are solved incrementally. 
\citet{yao2023tree} proposed tree-of-thoughts, enabling models to explore multiple reasoning branches simultaneously and perform deliberate decision-making through looking ahead or backtracking.
Collectively, these approaches leverage human prior knowledge and more structured reasoning frameworks to enhance the reasoning capabilities of LLMs.


\subsection{Scaling Inference-time Compute}

As unsupervised pre-training scaling might be constrained by the amount of available human data \citep{kaplan2020scaling, muennighoff2023scaling}, scaling compute during inference has become even more critical \citep{snell2025scaling}. Broadly, we define methods that improve model performance by increasing inference compute as forms of scaling inference-time compute.

% 1. multiple-generation
A straightforward approach trades compute for performance by generating multiple diverse reasoning chains and selecting the best answer. The optimal answer can be identified using a separate reranker \citep{gsm8k, brown2024large}, process-based reward models \citep{uesato2022solving, lightman2023let}, or simply by selecting the most common answer \citep{wang2022self}.
% 2. search
Search methods, such as Monte Carlo Tree Search and Beam Search, also guide exploration of the solution space more effectively \citep{hao2023reasoning, feng2024alphazeroliketreesearchguidelarge, xin2024deepseekproverv15harnessingproofassistant, AlphaGeometryTrinh2024}.
% 3. self-correct
Beyond parallel generation, self-correct techniques prompt or train models to iteratively critique and refine their outputs \citep{welleck2023generating, madaan2023selfrefine, kumar2024training}, often incorporating external feedback to enhance reliability \citep{yao2023react, gou2024critic}.
% 4. tool-use
Additionally, some methods improve performance by integrating tool use during testing, which is particularly effective for knowledge-intensive \citep{nakano2021webgpt} and compute-intensive tasks \citep{schick2023toolformer, gou2024tora, chen2025empirical}.
% 5. ttt and others
Test-time training (TTT) further updates the model during inference to boost performance \citep{sun2020test, akyurek2024surprising}. There are also various other inference-time scaling approaches that—either implicitly \citep{geiping2025scaling} or explicitly \citep{zelikman2024quietstar}—allocate more compute for each token.

In contrast, our work shows that LLMs can achieve scalable improvements through additional RL compute and increased test-time compute (i.e., more tokens). We integrate the benefits of scaling at test time into a broader framework that uses reinforcement learning to incentivize enhanced in-context search abilities.


\subsection{Reinforcement Learning for Reasoning Enhancement}

Reinforcement Learning plays a pivotal role in aligning LLMs with human preferences \citep{ouyang2022training, bai2022training}. Despite its importance, few studies have focused on using RL to enhance reasoning capabilities. Traditional RL pipelines begin with SFT on high-quality human demonstrations, which provides a strong initialization and prevents mode collapse. Following this, a reward model is trained on human preferences, and the language model is subsequently optimized using methods such as PPO \citep{schulman2017proximal} or DPO \citep{dpo}. Although this method works well for alignment, it risks constraining models to emulate human reasoning patterns, potentially hindering the discovery of novel problem-solving strategies.

% star 
Methods like STaR iteratively boost performance by fine-tuning on the model's self-generated chain-of-thought that leads to correct final answers \citep{zelikman2022star, yuan2023scaling, singh2024beyond}.
% prm
Recent studies have also investigated the use of process-based rewards that emphasize both the correctness of final answers and the soundness of the reasoning processes \citep{lightman2023let, mathshepherd,deepseekmath}.
% ours
Unlike these methods, our work applies outcome-based RL directly to base language models without an initial SFT phase. This design choice encourages the emergence of innovative and unconstrained reasoning strategies, enabling the model to develop diverse solutions beyond mere imitation of human examples.
% rep
Our approach also inspired further exploration in subsequent research \citep{tinyzero, liu2025oatzero, openr1}.


\section{Open Weights, Code, and Data} \label{sec:open}
To promote the development of the open-source community and industry ecosystem, we have made the model weights of \href{https://huggingface.co/deepseek-ai/DeepSeek-R1}{\dsri{}} and \href{https://huggingface.co/deepseek-ai/\dsro{}}{\dsro{}} publicly available on HuggingFace. In addition, we release   \href{https://huggingface.co/deepseek-ai/DeepSeek-R1-Distill-Qwen-1.5B}{DeepSeek-R1-Distill-Qwen-1.5B},  \href{https://huggingface.co/deepseek-ai/DeepSeek-R1-Distill-Qwen-7B}{DeepSeek-R1-Distill-Qwen-7B}, \href{https://huggingface.co/deepseek-ai/DeepSeek-R1-Distill-Qwen-14B}{DeepSeek-R1-Distill-Qwen-14B}, \href{https://huggingface.co/deepseek-ai/DeepSeek-R1-Distill-Qwen-32B}{DeepSeek-R1-Distill-Qwen-32B},        \href{https://huggingface.co/deepseek-ai/DeepSeek-R1-Distill-Llama-8B}{DeepSeek-R1-Distill-Llama-8B},        \href{https://huggingface.co/deepseek-ai/DeepSeek-R1-Distill-Llama-70B}{DeepSeek-R1-Distill-Llama-70B}.

Furthermore, we have released the fundamental model inference code (\url{https://github.com/deepseek-ai/DeepSeek-V3}) and provided detailed usage guidelines (\url{https://github.com/deepseek-ai/DeepSeek-R1}) on GitHub.

Here is an example of running the inference code to interact with DeepSeek-R1:

\begin{lstlisting}[language=bash,style=code]
# Download the model weights from Hugging Face
huggingface-cli download deepseek-ai/DeepSeek-R1 --local-dir /path/to/DeepSeek-R1

# Clone DeepSeek-V3 GitHub repository
git clone https://github.com/deepseek-ai/DeepSeek-V3.git

# Install necessary dependencies
cd DeepSeek-R1/inference
pip install -r requirements.txt

# Convert Hugging Face model weights to a specific format (for running the model on 16 H800 GPUs)
python convert.py --hf-ckpt-path /path/to/DeepSeek-R1 --save-path /path/to/DeepSeek-R1-Demo --n-experts 256 --model-parallel 16

# Run the model and interact with it
torchrun --nnodes 2 --nproc-per-node 8 --node-rank $RANK --master-addr $MASTER_ADDR generate.py --ckpt-path /path/to/DeepSeek-R1-Demo --config configs/config_671B.json --interactive --temperature 0.7 --max-new-tokens 8192
\end{lstlisting}


We also release SFT and RL data to the public at xxx. In the review process, we upload the data as an attachment.



\section{Evaluation Prompts and Settings}
\begin{table}[h]
    \centering \small
        \caption{MMLU assesses a model's factual and conceptual understanding across 57 tasks spanning STEM (science, technology, engineering, mathematics), humanities, social sciences, and professional fields (e.g., law, medicine). The benchmark is commonly used to evaluate a model's ability to perform general knowledge reasoning and multitask proficiency across a diverse range of subjects and tasks. Here is an example of MMLU. }
    \label{tab:mmlu_eval_format_example}
    \begin{tabular}{p{0.95\textwidth}}
    \hline
    \textbf{PROMPT}\\
    Answer the following multiple choice question. The last line of your response should be of the following format: 'Answer: \$LETTER' (without quotes) where LETTER is one of ABCD. Think step by step before answering.\\[1ex] 
    Which tool technology is associated with Neandertals?\\[1ex]
    A. Aurignacian\\
    B. Acheulean\\
    C. Mousterian\\
    D. both b and c \\ \hline
      \textbf{Evaluation}\\
      Parse the last line in response to judge if the  choice equals to ground truth. \\ \hline
    \end{tabular}
\end{table}

\begin{table}[h]
    \centering \small
        \caption{MMLU-Redux is a subset of 5,700 manually re-annotated questions across all 57 MMLU subjects. MMLU-Redux focuses on improving the quality, clarity, and robustness of the benchmark by reducing noise, ambiguities, and potential biases in the MMLU, while potentially adjusting the scope or difficulty of tasks to better align with modern evaluation needs. Here is an example of MMLU-Redux.}
    \label{tab:mmluredux_eval_format_example}
    \begin{tabular}{p{0.95\textwidth}}
    \hline
    \textbf{PROMPT}\\
    \#\# Question: 

Sauna use, sometimes referred to as "sauna bathing," is characterized by short-term passive exposure to extreme heat \ldots In fact, sauna use has been proposed as an alternative to exercise for people who are unable to engage in physical activity due to chronic disease or physical limitations.[13]

According to the article, which of the following is NOT a benefit of sauna use?

\#\# Choices:

- (A) Decreased risk of heart attacks. \\
- (B) Increase in stroke volume. \\
- (C) Improved mental health. \\
- (D) Decreased rate of erectile dysfunction.


\#\# Instruction 

Please answer this question by first reasoning and then selecting the correct choice. \\
Present your reasoning and solution in the following json format. \\
Please show your choice in the `answer` field with only the choice letter, e.g.,`"answer": "C"`. \\

\{ \\
    "reasoning": "\_\_\_", \\
    "answer": "\_\_\_" \\
\}\\ \hline
      \textbf{Evaluation}\\
      Parse the json output in response to judge if the answer equals to ground truth. \\ \hline
    \end{tabular}
\end{table}


\begin{table}[h]
    \centering \small
        \caption{LiveCodeBench aims to evaluate model performance on the algorithm competition task, which collects new problems over time from contests across three competition platforms, namely LeetCode, AtCoder, and CodeForces. }
    \label{tab:lcb_eval_format_example}
    \begin{tabular}{p{0.95\textwidth}}
    \hline
    \textbf{PROMPT}\\
Question:
There is a stack of N cards, and the i\-th card from the top has an integer $A_i$ written on it. \\
You take K cards from the bottom of the stack and place them on top of the stack, maintaining their order. \\
Print the integers written on the cards from top to bottom after the operation. \\
Input\\
The input is given from Standard Input in the following format:\\
N K\\
$A_1 A_2 \ldots A_N$\\
Output\\
Let $B_i$ be the integer written on the i\-th card from the top of the stack after the operation. Print $B_1,B_2,\ldots,B_N$ in this order, separated by spaces.\\
Constraints\\
$-1 \leq K < N \leq 100$\\
$-1 \leq A_i \leq 100$\\
All input values are integers.\\
Sample Input 1\\
5 3 \\
1 2 3 4 5 \\
Sample Output 1 \\
3 4 5 1 2\\
Initially, the integers written on the cards are 1,2,3,4,5 from top to bottom.
After taking three cards from the bottom of the stack and placing them on top, the integers written on the cards become 3,4,5,1,2 from top to bottom.\\
\\
Sample Input 2\\
\\
6 2\\
1 2 1 2 1 2\\
\\
Sample Output 2\\
\\
1 2 1 2 1 2\\
\\
The integers written on the cards are not necessarily distinct.\\
\\
Please write a python code to solve the above problem. Your code must read the inputs from stdin and output the results to stdout.\\ \hline
      \textbf{Evaluation}\\
      Extract the code wrapped by \texttt{```}python\texttt{```} in response to judge if the answer passes the test cases. \\ \hline
    \end{tabular}
\end{table}

\begin{table}[h]
    \centering \small
        \caption{Compared to MMLU, MMLU-Pro features a curated subset of tasks, but with significantly increased difficulty. Questions in MMLU-Pro are designed to require deeper reasoning, multi-step problem-solving, and advanced domain-specific knowledge. For example, STEM tasks may involve complex mathematical derivations or nuanced scientific concepts, while humanities tasks may demand intricate contextual analysis. }
    \label{tab:mmlu-pro}
    \begin{tabular}{p{0.95\textwidth}}
    \hline
    \textbf{PROMPT}\\
    The following are multiple choice questions (with answers) about business. Think step by step and then output the answer in the format of "The answer is (X)" at the end.

\ldots

Question: Typical advertising regulatory bodies suggest, for example that adverts must not: encourage \_\_\_, cause unnecessary \_\_\_ or \_\_\_, and must not cause \_\_\_ offence. \\
Options: A. Safe practices, Fear, Jealousy, Trivial \\
B. Unsafe practices, Distress, Joy, Trivial \\
C. Safe practices, Wants, Jealousy, Trivial \\
D. Safe practices, Distress, Fear, Trivial \\
E. Unsafe practices, Wants, Jealousy, Serious \\
F. Safe practices, Distress, Jealousy, Serious \\
G. Safe practices, Wants, Fear, Serious \\
H. Unsafe practices, Wants, Fear, Trivial \\
I. Unsafe practices, Distress, Fear, Serious \\
Answer: Let's think step by step.\\ \hline
      \textbf{Evaluation}\\
      Parse the capital letter following ``Answer: '' in response to judge if the answer equals to ground truth. \\ \hline
    \end{tabular}
\end{table}

\begin{table}[h]
    \centering \small
        \caption{DROP assesses a model's ability to understand and extract relevant information from extended textual passages. Unlike simpler question-answering benchmarks that focus on factual recall, DROP requires models to process and interpret context-rich paragraphs. }
    \label{tab:drop_eval_format_example}
    \begin{tabular}{p{0.95\textwidth}}
    \hline
    \textbf{PROMPT}\\
    You will be asked to read a passage and answer a question. Some examples of passages and Q\&A are provided below. 

\# Examples  
---
Passage:  Looking to avoid back-to-back divisional losses, the Patriots traveled to Miami to face the 6-4 Dolphins at Dolphin Stadium \ldots Cassel's 415 passing yards made him the second quarterback in Patriots history to throw for at least 400 yards in two or more games; Drew Bledsoe had four 400+ yard passing games in his Patriots career.

Question: How many points did the Dolphins lose by?
Answer: 20.

---
Passage:  In week 2, the Seahawks took on their division rivals, the San Francisco 49ers. Prior to the season, NFL analysts rated this rivalry as the top upcoming rivalry, as well as the top rivalry of the decade \ldots Seattle was now 2-0, and still unbeaten at home.

Question: How many field goals of at least 30 yards did Hauschka make?
Answer: 2.

---
Passage:  at Raymond James Stadium, Tampa, Florida TV Time: CBS 1:00pm eastern The Ravens opened the regular season on the road against the Tampa Bay Buccaneers on September 10. \ldots With the win, the Ravens were 1-0 and 1-0 against NFC Opponents.

Question: how many yards did lewis get
Answer: 4.
\# Your Task

---
Passage:  The Chargers (1-0) won their season opener 22-14 against the Oakland Raiders after five field goals by Nate Kaeding and three botched punts by the Raiders. The Raiders Pro Bowl long snapper Jon Condo suffered a head injury in the second quarter. He was replaced by linebacker Travis Goethel, who had not snapped since high school. Goethel rolled two snaps to punter Shane Lechler, each giving the Chargers the ball in Raiders territory, and Lechler had another punt blocked by Dante Rosario. The Chargers scored their only touchdown in the second quarter after a 13-play, 90-yard drive resulted in a 6-yard touchdown pass from Philip Rivers to wide receiver Malcom Floyd. The Chargers failed to score four out of five times in the red zone. San Diego led at halftime 10-6, and the Raiders did not scored a touchdown until 54 seconds remained in the game. Undrafted rookie Mike Harris made his first NFL start, filing in for left tackle for an injured Jared Gaither. San Diego protected Harris by having Rivers throw short passes; sixteen of Rivers' 24 completions were to running backs and tight ends, and he threw for 231 yards while only being sacked once. He did not have an interception after throwing 20 in 2011. The win was the Chargers' eighth in their previous nine games at Oakland. It improved Norv Turner's record to 4-2 in Chargers' season openers. Running back Ryan Mathews and receiver Vincent Brown missed the game with injuries.

Question: How many yards did Rivers pass?
Answer: 

Think step by step, then write a line of the form "Answer: \$ANSWER" at the end of your response. \\ \hline
      \textbf{Evaluation}\\
      Parse the capital letter following ``Answer: '' in response to judge if the answer equals to ground truth. \\ \hline
    \end{tabular}
\end{table}

\begin{table}[h]
    \centering \small
        \caption{Instruction‑Following Evaluation (IFEval) is a benchmark designed to assess a model’s ability to comply with explicit, verifiable instructions embedded within prompts. It targets a core competency of large language models (LLMs): producing outputs that meet multiple, clearly defined constraints specified by the user.}
    \label{tab:ifeval_eval_format_example}
    \begin{tabular}{p{0.95\textwidth}}
    \hline
    \textbf{PROMPT}\\
    Kindly summarize the text below in XML format. Make sure the summary contains less than 4 sentences.

Quantum entanglement is the phenomenon that occurs when a group of particles are generated, interact, or share spatial proximity in such a way that the quantum state of each particle of the group cannot be described independently of the state of the others, including when the particles are separated by a large distance. The topic of quantum entanglement is at the heart of the disparity between classical and quantum physics: entanglement is a primary feature of quantum mechanics not present in classical mechanics.

Measurements of physical properties such as position, momentum, spin, and polarization performed on entangled particles can, in some cases, be found to be perfectly correlated. For example, if a pair of entangled particles is generated such that their total spin is known to be zero, and one particle is found to have clockwise spin on a first axis, then the spin of the other particle, measured on the same axis, is found to be anticlockwise. However, this behavior gives rise to seemingly paradoxical effects: any measurement of a particle's properties results in an apparent and irreversible wave function collapse of that particle and changes the original quantum state. With entangled particles, such measurements affect the entangled system as a whole.

Such phenomena were the subject of a 1935 paper by Albert Einstein, Boris Podolsky, and Nathan Rosen, and several papers by Erwin Schrödinger shortly thereafter, describing what came to be known as the EPR paradox. Einstein and others considered such behavior impossible, as it violated the local realism view of causality (Einstein referring to it as "spooky action at a distance") and argued that the accepted formulation of quantum mechanics must therefore be incomplete.\\ \hline
      \textbf{Evaluation}\\
      Call official functions to check if the answer is consistent with the instructions. \\ \hline
    \end{tabular}
\end{table}

\begin{table}[h]
    \centering \small
        \caption{FRAMES (Factuality, Retrieval, And reasoning MEasurement Set) is a comprehensive benchmark designed to evaluate core components of retrieval-augmented generation (RAG) systems. Our evaluation employs the benchmark's official "Oracle Prompt" configuration. In this setting, each test prompt includes the question along with all the ground truth Wikipedia articles, thus eliminating the need for an external retrieval component (e.g., BM25). This setting allows us to specifically measure a model's ability to reason over and synthesize information from provided sources to generate correct and verifiable facts. }
    \label{tab:frames_eval_format_example}
    \begin{tabular}{p{0.95\textwidth}}
    \hline
    \textbf{PROMPT}\\
    Here are the relevant Wikipedia articles: \\
url: https:\/\/en.wikipedia.org\/wiki\/President\_of\_the\_United\_States \\
url content: The president of the United States (POTUS) is the head of state and head of government of the United States of America. The president directs the executive branch of the federal government and is the commander-in-chief of the United States Armed Forces.
\ldots

Based on all the information, answer the query. 

Query: If my future wife has the same first name as the 15th first lady of the United States' mother and her surname is the same as the second assassinated president's mother's maiden name, what is my future wife's name? \\ \hline
      \textbf{Evaluation}\\
      ===Task===\\
I need your help in evaluating an answer provided by an LLM against a ground truth answer. Your task is to determine if the ground truth answer is present in the LLM's response. Please analyze the provided data and make a decision. \\
===Instructions=== \\
1. Carefully compare the "Predicted Answer" with the "Ground Truth Answer".\\
2. Consider the substance of the answers - look for equivalent information or correct answers.\\
Do not focus on exact wording unless the exact wording is crucial to the meaning.\\
3. Your final decision should be based on whether the meaning and the vital facts of the
"Ground Truth Answer" are present in the "Predicted Answer:"\\
===Input Data===
- Question: If my future wife has the same first name as the 15th first lady of the United States' mother and her surname is the same as the second assassinated president's mother's maiden name, what is my future wife's name? \\
- Predicted Answer: \ldots \\
- Ground Truth Answer: Jane Ballou\\
===Output Format===\\
Provide your final evaluation in the following format:\\
Explanation: xxx\\
Decision: "TRUE" or "FALSE"\\
Please proceed with the evaluation.\\ \hline
    \end{tabular}
\end{table}


\begin{table}[h]
    \centering \small
        \caption{Arena-Hard is an open-ended evaluation benchmark specifically designed to assess the capabilities of LLMs. It presents models with challenging, novel, and diverse prompts curated from Chatbot Arena, a continuously evolving, crowd-sourced platform. It focuses on measuring model performance in open-ended tasks, with particular emphasis on coding and mathematics-related prompts.
        Given the inherently subjective nature of open-ended tasks, where multiple valid responses may exist, the benchmark necessitates the use of an evaluation model to approximate human judgment effectively. Higher evaluation scores suggest that the model is more likely to be favored by human users in real-world scenarios.   }
    \label{tab:arenahard_eval_format_example}
    \begin{tabular}{p{0.95\textwidth}}
    \hline
    \textbf{PROMPT}\\
    To write an SQL query that selects the top 10 rows in a database and joins to 3 different tables based on a field called "code", you would need to know the names of the tables and the specific fields you want to select. Assuming the tables are named "table1", "table2", and "table3", and you want to select all fields from each table, the query would look like this:

```sql
SELECT *
FROM (
  SELECT * 
  FROM table1
  LIMIT 10
) AS t1
JOIN table2 ON t1.code = table2.code
JOIN table3 ON t1.code = table3.code;
```

This query first selects the top 10 rows from "table1" and then joins "table2" and "table3" based on the "code" field. If you want to select specific fields from each table, you can replace the asterisks (*) with the field names separated by commas.\\ \hline
      \textbf{Evaluation}\\
      Please act as an impartial judge and evaluate the quality of the responses provided by two AI assistants to the user prompt displayed below. You will be given assistant A's answer and assistant B's answer. Your job is to evaluate which assistant's answer is better.

Begin your evaluation by generating your own answer to the prompt. You must provide your answers before judging any answers.

When evaluating the assistants' answers, compare both assistants' answers with your answer. You must identify and correct any mistakes or inaccurate information.

Then consider if the assistant's answers are helpful, relevant, and concise. Helpful means the answer correctly responds to the prompt or follows the instructions. Note when user prompt has any ambiguity or more than one interpretation, it is more helpful and appropriate to ask for clarifications or more information from the user than providing an answer based on assumptions. Relevant means all parts of the response closely connect or are appropriate to what is being asked. Concise means the response is clear and not verbose or excessive.

Then consider the creativity and novelty of the assistant's answers when needed. Finally, identify any missing important information in the assistants' answers that would be beneficial to include when responding to the user prompt.

After providing your explanation, you must output only one of the following choices as your final verdict with a label:

1. Assistant A is significantly better: [[A$>>$B]]\\
2. Assistant A is slightly better: [[$A>$B]]\\
3. Tie, relatively the same: [[A$=$B]]\\
4. Assistant B is slightly better: [[B$>$A]]\\
5. Assistant B is significantly better: [[B$>>$A]]\\

Example output: "My final verdict is tie: [[A$=$B]]". \\ \hline
    \end{tabular}
\end{table}


\begin{table}[h]
    \centering \small
        \caption{AlpacaEval 2.0 is an open-ended evaluation dataset, similar in nature to ArenaHard, and leverages an LLM to assess model performance on subjective tasks. However, in contrast to ArenaHard, the prompts in AlpacaEval 2.0 are generally less challenging and only a small subset necessitates the deployment of reasoning capabilities by the evaluated models. }
    \label{tab:alpacaeval_eval_format_example}
    \begin{tabular}{p{0.95\textwidth}}
    \hline
    \textbf{PROMPT}\\
    What are the names of some famous actors that started their careers on Broadway?
    \\ \hline
      \textbf{Evaluation}\\
      $<|im_start|>$system \\
You are a highly efficient assistant, who evaluates and selects the best large language model (LLMs) based on the quality of their responses to a given instruction. This process will be used to create a leaderboard reflecting the most accurate and human-preferred answers.\\
$<|im_end|>$\\
$<|im_start|>$user\\
I require a leaderboard for various large language models. I'll provide you with prompts given to these models and their corresponding outputs. Your task is to assess these responses, and select the model that produces the best output from a human perspective.

\#\# Instruction

\{\\
    "instruction": """\{instruction\}""",\\
\}\\

\#\# Model Outputs

Here are the unordered outputs from the models. Each output is associated with a specific model, identified by a unique model identifier.

\{\\
    \{\\
        "model\_identifier": "m",\\
        "output": """\{output\_1\}"""\\
    \},\\
    \{\\
        "model\_identifier": "M",\\
        "output": """\{output\_2\}"""\\
    \}\\
\}

\#\# Task

Evaluate the models based on the quality and relevance of their outputs, and select the model that generated the best output. Answer by providing the model identifier of the best model. We will use your output as the name of the best model, so make sure your output only contains one of the following model identifiers and nothing else (no quotes, no spaces, no new lines, ...): m or M.

\#\# Best Model Identifier\\
$<|im_end|>$\\ \hline
    \end{tabular}
\end{table}

\begin{table}[h]
    \centering \small
        \caption{The CLUEWSC (Chinese Language Understanding Evaluation Benchmark - Winograd Schema Challenge) is a specialized task within the CLUE benchmark suite designed to evaluate a model's commonsense reasoning and contextual understanding capabilities in Chinese. }
    \label{tab:cluewsc_eval_format_example}
    \begin{tabular}{p{0.95\textwidth}}
    \hline
    \textbf{PROMPT}\\
    请参考示例的格式，完成最后的测试题。

下面是一些示例：
他伯父还有许多女弟子，大半是富商财主的外室；这些财翁白天忙着赚钱，怕小公馆里的情妇长日无聊，要不安分，常常叫她们学点玩艺儿消遣。\\
上面的句子中的"她们"指的是\\
情妇

耶律克定说到雁北义军时，提起韦大哥，就连声说不可挡、不可挡，似有谈虎色变之味。后来又听说粘罕在云中，特派人厚币卑词，要与‘韦义士修好’。韦大哥斩钉截铁地回绝了，大义凛然，端的是条好汉。如今张孝纯也想结识他，几次三番派儿子张浃上门来厮缠，定要俺引他上雁门山去见韦大哥。\\
上面的句子中的"他"指的是\\
张浃

“你何必把这事放在心上？何况你的还不过是手稿，并没有发表出来。”龙点睛越发坦率：“如果发表出来了，那倒也就算了。不过既然没发表出来，我何必还让它飘在外头呢？你给我找一找吧，我要收回。”\\
上面的句子中的"它"指的是 \\
手稿

这个身材高大，曾经被孙光平拿着菜刀追赶得到处乱窜的年轻人，那天早晨穿上了全新的卡其布中山服，像一个城里来的干部似的脸色红润，准备过河去迎接他的新娘。\\
上面的句子中的"他"指的是 \\
年轻人

负责接待我们的是两位漂亮的朝鲜女导游，身材高挑，露出比例完美的小腿。一个姓韩，一个姓金，自称小韩和小金。她们的中文说得毫无口音，言谈举止也相当亲切。\\
上面的句子中的"她们"指的是 \\
两位漂亮的朝鲜女导游

下面是测试题，请在思考结束后（</think>后）用一句话输出答案，不要额外的解释。

崩龙珍夫妻康健和美；鞠琴十年前丧偶，两年前重结良缘，现在的老伴是一位以前未曾有过婚史的高级工程师；崩龙珍和鞠琴都尽量避免谈及自己的爱人，也尽量回避提及蒋盈波的亡夫屈晋勇——尽管她们对他都很熟悉；当然也绝不会愚蠢地提出蒋盈波今后是一个人过到底还是再找个老伴的问题来加以讨论，那无论如何还为时过早。\\
上面的句子中的"他"指的是
\\ \hline
      \textbf{Evaluation}\\
      Parse the last line in response to judge if the answer equals to ground truth.
       \\ \hline
    \end{tabular}
\end{table}


\begin{table}[h]
    \centering \small
        \caption{C-EVAL evaluates a model's breadth and depth of knowledge across 52 diverse academic disciplines, spanning humanities, social sciences, STEM (Science, Technology, Engineering, and Mathematics), and other professional fields (e.g., medicine, law). All question in C-Eval are Chinese. }
    \label{tab:ceval_eval_format_example}
    \begin{tabular}{p{0.95\textwidth}}
    \hline
    \textbf{PROMPT}\\
以下是中国关于逻辑学考试的单项选择题，请选出其中的正确答案。

1991年6月15日，菲律宾吕宋岛上的皮纳图博火山突然大喷发，2000万吨二氧化硫气体冲入平流层，形成的霾像毯子一样盖在地球上空，把部分要照射到地球的阳光反射回太空几年之后，气象学家发现这层霾使得当时地球表面的温度累计下降了0．5℃，而皮纳图博火山喷发前的一个世纪，因人类活动而造成的温室效应已经使地球表面温度升高1℃。某位持“人工气候改造论”的科学家据此认为，可以用火箭弹等方式将二氧化硫充入大气层，阻挡部分阳光，达到地球表面降温的目的。以下哪项如果为真,最能对该科学家的提议构成质疑?\_\_\_ \\
A. 如果利用火箭弹将二氧化硫充入大气层，会导致航空乘客呼吸不适。\\
B. 火山喷发形成的降温效应只是暂时的，经过一段时间温度将再次回升。\\
C. 可以把大气层中的碳取出来存储在地下，减少大气层的碳含量。\\
D. 不论何种方式，“人工气候改造”都将破坏地区的大气层结构。\\
答案： B

\ldots

新疆的哈萨克人用经过训练的金雕在草原上长途追击野狼。某研究小组为研究金雕的飞行方向和判断野狼群的活动范围，将无线电传导器放置在一只金雕身上进行追踪。野狼为了觅食，其活动范围通常很广。因此，金雕追击野狼的飞行范围通常也很大。然而两周以来，无线电传导器不断传回的信号显示，金雕仅在放飞地3公里的范围内飞行。以下哪项如果为真，最有助于解释上述金雕的行为?\_\_\_ \\
A. 金雕放飞地周边重峦叠嶂，险峻异常。\\
B. 金雕的放飞地2公里范围内有一牧羊草场，成为狼群袭击的目标。\\
C. 由于受训金雕的捕杀，放飞地广阔草原的野狼几乎灭绝了。\\
D. 无线电传导信号仅能在有限的范围内传导。\\
\\ \hline
      \textbf{Evaluation}\\
      Parse the last line in response to judge if the choice equals to ground truth.
       \\ \hline
    \end{tabular}
\end{table}




\begin{table}[h]
    \centering \small
        \caption{GPQA (Graduate‑Level Google‑Proof QA Benchmark) is a rigorous evaluation framework designed to measure an LLM’s ability to tackle complex, graduate-level multiple‑choice problems in STEM domains—specifically biology, physics, and chemistry. }
    \label{tab:gpqa_eval_format_example}
    \begin{tabular}{p{0.95\textwidth}}
    \hline
    \textbf{PROMPT}\\
    Answer the following multiple choice question. The last line of your response should be of the following format: 'ANSWER: \$LETTER' (without quotes) where LETTER is one of ABCD. Think step by step before answering.

Two quantum states with energies E1 and E2 have a lifetime of $10^{-9}$ sec and $10^{-8}$ sec, respectively. We want to clearly distinguish these two energy levels. Which one of the following options could be their energy difference so that they can be clearly resolved?

A) $10^{-9}$ eV \\
B) $10^{-8}$ eV \\
C) $10^{-4}$ eV \\
D) $10^{-11}$ eV \\ \hline
      \textbf{Evaluation}\\
      Parse the capital letter following ``ANSWER: '' in response to judge if the choice equals to ground truth. \\ \hline
    \end{tabular}
\end{table}

\begin{table}[h]
    \centering \small
        \caption{SimpleQA is a factuality evaluation benchmark that measures a model’s ability to answer short, fact-seeking questions with precise, verifiable correctness. }
    \label{tab:simpleqa_eval_format_example}
    \begin{tabular}{p{0.95\textwidth}}
    \hline
    \textbf{PROMPT}\\
    Who received the IEEE Frank Rosenblatt Award in 2010? \\ \hline
      \textbf{Evaluation}\\
      Your job is to look at a question, a gold target, and a predicted answer, and then assign a grade of either ["CORRECT", "INCORRECT", "NOT\_ATTEMPTED"].
First, I will give examples of each grade, and then you will grade a new example.


The following are examples of CORRECT predicted answers.

Question: What are the names of Barack Obama's children? \\
Gold target: Malia Obama and Sasha Obama \\
Predicted answer 1: sasha and malia obama \\
Predicted answer 2: most people would say Malia and Sasha, but I'm not sure and would have to double check \\
\ldots


The following are examples of INCORRECT predicted answers.

Question: What are the names of Barack Obama's children? \\
Gold target: Malia and Sasha \\
Predicted answer 1: Malia. \\
Predicted answer 2: Malia, Sasha, and Susan. \\
\ldots


The following are examples of NOT\_ATTEMPTED predicted answers.

Question: What are the names of Barack Obama's children? \\
Gold target: Malia and Sasha \\
Predicted answer 1: I don't know. \\
Predicted answer 2: I need more context about which Obama you are 
\ldots


Also note the following things: \\
\ldots

Here is a new example. Simply reply with either CORRECT, INCORRECT, NOT\_ATTEMPTED. Don't apologize or correct yourself if there was a mistake; we are just trying to grade the answer.\\

Question: Who received the IEEE Frank Rosenblatt Award in 2010?\\
Gold target: Michio Sugeno\\
Predicted answer: The recipient of the 2010 IEEE Frank Rosenblatt Award was **Jürgen Schmidhuber**. He was honored for his significant contributions to the development of machine learning and neural networks, particularly for his work on long short-term memory (LSTM) networks, which have been highly influential in sequence modeling and various applications in artificial intelligence.


Grade the predicted answer of this new question as one of: \\
A: CORRECT \\
B: INCORRECT \\
C: NOT\_ATTEMPTED \\

Just return the letters "A", "B", or "C", with no text around it. \\ \hline
    \end{tabular}
\end{table}


\begin{table}[h]
    \centering \small
        \caption{An example of C-SimpleQA. It measures a model’s ability to answer short, fact-seeking questions in Chinese with precise, verifiable correctness. }
    \label{tab:csimpleqa_eval_format_example}
    \begin{tabular}{p{0.95\textwidth}}
    \hline
    \textbf{PROMPT}\\
显脉香茶菜可以用来治疗急性的什么类型的黄疸型肝炎？
\\ \hline
      \textbf{Evaluation}\\
      请根据给定问题、标准答案和模型预测的答案来评估模型的回答是否正确。您的任务是将结果评定为：【正确】、【错误】或【未尝试】。

首先，我们将列出每个评定类别的示例，然后请您对新问题的预测答案进行评定。

以下是【正确】的答复示例：

问题：贝拉克·奥巴马的孩子叫什么名字？\\
标准答案：玛丽亚·奥巴马和萨莎·奥巴马\\
模型预测1：Malia Obama and Sasha Obama\\
模型预测2：玛丽亚和萨沙\\
\ldots

以下是【错误】的答复示例：

问题：巴拉克·奥巴马的孩子叫什么名字？\\
标准答案：玛丽亚·奥巴马和萨莎·奥巴马\\
模型预测1：玛丽亚\\
模型预测2：玛丽亚、萨莎和苏珊\\
\dots

以下是【未尝试】的答复示例：

问题：巴拉克·奥巴马的孩子叫什么名字？\\
标准答案：玛丽亚·奥巴马和萨莎·奥巴马\\
模型预测1：我不知道。\\
模型预测2：我需要更多关于您所指奥巴马的上下文。\\
\ldots

下面是一个新的问题示例。请只回复A、B、C之一，不要道歉或纠正自己的错误，只需要评估该回答。

问题: 显脉香茶菜可以用来治疗急性的什么类型的黄疸型肝炎？\\
正确答案: 黄疸型肝炎\\
预测答案: \dots


将此新问题的预测答案评定为以下之一：\\
A:【正确】\\
B:【错误】\\
C:【未尝试】\\

只返回字母"A"、"B"或"C"，无须添加其他文本。
       \\ \hline
    \end{tabular}
\end{table}

% =============== math ===============

\begin{table}[h]
    \centering \small
        \caption{An example of math evaluation, which applies to AIME, MATH, and CNMO. These benchmarks evaluate model performance on mathematical tasks. } 
    \label{tab:math_eval_format_example}
    \begin{tabular}{p{0.95\textwidth}}
    \hline
    \textbf{PROMPT}\\
    Let $b \geq 2$ be an integer. Call a positive integer $n$ $b\textit{-eautiful}$ if it has exactly two digits when expressed in base $b$, and these two digits sum to $\sqrt{n}$. For example, $81$ is $13$-eautiful because $81=\underline{6}$ $\underline{3}_{13}$ and $6+3=\sqrt{81}$. Find the least integer $b\geq 2$ for which there are more than ten $b$-eautiful integers. \\
    Please reason step by step, and put your final answer within \textbackslash boxed\{\}. \\
    \hline
  \textbf{Evaluation}\\
  
  Parse the final answer within \textbackslash boxed\{\} and use a rule-based grader to determine if it equals the ground truth. Round numerical values as needed, and use `SymPy'\footnote{\url{https://www.sympy.org}} to parse expressions. \\
      \hline
    \end{tabular}
\end{table}

\clearpage

\bibliography{sn-bibliography}% common bib file


\end{CJK*}
\end{document}
